\documentclass[11pt,a4paper,oneside,openany]{book}
\usepackage[T1]{fontenc}
\usepackage[utf8]{inputenc}
\usepackage[polish]{babel}
\usepackage{lmodern}
\usepackage{amsmath,amssymb,amsthm,mathtools}
\usepackage[margin=2.5cm,headheight=14pt]{geometry}
\usepackage{array,enumitem}
\usepackage{tikz}
\usetikzlibrary{arrows.meta,calc,positioning}
\usepackage{placeins}
\usepackage{microtype}
\usepackage{xcolor}
\usepackage[unicode,colorlinks=true,linkcolor=blue!45!black]{hyperref}
\theoremstyle{plain}
\newtheorem{twierdzenie}{Twierdzenie}[section]
\newtheorem{lemat}[twierdzenie]{Lemat}
\newtheorem{stwierdzenie}[twierdzenie]{Stwierdzenie}
\newtheorem{fakt}[twierdzenie]{Fakt}
\newtheorem{wniosek}[twierdzenie]{Wniosek}
\theoremstyle{definition}
\newtheorem{definicja}[twierdzenie]{Definicja}
\newtheorem{przyklad}[twierdzenie]{Przykład}
\theoremstyle{remark}
\newtheorem{uwaga}[twierdzenie]{Uwaga}
\renewcommand{\proofname}{Dowód}
\newcommand{\N}{\mathbb N}
\newcommand{\Z}{\mathbb Z}
\newcommand{\Q}{\mathbb Q}
\newcommand{\R}{\mathbb R}
\newcommand{\C}{\mathbb C}
\newcommand{\id}{\mathrm{id}}
\newcommand{\dd}{\mathrm d}
\setlength{\parindent}{0pt}
\setlength{\parskip}{4pt}
\setlength{\emergencystretch}{2em}
\usepackage{fancyhdr}
\pagestyle{fancy}
\fancyhf{}
\fancyhead[L]{\small Analiza matematyczna II}
\fancyhead[R]{\small\ifnum\value{chapter}>0 Rozdział \thechapter\else Spis treści\fi}
\fancyfoot[C]{\thepage}
\renewcommand{\headrulewidth}{0.4pt}

% BEGIN GLOBAL CODE LISTINGS (generated from latex/pl/code-listings.tex)
\usepackage{xcolor}
\usepackage{listings}
\definecolor{codebg}{RGB}{248,249,251}
\definecolor{codeborder}{RGB}{180,191,203}
\definecolor{codecomment}{RGB}{61,113,79}
\definecolor{codekeyword}{RGB}{41,78,132}
\lstdefinelanguage{Julia}{
  morekeywords={abstract,baremodule,begin,break,catch,const,continue,do,else,elseif,
    end,export,false,finally,for,function,global,if,import,in,isa,let,local,macro,
    module,mutable,primitive,quote,return,struct,true,try,type,using,where,while,
    missing,nothing,println,print},
  sensitive=true,morecomment=[l]{\#},morecomment=[n]{\#=}{=\#},
  morestring=[b]",
}
\lstdefinestyle{skrypt}{
  backgroundcolor=\color{codebg},frame=single,rulecolor=\color{codeborder},
  framesep=5pt,frameround=tttt,framerule=.5pt,
  basicstyle=\ttfamily\footnotesize,keywordstyle=\color{codekeyword}\bfseries,
  commentstyle=\color{codecomment}\itshape,stringstyle=\color{purple!65!black},
  numbers=left,numberstyle=\tiny\color{gray!75!black},numbersep=8pt,
  stepnumber=1,breaklines=true,breakatwhitespace=false,
  columns=fullflexible,keepspaces=true,showstringspaces=false,tabsize=4,
  linewidth=.90\linewidth,captionpos=b,aboveskip=11pt,belowskip=12pt,
  xleftmargin=12pt,xrightmargin=2pt,
}
\lstdefinestyle{pythonpde}{style=skrypt,language=Python}
\lstdefinestyle{python}{style=skrypt,language=Python}
\lstdefinestyle{julia}{style=skrypt,language=Julia}
\lstset{style=skrypt}
\renewcommand{\lstlistingname}{Kod}
\renewcommand{\lstlistlistingname}{Spis kodów}
% END GLOBAL CODE LISTINGS
\begin{document}
\frontmatter
\hypersetup{pageanchor=false}
\begin{titlepage}
\centering
\vspace*{4cm}
{\Huge\bfseries Analiza matematyczna\par}
\vspace{.7cm}
{\Large Tom II\par}
\vspace{.7cm}
{\large Szeregi, rachunek różniczkowy i całkowy, analiza Fouriera\par}
\vfill
{\normalsize Biblioteka matematyczna\par}
\end{titlepage}
\hypersetup{pageanchor=true}
\begingroup
\footnotesize\linespread{0.90}\selectfont
\setlength{\parskip}{0pt}
\tableofcontents
\endgroup
\mainmatter
% Adaptacja AF-12 i AM I-07; oryginalne etykiety zachowane.
\chapter{Przestrzenie unormowane i przestrzenie Banacha}
\label{ma2:ch:normed}
\providecommand{\K}{\mathbb K}
Niech $\K\in\{\R,\C\}$. Zakładamy znajomość elementarnej
algebry liniowej, ciągów i zupełności metrycznej z tomu I.
\section{Normy i operatory}
\begin{definicja}[Przestrzeń unormowana i przestrzeń Banacha]
\label{def:banach12}
Przestrzenią unormowaną nazywamy parę $(E,\|\cdot\|)$, gdzie $E$
jest przestrzenią wektorową nad $\K$, a norma spełnia
$\|x\|=0\iff x=0$, $\|\alpha x\|=|\alpha|\|x\|$ oraz
$\|x+y\|\le\|x\|+\|y\|$. Przestrzeń unormowaną nazywamy
\emph{przestrzenią Banacha}, jeżeli jest zupełna względem metryki
$d(x,y)=\|x-y\|$.
\end{definicja}

\begin{stwierdzenie}[Norma operatora i przestrzeń sprzężona]
\label{prop:operator-norm12}
Niech $E$ będzie przestrzenią unormowaną i $F$ przestrzenią Banacha.
Wtedy przestrzeń $\mathcal L(E,F)$ wszystkich ciągłych operatorów
liniowych $E\to F$, z normą
$\|T\|=\sup_{\|x\|_E\le1}\|Tx\|_F$, jest przestrzenią Banacha.
W szczególności $E'=\mathcal L(E,\K)$ jest zawsze przestrzenią
Banacha, nawet gdy $E$ nie jest zupełna.
\end{stwierdzenie}
\begin{proof}
Ciągłość liniowego $T$ w zerze jest równoważna istnieniu $C\ge0$
takiego, że $\|Tx\|_F\le C\|x\|_E$: dla implikacji nietrywialnej
wystarczy przeskalować niezerowe $x$ do kuli, na której $T$ jest
ograniczony. Norma operatorowa jest zatem skończona i spełnia
aksjomaty normy.

Jeśli $(T_n)$ jest ciągiem Cauchy'ego w normie operatorowej,
to dla każdego $x\in E$ ciąg $(T_nx)$ jest Cauchy'ego w $F$, ponieważ
$\|T_nx-T_mx\|\le\|T_n-T_m\|\|x\|$.
Niech $Tx=\lim_nT_nx$. Granice zachowują liniowość. Ustalmy
$\varepsilon>0$ i wybierzmy $N$ tak, że
$\|T_n-T_m\|<\varepsilon$ dla $m,n\ge N$.
Po przejściu do granicy $m\to\infty$ otrzymujemy
$\|(T_n-T)x\|\le\varepsilon\|x\|$ dla $n\ge N$.
Stąd $T$ jest ciągły i $\|T_n-T\|\to0$.
\end{proof}

\section{Przestrzenie ciągowe}

\begin{definicja}[Przestrzenie $\ell^p$, $c_0$ i $\ell^\infty$]
\label{def:sequence-spaces12}
Dla $1\le p<\infty$ określamy
$\ell^p=\{x=(x_n)_{n\ge1}\in\K^{\N}:\sum_{n\ge1}|x_n|^p<\infty\}$
i $\|x\|_p=(\sum_{n\ge1}|x_n|^p)^{1/p}$.
Ponadto $\ell^\infty$ jest przestrzenią ciągów ograniczonych
z normą $\|x\|_\infty=\sup_{n\ge1}|x_n|$, a $c_0$ oznacza
podprzestrzeń ciągów zbieżnych do zera, z tą samą normą.
Przez $c_{00}$ oznaczamy przestrzeń ciągów o skończonym nośniku.
\end{definicja}

\begin{lemat}[Nierówność H\"oldera dla ciągów]
\label{lem:holder-seq12}
Dla $1<p<\infty$, $q=p/(p-1)$, $x\in\ell^p$ i $y\in\ell^q$
zachodzi
$\sum_{n\ge1}|x_ny_n|\le\|x\|_p\|y\|_q$.
Dla $p=1$, $q=\infty$, ta sama nierówność jest prawdziwa.
\end{lemat}
\begin{proof}
Dla $a,b\ge0$ nierówność Younga
$ab\le a^p/p+b^q/q$ wynika z wypukłości funkcji
$t\mapsto e^t$ (równoważnie, z minimalizacji
$t^p/p-tb$ względem $t\ge0$). Jeżeli obie normy są niezerowe,
stosujemy ją do $a=|x_n|/\|x\|_p$ i
$b=|y_n|/\|y\|_q$, sumujemy po $n\le N$,
a następnie przechodzimy do granicy $N\to\infty$.
Jeśli jedna norma jest zerowa, wszystkie współrzędne tego ciągu są zerowe, więc obie strony nierówności są zerowe. Dla $p=1$
oszacowujemy bezpośrednio każdy wyraz przez
$|x_n|\|y\|_\infty$.
\end{proof}

\begin{twierdzenie}[Zupełność klasycznych przestrzeni ciągowych]
\label{thm:seq-banach12}
Dla każdego $1\le p\le\infty$ przestrzeń $\ell^p$ jest Banacha.
Przestrzeń $c_0$ jest domkniętą podprzestrzenią $\ell^\infty$,
a zatem także przestrzenią Banacha. Ponadto $c_{00}$ jest gęste
w $\ell^p$ dla $1\le p<\infty$ i w $c_0$, lecz nie jest gęste
w $\ell^\infty$.
\end{twierdzenie}
\begin{proof}
Dla $p=\infty$ norma supremum jest normą: nieujemność i
rozdzielanie punktów wynikają z modułu, jednorodność z
$|\alpha x_j|=|\alpha||x_j|$, a nierówność trójkąta z
$|x_j+y_j|\le\|x\|_\infty+\|y\|_\infty$ dla każdego $j$.
Dla $p<\infty$ jednorodność i rozdzielanie punktów wynikają
bezpośrednio z definicji. Dla $p=1$ nierówność trójkąta
otrzymujemy przez zsumowanie nierówności dla modułu.
Dla $p>1$ wypukłość $t\mapsto t^p$ na $[0,\infty)$ daje
$|x_j+y_j|^p\le2^{p-1}(|x_j|^p+|y_j|^p)$, więc $z=x+y\in\ell^p$.
Niech $q=p/(p-1)$ i $w_j=|z_j|^{p-1}$.
Ponieważ $(p-1)q=p$, mamy $w\in\ell^q$ i
$\|w\|_q=\|z\|_p^{p-1}$. Z lematu \ref{lem:holder-seq12},
\[
 \|z\|_p^p\le\sum_j|x_j|w_j+\sum_j|y_j|w_j
 \le(\|x\|_p+\|y\|_p)\|z\|_p^{p-1}.
\]
Gdy $\|z\|_p=0$, nierówność trójkąta wynika z nieujemności
prawej strony; inaczej dzielimy przez $\|z\|_p^{p-1}>0$.
Wykazaliśmy również domknięcie na dodawanie; domknięcie na
mnożenie przez skalar wynika z jednorodności.
Niech $(x^{(n)})$ będzie ciągiem Cauchy'ego w $\ell^p$.
Dla każdego $j$ zachodzi
$|x_j^{(n)}-x_j^{(m)}|\le\|x^{(n)}-x^{(m)}\|_p$.
Z zupełności $\K$ istnieje $x_j=\lim_m x_j^{(m)}$.
Ustalmy $\varepsilon>0$ i $N$ tak, że norma różnicy jest
mniejsza niż $\varepsilon/2$ dla $m,n\ge N$.
Dla $p<\infty$ i dowolnego $J\ge1$ mamy
\[
 \sum_{j=1}^J|x_j^{(n)}-x_j^{(m)}|^p\le(\varepsilon/2)^p.
\]
Przechodząc do granicy $m\to\infty$ w skończonej sumie,
a następnie do granicy niemalejących sum po $J$, dostajemy
$x^{(n)}-x\in\ell^p$ oraz
$\|x^{(n)}-x\|_p\le\varepsilon/2$ dla $n\ge N$.
Dla $p=\infty$ przechodzimy do granicy w każdej nierówności
$|x_j^{(n)}-x_j^{(m)}|\le\varepsilon/2$, po czym bierzemy
supremum po $j$; otrzymujemy identyczne oszacowanie.
W obu przypadkach $x=x^{(N)}-(x^{(N)}-x)\in\ell^p$,
a oszacowanie daje zbieżność w normie.
Jeśli $x^{(n)}\in c_0$ i $x^{(n)}\to x$ w normie supremum,
wybieramy $n$ z $\|x-x^{(n)}\|_\infty<\varepsilon/2$
oraz $J$ z $|x_j^{(n)}|<\varepsilon/2$ dla $j\ge J$.
Wówczas $|x_j|<\varepsilon$ dla $j\ge J$, więc $x\in c_0$.
Z rachunku granic skalarnych $c_0$ jest podprzestrzenią liniową.
Jako domknięta podprzestrzeń zupełnej przestrzeni jest zupełna.
Obcięcie $x$ po $J$ współrzędnych daje $x^{[J]}\in c_{00}$.
Dla $p<\infty$ norma różnicy wynosi
$(\sum_{j>J}|x_j|^p)^{1/p}\to0$, a dla $x\in c_0$
wynosi $\sup_{j>J}|x_j|\to0$. Dla $u=(1,1,\ldots)$
i każdego $v\in c_{00}$ pewna współrzędna $v_j$ jest zerowa,
więc $\|u-v\|_\infty\ge1$, co wyklucza gęstość w $\ell^\infty$.
\end{proof}



\section{Metryka, równoważność norm i zupełność}
\begin{twierdzenie}
\label{ma2:thm:metric}
Norma wyznacza metrykę $d(x,y)=\|x-y\|$. Norma jest ciągła,
a dodawanie i mnożenie przez skalar są ciągłe w topologiach
produktowych.
\end{twierdzenie}
\begin{proof}
Z $0=\|x+(-x)\|\le2\|x\|$ wynika nieujemność normy.
Rozdzielanie punktów i symetria wynikają z aksjomatów normy,
w tym jednorodności dla skalara $-1$. Nierówność trójkąta
dla metryki wynika z $x-z=(x-y)+(y-z)$.
Z $\|x\|\le\|x-y\|+\|y\|$ oraz nierówności z zamienionymi
$x,y$ otrzymujemy $|\|x\|-\|y\||\le\|x-y\|$.
Jeśli $x_n\to x$, $y_n\to y$, to
$\|x_n+y_n-x-y\|\le\|x_n-x\|+\|y_n-y\|\to0$.
Jeżeli $\alpha_n\to\alpha$, ciąg $(|\alpha_n|)$ jest
ograniczony przez $M$ i
$\|\alpha_n x_n-\alpha x\|\le M\|x_n-x\|+|\alpha_n-\alpha|\|x\|\to0$.
Kryterium ciągowe w przestrzeniach metrycznych daje ciągłość działań.
\end{proof}
\begin{twierdzenie}[Równoważność norm]
\label{ma2:thm:equiv}
Jeżeli dla norm $p,q$ na $E$ istnieją $a,b>0$ z
$a p(x)\le q(x)\le b p(x)$ dla każdego $x\in E$, to normy
wyznaczają te same topologie, zbieżności i ciągi Cauchy'ego.
Zupełność jednej jest równoważna zupełności drugiej.
\end{twierdzenie}
\begin{proof}
Inkluzje $B_p(x,r/b)\subseteq B_q(x,r)$ i
$B_q(x,ar)\subseteq B_p(x,r)$ dają równość topologii.
Stosując oszacowania do $x_n-x$ oraz $x_n-x_m$,
otrzymujemy równoważność odpowiednio warunków zbieżności
i warunków Cauchy'ego. Jeśli jedna norma jest zupełna,
każdy ciąg Cauchy'ego w drugiej ma granicę w pierwszej,
a zatem także w drugiej.
\end{proof}
\begin{twierdzenie}[Wymiar skończony]
\label{ma2:thm:finite}
W skończenie wymiarowej przestrzeni wszystkie normy są
równoważne i każda z nich jest zupełna.
\end{twierdzenie}
\begin{proof}
Dla przestrzeni zerowej tezy wynikają z istnienia wyłącznie
wektora zerowego. Dla wymiaru $d\ge1$ wybieramy bazę
$(e_j)_{j=1}^d$ i izomorfizm $Tz=\sum_j z_je_j$ z $\K^d$.
Dla dowolnej normy $p$ mamy $p(Tz)\le C\max_j|z_j|$,
gdzie $C=\sum_jp(e_j)>0$. Z odwrotnej nierówności trójkąta
funkcja $z\mapsto p(Tz)$ jest ciągła względem normy maksimum.
Na zwartej sferze $\max_j|z_j|=1$ (twierdzenie Heinego--Borela)
osiąga minimum $c>0$, bo $T$ jest różnowartościowe.
Skalowanie każdego niezerowego $z$ do tej sfery daje
$c\max_j|z_j|\le p(Tz)$; dla $z=0$ nierówność też zachodzi.
Porównując dwie dowolne normy z normą maksimum, dostajemy
ich równoważność. Jeśli $(Tz_n)$ jest ciągiem Cauchy'ego,
każda współrzędna $z_n$ jest ciągiem Cauchy'ego w $\K$.
Istnieje granica współrzędnościowa $z\in\K^d$; skończoność
liczby współrzędnych daje $\max_j|(z_n-z)_j|\to0$.
Górne oszacowanie daje $p(Tz_n-Tz)\to0$.
\end{proof}
\begin{przyklad}[Niezupełna przestrzeń unormowana]
Przestrzeń $c_{00}$ z normą supremum nie jest zupełna.
Ciąg $x^{(n)}=(1,1/2,\ldots,1/n,0,\ldots)$ jest ciągiem
Cauchy'ego, gdyż dla $m>n$ norma różnicy wynosi $1/(n+1)$.
W $\ell^\infty$ zbiega do $(1/j)_{j\ge1}$, ponieważ norma
różnicy od tego ciągu wynosi $1/(n+1)$. Jednoznaczność
granicy wyklucza granicę w $c_{00}$.
\end{przyklad}

% Adaptacja merytoryczna AM I-09; nowe etykiety dla twierdzeń wektorowych.
\chapter{Szeregi w przestrzeniach Banacha}
\label{ma2:ch:banach-series}

W tym rozdziale $E$ jest przestrzenią unormowaną nad
$\K\in\{\R,\C\}$. Zupełność będziemy zakładać wyraźnie tam,
gdzie jest potrzebna. Wykorzystujemy rachunek szeregów liczbowych
z tomu I (rozdział 09) i wyniki rozdziału \ref{ma2:ch:normed}.
W szczególności zbieżność będzie rozumiana w normie, a nie
względem poszczególnych współrzędnych.

\section{Sumy częściowe i warunek Cauchy'ego}
\begin{definicja}[Szereg wektorowy]
\label{ma2:def:vector-series}
Dla ciągu $(x_n)_{n\ge1}$ w $E$ określamy sumy częściowe
$S_0=0$ i $S_N=\sum_{n=1}^N x_n$ dla $N\ge1$.
Szereg $\sum_{n=1}^\infty x_n$ nazywamy zbieżnym w $E$,
jeżeli istnieje $S\in E$ takie, że $\|S_N-S\|\to0$.
Piszemy wówczas $\sum_{n=1}^\infty x_n=S$.
W przeciwnym razie szereg jest rozbieżny.
\end{definicja}
Suma jest jednoznaczna: gdyby $\|S_N-S\|\to0$ i
$\|S_N-T\|\to0$, nierówność
$\|S-T\|\le\|S-S_N\|+\|S_N-T\|$ dawałaby $\|S-T\|=0$,
czyli $S=T$. Dla szeregu od indeksu $m\ge0$ stosujemy
analogiczną definicję, a sumę pustą przyjmujemy równą $0$.

\begin{twierdzenie}[Warunek konieczny i kryterium Cauchy'ego]
\label{ma2:thm:series-cauchy}
Jeżeli szereg $\sum_{n\ge1}x_n$ jest zbieżny w $E$, to
$x_n\to0$ i dla każdego $\varepsilon>0$ istnieje $N\ge1$
takie, że
\[
 \left\|\sum_{n=p}^{q}x_n\right\|<\varepsilon
 \quad\text{dla wszystkich całkowitych }q\ge p\ge N.
\]
Jeżeli $E$ jest przestrzenią Banacha, powyższy warunek
dotyczący sum skończonych jest również wystarczający.
\end{twierdzenie}
\begin{proof}
Niech $S_N\to S$. Dla $n\ge1$ mamy $x_n=S_n-S_{n-1}$,
więc
$\|x_n\|\le\|S_n-S\|+\|S_{n-1}-S\|\to0$.
Dla $\varepsilon>0$ istnieje $M$ takie, że
$\|S_j-S\|<\varepsilon/2$ dla $j\ge M$.
Jeśli $q\ge p\ge M+1$, to
\[
 \left\|\sum_{n=p}^q x_n\right\|=\|S_q-S_{p-1}\|
 \le\|S_q-S\|+\|S_{p-1}-S\|<\varepsilon.
\]
Odwrotnie, niech warunek sum skończonych zachodzi.
Dla $\varepsilon>0$ wybieramy odpowiadające mu $N$.
Jeśli $i>j\ge N$, to
$\|S_i-S_j\|=\|\sum_{n=j+1}^i x_n\|<\varepsilon$,
bo $i\ge j+1\ge N$. Dla $i<j$ korzystamy z symetrii
normy różnicy, a dla $i=j$ różnica jest zerowa.
Ciąg $(S_n)$ jest zatem ciągiem Cauchy'ego w $E$.
Jeśli $E$ jest zupełna, istnieje $S\in E$ z $S_n\to S$,
co jest definicją zbieżności szeregu.
\end{proof}

\begin{przyklad}[Wyrazy dążące do zera nie wystarczają]
\label{ma2:ex:vector-harmonic}
Jeżeli $v\in E$ i $v\ne0$, szereg $\sum_{n\ge1}v/n$
jest rozbieżny, chociaż $\|v/n\|=\|v\|/n\to0$.
Istotnie dla każdego $N\ge1$ mamy
\[
 \left\|\sum_{n=N+1}^{2N}\frac vn\right\|
 =\|v\|\sum_{n=N+1}^{2N}\frac1n\ge\frac{\|v\|}{2}.
\]
Wybierając $\varepsilon=\|v\|/2>0$, widzimy, że dla
każdego proponowanego indeksu w kryterium Cauchy'ego istnieje
ogon skończony o normie co najmniej $\varepsilon$.
Z twierdzenia \ref{ma2:thm:series-cauchy} szereg jest rozbieżny.
\end{przyklad}

\section{Rachunek sum i operatory ciągłe}
\begin{twierdzenie}[Pominięcie skończonego początku i liniowość]
\label{ma2:thm:series-linear-tail}
Dla każdego całkowitego $m\ge1$ szereg $\sum_{n\ge1}x_n$
jest zbieżny wtedy i tylko wtedy, gdy zbieżny jest
$\sum_{n\ge m}x_n$. W przypadku zbieżności
\[
 \sum_{n=1}^\infty x_n=\sum_{n=1}^{m-1}x_n+
                       \sum_{n=m}^\infty x_n.
\]
Jeśli szeregi $\sum x_n$ i $\sum y_n$ w $E$ są zbieżne
oraz $\alpha,\beta\in\K$, to szereg $\sum(\alpha x_n+\beta y_n)$
jest zbieżny i
\[
 \sum_{n=1}^\infty(\alpha x_n+\beta y_n)
 =\alpha\sum_{n=1}^\infty x_n+\beta\sum_{n=1}^\infty y_n.
\]
\end{twierdzenie}
\begin{proof}
Dla $N\ge m$ suma od $m$ do $N$ jest równa $S_N-S_{m-1}$.
Jeśli $S_N\to S$, to
$\|(S_N-S_{m-1})-(S-S_{m-1})\|=\|S_N-S\|\to0$.
Odwrotnie, jeśli sumy od $m$ do $N$ dążą do $A$, to
$\|S_N-(A+S_{m-1})\|=\|(S_N-S_{m-1})-A\|\to0$.
Daje to obie implikacje i wzór, również dla $m=1$,
gdy $S_0=0$.
Niech $S_N\to S$ i $T_N=\sum_{n=1}^N y_n\to T$.
Suma częściowa kombinacji liniowej wynosi
$\alpha S_N+\beta T_N$, a jej odległość od
$\alpha S+\beta T$ jest nie większa niż
$|\alpha|\|S_N-S\|+|\beta|\|T_N-T\|\to0$.
Nie potrzeba tu zupełności $E$, ponieważ obie granice
istnieją z założenia.
\end{proof}

\begin{twierdzenie}[Przenoszenie sumy przez operator]
\label{ma2:thm:series-operator}
Niech $F$ będzie przestrzenią unormowaną nad $\K$,
a $T:E\to F$ ciągłym operatorem liniowym.
Jeśli $\sum x_n=S$ w $E$, to $\sum Tx_n=TS$ w $F$.
\end{twierdzenie}
\begin{proof}
Z liniowości $\sum_{n=1}^N Tx_n=TS_N$.
Z twierdzenia \ref{prop:operator-norm12} mamy
$\|TS_N-TS\|_F\le\|T\|\|S_N-S\|_E\to0$.
Granica $TS$ należy do $F$ niezależnie od zupełności $F$.
\end{proof}

\begin{twierdzenie}[Współrzędne w skończonym wymiarze]
\label{ma2:thm:series-coordinates}
Niech $E$ będzie skończenie wymiarowa, z bazą
$(e_j)_{j=1}^d$, gdzie $d\ge0$, i niech
$x_n=\sum_{j=1}^d a_{nj}e_j$, $a_{nj}\in\K$.
Szereg $\sum_n x_n$ jest zbieżny wtedy i tylko wtedy,
gdy każdy z $d$ szeregów skalarnych $\sum_n a_{nj}$
jest zbieżny. Wtedy
\[
 \sum_{n=1}^\infty x_n
 =\sum_{j=1}^d\left(\sum_{n=1}^\infty a_{nj}\right)e_j.
\]
\end{twierdzenie}
\begin{proof}
Dla $d=0$ każdy wyraz i każda suma są zerowe, a warunek
zbieżności szeregów współrzędnych jest spełniony w sposób pusty.
Załóżmy więc $d\ge1$.
Z dowodu twierdzenia \ref{ma2:thm:finite} istnieją $c,C>0$
takie, że dla każdego $z\in\K^d$ zachodzi
$c\max_j|z_j|\le\|\sum_j z_je_j\|\le C\max_j|z_j|$.
Oznaczmy $s_{Nj}=\sum_{n=1}^N a_{nj}$.
Z rozdzielności skończonych sum $S_N=\sum_j s_{Nj}e_j$.
Jeśli $S_N\to S=\sum_j s_je_j$, dolne oszacowanie daje
$|s_{Nj}-s_j|\le\|S_N-S\|/c\to0$ dla każdego $j$.
Odwrotnie, jeśli $s_{Nj}\to s_j$ dla wszystkich $j$, to
dla $\varepsilon>0$ dobieramy dla każdego $j$ indeks $N_j$
z $|s_{Nj}-s_j|<\varepsilon/C$ dla $N\ge N_j$.
Dla $N\ge\max_j N_j$ górne oszacowanie daje
$\|S_N-\sum_j s_je_j\|<\varepsilon$.
To dowodzi zbieżności i wzoru.
\end{proof}

\section{Szeregi geometryczne, teleskopowe i grupowanie}
\begin{twierdzenie}[Szereg geometryczny pomnożony przez wektor]
\label{ma2:thm:vector-geometric}
Niech $v\in E$ i $q\in\K$, przy czym $q^0=1$ również dla $q=0$.
Jeżeli $|q|<1$, to
\[
 \sum_{n=0}^{\infty}q^n v=\frac{1}{1-q}v.
\]
Jeśli $v\ne0$, warunek $|q|<1$ jest także konieczny.
Dla $v=0$ szereg jest zbieżny dla każdego $q$.
\end{twierdzenie}
\begin{proof}
Dla $|q|<1$ mamy $q\ne1$; tożsamość skończonego szeregu
geometrycznego daje
$\sum_{n=0}^N q^n v=(1-q^{N+1})v/(1-q)$.
Odległość od $v/(1-q)$ wynosi
$|q|^{N+1}\|v\|/|1-q|$, co dąży do zera, ponieważ
$0\le |q|<1$. Nie używamy zupełności $E$.
Jeśli $v\ne0$ i $|q|\ge1$, normy wyrazów
$|q|^n\|v\|$ są nie mniejsze niż $\|v\|>0$.
Warunek konieczny z twierdzenia \ref{ma2:thm:series-cauchy}
(po pominięciu wyrazu o indeksie $0$) nie jest spełniony.
Jeśli $v=0$, wszystkie sumy częściowe są zerowe.
\end{proof}

\begin{twierdzenie}[Szereg teleskopowy]
\label{ma2:thm:vector-telescoping}
Dla ciągu $(u_n)$ w $E$ szereg $\sum_{n\ge1}(u_n-u_{n+1})$
jest zbieżny wtedy i tylko wtedy, gdy $(u_n)$ jest zbieżny.
Jeśli $u_n\to u$, suma wynosi $u_1-u$.
\end{twierdzenie}
\begin{proof}
Dla $N\ge1$ sumy pośrednich wyrazów znoszą się, więc
$\sum_{n=1}^N(u_n-u_{n+1})=u_1-u_{N+1}$.
Jeśli $u_n\to u$, norma różnicy tej sumy od $u_1-u$
wynosi $\|u_{N+1}-u\|\to0$.
Jeśli natomiast suma częściowa dąży do $S$, to
$\|u_{N+1}-(u_1-S)\|=\|S-(u_1-u_{N+1})\|\to0$.
Pominięcie pierwszego wyrazu nie zmienia granicy ciągu,
zatem $u_n\to u_1-S$.
\end{proof}

\begin{twierdzenie}[Grupowanie kolejnych wyrazów]
\label{ma2:thm:series-grouping}
Niech $0=k_0<k_1<k_2<\cdots$ będą liczbami całkowitymi
i niech $b_j=\sum_{n=k_{j-1}+1}^{k_j}x_n$ dla $j\ge1$.
Jeżeli $\sum x_n=S$, to $\sum b_j=S$.
Odwrotnie, jeśli $\sum b_j=S$ oraz
\[
 \max_{k_{j-1}\le m\le k_j}
 \left\|\sum_{n=k_{j-1}+1}^{m}x_n\right\|\longrightarrow0,
\]
to $\sum x_n=S$.
\end{twierdzenie}
\begin{proof}
Z rozłączności kolejnych przedziałów indeksów
$\sum_{j=1}^J b_j=\sum_{n=1}^{k_J}x_n=S_{k_J}$.
Ponieważ $k_J\ge J$, jest to podciąg sum częściowych,
a z $S_N\to S$ wynika $S_{k_J}\to S$.
Dla implikacji odwrotnej oznaczmy maksimum z tezy przez $d_j$.
Dla każdego $N\ge1$ istnieje dokładnie jeden $j$ z
$k_{j-1}<N\le k_j$, ponieważ $(k_j)$ jest ściśle rosnący
i nieograniczony. Wtedy
\[
 \|S_N-S\|\le\|S_{k_{j-1}}-S\|+
              \|S_N-S_{k_{j-1}}\|
 \le\|S_{k_{j-1}}-S\|+d_j.
\]
Dla $\varepsilon>0$ wybieramy $J$ tak duże, że dla $j\ge J$
obie wielkości po prawej stronie są mniejsze niż
$\varepsilon/2$ (pierwsza z założenia zbieżności szeregu
zgrupowanego, druga z $d_j\to0$).
Dla $N>k_{J-1}$ odpowiedni indeks $j$ jest co najmniej $J$,
więc $\|S_N-S\|<\varepsilon$.
\end{proof}
\begin{przyklad}[Grupowanie nie musi zachowywać rozbieżności]
\label{ma2:ex:grouping-divergence}
Dla $v\ne0$ połóżmy $x_{2j-1}=v$, $x_{2j}=-v$.
Sumy częściowe o indeksach parzystych wynoszą $0$,
a o nieparzystych $v$. Ponieważ $v\ne0$, oba podciągi
mają różne granice, więc szereg nie jest zbieżny.
Przy $k_j=2j$ każdy wyraz $b_j$ jest zerowy,
a szereg zgrupowany jest zbieżny do $0$.
Warunek maksimum w twierdzeniu \ref{ma2:thm:series-grouping}
nie zachodzi, ponieważ w każdym bloku maksimum wynosi $\|v\|$.
\end{przyklad}

\section{Zbieżność bezwzględna i oszacowania reszty}
\begin{definicja}[Zbieżność bezwzględna]
\label{ma2:def:absolute-series}
Szereg $\sum x_n$ nazywamy bezwzględnie zbieżnym, jeżeli
zbieżny jest szereg liczbowy $\sum\|x_n\|$.
Sama definicja nie wymaga zupełności $E$.
\end{definicja}
\begin{twierdzenie}
\label{ma2:thm:absolute-series}
Jeżeli $E$ jest przestrzenią Banacha i $\sum\|x_n\|<\infty$,
to $\sum x_n$ jest zbieżny. Dla jego sumy $S$ i każdego
$N\ge0$ zachodzi oszacowanie
\[
 \|S-S_N\|\le\sum_{n=N+1}^\infty\|x_n\|.
\]
\end{twierdzenie}
\begin{proof}
Dla $q\ge p$ nierówność
$\|\sum_{n=p}^q x_n\|\le\sum_{n=p}^q\|x_n\|$
wynika przez indukcję: dla jednego składnika mamy równość,
a dołączenie następnego wyrazu zwiększa prawą stronę o jego normę
na mocy nierówności trójkąta.
Dla $\varepsilon>0$ kryterium Cauchy'ego dla zbieżnego
szeregu liczbowego $\sum\|x_n\|$ daje $M$, dla którego
$\sum_{n=p}^q\|x_n\|<\varepsilon$ przy $q\ge p\ge M$.
Stąd spełniony jest warunek z twierdzenia
\ref{ma2:thm:series-cauchy}. Zupełność $E$ daje istnienie $S$.
Dla $L>N$ mamy
$\|S_L-S_N\|\le\sum_{n=N+1}^L\|x_n\|$.
Przy ustalonym $N$ lewa strona dąży do $\|S-S_N\|$
z ciągłości normy (twierdzenie \ref{ma2:thm:metric}),
a prawa do sumy ogona zbieżnego szeregu liczbowego.
Przechodząc do granicy w nierówności nieostrej, dostajemy tezę.
\end{proof}

\begin{twierdzenie}[Charakteryzacja zupełności]
\label{ma2:thm:series-completeness}
Przestrzeń unormowana $E$ jest przestrzenią Banacha wtedy
i tylko wtedy, gdy każdy szereg bezwzględnie zbieżny w $E$
jest zbieżny w $E$.
\end{twierdzenie}
\begin{proof}
Pierwsza implikacja wynika z twierdzenia \ref{ma2:thm:absolute-series}.
Załóżmy własność zbieżności wszystkich takich szeregów
i niech $(y_n)$ będzie dowolnym ciągiem Cauchy'ego w $E$.
Dla każdego $j\ge1$ istnieje $M_j\ge1$ takie, że
$\|y_n-y_m\|<2^{-j}$ dla $n,m\ge M_j$.
Możemy wybrać $M_j$ jako najmniejszy dodatni indeks
o podanej własności. Wybieramy $n_1=M_1$ i rekurencyjnie
$n_{j+1}=\max\{n_j+1,M_{j+1},M_j\}$.
Z definicji $n_{j+1}>n_j$ i $n_{j+1}\ge M_{j+1}$.
Ponieważ $n_1=M_1$, indukcyjnie $n_j\ge M_j$ dla każdego $j$.
Ponadto $n_{j+1}\ge M_j$, zatem $n_j,n_{j+1}\ge M_j$,
więc dla $h_j=y_{n_{j+1}}-y_{n_j}$ mamy $\|h_j\|<2^{-j}$.
Szereg $\sum\|h_j\|$ jest zbieżny przez porównanie
z szeregiem geometrycznym. Z założenia istnieje
$h=\sum_{j=1}^\infty h_j$ w $E$.
Suma teleskopowa daje $y_{n_{J+1}}=y_{n_1}+\sum_{j=1}^J h_j$,
a zatem $y_{n_{J+1}}\to y=y_{n_1}+h$.
Ustalmy $\varepsilon>0$. Wybieramy $M$ z
$\|y_n-y_m\|<\varepsilon/2$ dla $n,m\ge M$.
Ponieważ $n_{J+1}\to\infty$ i $y_{n_{J+1}}\to y$,
wybieramy $J$ z $n_{J+1}\ge M$ oraz
$\|y_{n_{J+1}}-y\|<\varepsilon/2$.
Dla każdego $n\ge M$ dostajemy
$\|y_n-y\|\le\|y_n-y_{n_{J+1}}\|+\|y_{n_{J+1}}-y\|<\varepsilon$.
Dowolny ciąg Cauchy'ego ma więc granicę w $E$.
\end{proof}

\begin{przyklad}[Konieczność zupełności]
\label{ma2:ex:incomplete-series}
W przestrzeni $c_{00}$ z normą supremum niech $e_n$
oznacza ciąg mający $1$ na $n$-tej współrzędnej i zera na
pozostałych. Wtedy $x_n=2^{-n}e_n$ należy do $c_{00}$ i
$\sum\|x_n\|_\infty=\sum 2^{-n}=1$.
Sumy częściowe $S_N$ mają w $\ell^\infty$ granicę
$x=(2^{-j})_{j\ge1}$, ponieważ
$\|S_N-x\|_\infty=\sup_{j>N}2^{-j}=2^{-(N+1)}\to0$.
Gdyby szereg był zbieżny w $c_{00}$, granica byłaby
również jego granicą w $\ell^\infty$, gdzie norma jest ta sama.
Jednoznaczność dawałaby $x\in c_{00}$, co jest niemożliwe,
ponieważ $x_j>0$ dla każdego $j$.
\end{przyklad}

\section{Kryteria porównawcze, ilorazowe i pierwiastkowe}
\begin{wniosek}[Majoranta liczbowa]
\label{ma2:cor:series-majorant}
Niech $E$ będzie przestrzenią Banacha. Jeśli $a_n\ge0$,
$\sum a_n<\infty$ i od pewnego indeksu $\|x_n\|\le a_n$,
to $\sum x_n$ jest bezwzględnie zbieżny i zbieżny.
Jeśli dla całkowitego $N\ge0$ nierówność zachodzi dla wszystkich $n>N$, to
$\|S-S_N\|\le\sum_{n>N}a_n$.
\end{wniosek}
\begin{proof}
Wybierzmy $m\ge1$ tak, aby $\|x_n\|\le a_n$ dla $n\ge m$.
Dla $L\ge m$ mamy
\[
 \sum_{n=1}^L\|x_n\|\le
 \sum_{n=1}^{m-1}\|x_n\|+\sum_{n=m}^L a_n
 \le \sum_{n=1}^{m-1}\|x_n\|+\sum_{n=1}^\infty a_n.
\]
To samo skończone ograniczenie działa dla $L<m$,
bo normy są nieujemne. Sumy częściowe szeregu norm są
niemalejące i ograniczone z góry, więc są zbieżne
na mocy twierdzenia o zbieżności ciągu monotonicznego
z tomu I. Z twierdzenia \ref{ma2:thm:absolute-series}
istnieje suma $S$ szeregu wektorowego.
Jeżeli $\|x_n\|\le a_n$ dla wszystkich $n>N$, to dla $L>N$
mamy $\sum_{n=N+1}^L\|x_n\|\le\sum_{n=N+1}^L a_n$.
Obie strony mają skończone granice przy $L\to\infty$,
więc oszacowanie reszty z twierdzenia
\ref{ma2:thm:absolute-series} daje
$\|S-S_N\|\le\sum_{n=N+1}^\infty\|x_n\|
\le\sum_{n=N+1}^\infty a_n$.
\end{proof}

\begin{twierdzenie}[Kryteria d'Alemberta i Cauchy'ego]
\label{ma2:thm:vector-ratio-root}
Niech $E$ będzie przestrzenią Banacha.
\begin{enumerate}[label=(\roman*)]
\item Jeżeli od pewnego indeksu $x_n\ne0$, to dla
$r_n=\|x_{n+1}\|/\|x_n\|$ warunek $\limsup r_n<1$
daje zbieżność bezwzględną i zbieżność szeregu.
Warunek $\liminf r_n>1$ daje rozbieżność.
\item Dla $L=\limsup_{n\to\infty}\|x_n\|^{1/n}$
warunek $L<1$ daje zbieżność bezwzględną i zbieżność,
a $L>1$ daje rozbieżność. W tym przypadku dopuszczamy
wyrazy zerowe.
\end{enumerate}
\end{twierdzenie}
\begin{proof}
Granice dolną i górną rozumiemy w rozszerzonej prostej
rzeczywistej, zgodnie z tomem I (rozdział 08). Korzystamy
z kryterium progowego: jeżeli $\limsup b_n<q$, to od
pewnego indeksu $b_n<q$; jeżeli $q<\liminf b_n$, to od
pewnego indeksu $q<b_n$; jeżeli $q<\limsup b_n$, to
dla każdego $M$ istnieje $n\ge M$ z $q<b_n$.
W części (i), jeśli $\limsup r_n<1$, wybieramy
$0<q<1$ z $\limsup r_n<q$. Od pewnego $m$ mamy
$\|x_{n+1}\|\le q\|x_n\|$ dla $n\ge m$.
Indukcyjnie $\|x_{m+k}\|\le q^k\|x_m\|$ dla $k\ge0$:
dla $k=0$ mamy równość, a przejście do $k+1$ wynika
z powyższego oszacowania kolejnych norm.
Majoranta geometryczna i wniosek \ref{ma2:cor:series-majorant}
dają zbieżność.
Jeśli $\liminf r_n>1$, wybieramy skończone $q>1$ mniejsze
od granicy dolnej (również gdy jest nieskończona).
Od pewnego $m$ mamy $r_n>q$ i $x_n\ne0$.
Dla $k=0$ mamy równość $\|x_{m+k}\|=\|x_m\|$.
Jeżeli $\|x_{m+k}\|\ge q^k\|x_m\|$, to
$\|x_{m+k+1}\|>q\|x_{m+k}\|\ge q^{k+1}\|x_m\|$.
Zatem indukcyjnie $\|x_{m+k}\|\ge q^k\|x_m\|\ge\|x_m\|>0$
dla każdego $k\ge0$. Z twierdzenia \ref{ma2:thm:series-cauchy}
szereg jest rozbieżny.
W części (ii), jeśli $L<1$, wybieramy $0<q<1$ z $L<q$.
Od pewnego $m$ mamy $\|x_n\|^{1/n}<q$, czyli
$\|x_n\|<q^n$ dla $n\ge m$ (potęgowanie stopnia $n$ zachowuje
porządek na liczbach nieujemnych).
Stosujemy wniosek \ref{ma2:cor:series-majorant}.
Jeśli $L>1$, wybieramy skończone $q>1$ z $q<L$.
Dla każdego $M$ istnieje $n\ge M$ z $\|x_n\|^{1/n}>q$,
stąd $\|x_n\|>q^n\ge1$. Wyrazy nie dążą do zera,
więc szereg jest rozbieżny.
\end{proof}
\begin{przyklad}[Kontrola błędu obcięcia]
\label{ma2:ex:series-geometric-bound}
Jeśli $C\ge0$, $0\le q<1$ i $\|x_n\|\le Cq^n$ dla $n\ge1$,
to w przestrzeni Banacha z wniosku
\ref{ma2:cor:series-majorant} istnieje suma $S$ i
\[
 \|S-S_N\|\le C\sum_{n=N+1}^\infty q^n
              =\frac{Cq^{N+1}}{1-q}\qquad(N\ge0).
\]
Dla $q=0$ wszystkie $x_n$ są zerowe, więc oszacowanie
nadal zachodzi. Dla $C=1$, $q=1/2$ i $N=10$ błąd jest
nie większy niż $2^{-10}<1/1000$.
\end{przyklad}

\begin{uwaga}
Rozbieżność wykazana przez warunek konieczny nie wymaga
zupełności. W tym rozdziale nie utożsamiamy zbieżności
bezwzględnej ze zbieżnością bezwarunkową i nie zmieniamy
kolejności wyrazów. Przestawienia, szeregi warunkowo zbieżne
i szeregi bezwarunkowo zbieżne będą tematem rozdziału 03.
\end{uwaga}

% Wyniki skalarne przejęte z AM I-10 z oryginalnymi etykietami.
% Kryterium sum skończonych rozwija AF-27, lem:net-unconditional27.
\chapter{Zbieżność bezwzględna, warunkowa i bezwarunkowa}
\label{ma2:ch:unconditional-series}

Niech $E$ będzie przestrzenią unormowaną nad $\K=\R$ lub
$\K=\C$. Przez $\N_+$ rozumiemy zbiór dodatnich liczb
naturalnych.
Zbieżność szeregu zawsze oznacza zbieżność jego sum
częściowych w normie. Korzystamy z rozdziału
\ref{ma2:ch:banach-series} oraz teorii skalarnej z tomu I.

\section{Trzy rodzaje zbieżności}
\begin{definicja}
\label{ma2:def:unconditional-series}
Szereg $\sum_{n\ge1}x_n$ nazywamy zbieżnym warunkowo
w sensie zbieżności niebezwzględnej, gdy jest zbieżny,
ale $\sum_{n\ge1}\|x_n\|=+\infty$.
Przestawieniem szeregu nazywamy szereg
$\sum_{k\ge1}x_{\pi(k)}$, gdzie $\pi:\N_+\to\N_+$
jest bijekcją.
Szereg nazywamy zbieżnym bezwarunkowo do $S\in E$,
gdy dla każdej takiej bijekcji szereg przestawiony
jest zbieżny do $S$.
\end{definicja}
W literaturze spotyka się również określenie zbieżności
warunkowej dla szeregu zbieżnego, ale nie bezwarunkowo.
Aby uniknąć niejednoznaczności, będziemy dalej mówić
wprost o szeregu zbieżnym niebezwzględnie albo
zbieżnym, lecz nie bezwarunkowo. W skończonym wymiarze
te dwie własności są równoważne, co wykażemy w
twierdzeniu \ref{ma2:thm:finite-unconditional-absolute}.
W nieskończonym wymiarze mogą się różnić.

\begin{twierdzenie}[Przestawienia szeregów bezwzględnie zbieżnych]
\label{ma2:thm:absolute-rearrangement}
Jeżeli $E$ jest przestrzenią Banacha, a szereg
$\sum x_n$ jest bezwzględnie zbieżny, to jest bezwarunkowo
zbieżny. Każde przestawienie jest również bezwzględnie
zbieżne, a suma norm jest niezależna od przestawienia.
W dowolnej przestrzeni unormowanej ten sam wniosek
zachodzi, jeśli oprócz zbieżności szeregu norm założymy
istnienie sumy $\sum x_n$ w $E$.
\end{twierdzenie}
\begin{proof}
Niech $A=\sum_{n\ge1}\|x_n\|$ oraz $S=\sum_{n\ge1}x_n$;
istnienie $S$ wynika z twierdzenia \ref{ma2:thm:absolute-series}
w przypadku przestrzeni Banacha, a w pozostałym przypadku
jest założeniem. Oszacowanie
$\|S-\sum_{n=1}^N x_n\|\le\sum_{n>N}\|x_n\|$
w obu przypadkach wynika z nierówności trójkąta dla
sum skończonych i przejścia do granicy, jak w dowodzie
twierdzenia \ref{ma2:thm:absolute-series}.
Dla bijekcji $\pi$ niech $F_k=\{\pi(1),\ldots,\pi(k)\}$.
Z iniektywności i nieujemności norm mamy
$\sum_{j=1}^k\|x_{\pi(j)}\|\le
\sum_{n=1}^{\max F_k}\|x_n\|\le A$.
Sumy te są niemalejące, więc mają granicę $A'\le A$.
Dla każdego $N\ge1$ istnieje
$k_0=\max_{1\le n\le N}\pi^{-1}(n)$.
Dla $k\ge k_0$ zachodzi $\{1,\ldots,N\}\subseteq F_k$,
więc $A'\ge\sum_{n=1}^N\|x_n\|$.
Przechodząc do granicy $N\to\infty$, dostajemy $A'\ge A$,
czyli $A'=A$.
Dla $\varepsilon>0$ wybieramy $N$ tak, by
$\sum_{n>N}\|x_n\|<\varepsilon/2$.
Dla $k\ge k_0$ nierówność trójkąta oraz powyższe
oszacowanie reszty dają
\[
 \left\|\sum_{j=1}^k x_{\pi(j)}-S\right\|
 \le\left\|\sum_{n\in F_k\setminus\{1,\ldots,N\}}x_n\right\|
 +\left\|\sum_{n=1}^N x_n-S\right\|
 \le 2\sum_{n>N}\|x_n\|<\varepsilon.
\]
Zatem wszystkie przestawienia mają sumę $S$.
\end{proof}

\section{Szeregi naprzemienne}
\begin{twierdzenie}[Kryterium Leibniza]
\label{ma1:thm:leibniz-series}
Niech $(b_n)_{n\ge1}$ będzie nierosnącym ciągiem nieujemnych liczb rzeczywistych takim, że $b_n\to0$. Wówczas szereg
\[
 \sum_{n=1}^\infty (-1)^{n-1}b_n
\]
jest zbieżny. Jeżeli $S$ oznacza jego sumę, $S_0=0$, a $S_N=\sum_{n=1}^N(-1)^{n-1}b_n$ dla $N\ge1$, to dla każdego całkowitego $N\ge0$ zachodzi
\[
 0\le(-1)^N(S-S_N)\le b_{N+1}.
\]
W szczególności $|S-S_N|\le b_{N+1}$.
\end{twierdzenie}
\begin{proof}
Dla $k\ge0$ mamy $S_{2k+2}-S_{2k}=b_{2k+1}-b_{2k+2}\ge0$, więc ciąg parzystych sum częściowych jest niemalejący. Dla $k\ge0$ mamy również $S_{2k+3}-S_{2k+1}=-b_{2k+2}+b_{2k+3}\le0$, czyli ciąg nieparzystych sum częściowych jest nierosnący. Ponadto $S_{2k+1}-S_{2k}=b_{2k+1}\ge0$. Dla każdego $k\ge0$ otrzymujemy
\[
 0=S_0\le S_{2k}\le S_{2k+1}\le S_1=b_1.
\]
Z twierdzenia o zbieżności ciągów monotonicznych (tom I, rozdział o granicach ciągów) istnieją skończone granice $E=\lim_k S_{2k}$ i $O=\lim_k S_{2k+1}$. Ponieważ $2k+1\to\infty$, z założenia $b_n\to0$ otrzymujemy $S_{2k+1}-S_{2k}=b_{2k+1}\to0$, rachunek granic daje $O-E=0$. Oznaczmy wspólną granicę przez $S$. Dla dowolnego $\varepsilon>0$ istnieje całkowity $K\ge0$ taki, że dla każdego $k\ge K$ zarówno $|S_{2k}-S|<\varepsilon$, jak i $|S_{2k+1}-S|<\varepsilon$. Każdy całkowity indeks $N\ge2K$ ma postać $2k$ lub $2k+1$ z $k\ge K$. Zatem $S_N\to S$.

Monotoniczność i twierdzenie o granicach ciągów monotonicznych dają $S=\sup_{k\ge0}S_{2k}=\inf_{k\ge0}S_{2k+1}$, a więc $S_{2k}\le S\le S_{2k+1}$ dla każdego $k\ge0$. Dla parzystego $N=2k$ wynika stąd $0\le S-S_N\le S_{2k+1}-S_{2k}=b_{N+1}$. Dla nieparzystego $N=2k+1$ korzystamy z $S_{2k+2}\le S\le S_{2k+1}$, więc $0\le S_N-S\le S_{2k+1}-S_{2k+2}=b_{N+1}$. Oba przypadki dają żądane oszacowanie.
\end{proof}
\begin{przyklad}
\label{ma1:ex:alternating-harmonic}
Szereg $\sum_{n=1}^\infty(-1)^{n-1}/n$ jest zbieżny, ponieważ $(1/n)$ jest nierosnący i zbiega do zera. Nie jest zbieżny bezwzględnie: szereg modułów jest szeregiem harmonicznym (tom I, rozdział 09). Dla jego sumy $S$ mamy
$S_4=7/12\le S\le S_5=47/60$.
Dla błędu mniejszego od $1/100$ wystarczy obliczyć $S_{100}$, ponieważ $|S-S_{100}|\le1/101<1/100$. Nie potrzebujemy tu identyfikacji $S$ z $\log 2$.
\end{przyklad}

\begin{figure}[htbp]
\centering
\begin{tikzpicture}[x=.75cm,y=4.2cm,>=Stealth]
\draw[->] (-.15,0)--(10.9,0) node[right] {$N$};
\draw[->] (0,-.04)--(0,1.13) node[above] {$S_N$};
\foreach \n in {1,...,10} {
 \draw (\n,.012)--(\n,-.012) node[below] {\small $\n$};
}
\draw[dashed,gray] (0,1)--(10.5,1);
\node[left] at (0,1) {$1$};
\draw[blue!70!black,thick] (0,0)--(2,.5)--(4,.583333)--(6,.616667)--(8,.634524)--(10,.645635);
\draw[orange!85!black,thick] (1,1)--(3,.833333)--(5,.783333)--(7,.759524)--(9,.745635);
\foreach \n/\s in {0/0,2/.5,4/.583333,6/.616667,8/.634524,10/.645635}
 \fill[blue!70!black] (\n,\s) circle (1.8pt);
\foreach \n/\s in {1/1,3/.833333,5/.783333,7/.759524,9/.745635}
 \fill[orange!85!black] (\n,\s) circle (1.8pt);
\node[blue!70!black,below] at (8,.60) {$S_{2k}$};
\node[orange!85!black,above] at (8,.81) {$S_{2k+1}$};
\draw[densely dotted,gray] (9,.745635)--(10.45,.745635);
\draw[densely dotted,gray] (10,.645635)--(10.45,.645635);
\draw[<->,gray] (10.35,.645635)--(10.35,.745635);
\node[right] at (10.45,.695635) {\small $1/10$};
\end{tikzpicture}
\caption{Parzyste i nieparzyste sumy częściowe szeregu naprzemiennego harmonicznego. Odcinki łączą punkty wyłącznie dla czytelności. Długość przedziału między $S_{10}$ a $S_9$ wynosi $1/10$.}
\label{ma1:fig:leibniz-sums}
\end{figure}
\FloatBarrier

\section{Część dodatnia i ujemna szeregu rzeczywistego}
\begin{definicja}
\label{ma2:def:positive-negative-parts}
Dla liczby rzeczywistej $a$ określamy $a^+=\max(a,0)$ i $a^-=\max(-a,0)$. Obie liczby są nieujemne. Dla $a\ge0$ mamy $a^+=a$, $a^-=0$, a dla $a<0$ mamy $a^+=0$, $a^-=-a$. W obu przypadkach zachodzą $a=a^+-a^-$ oraz $|a|=a^++a^-$.
\end{definicja}
\begin{twierdzenie}
\label{ma1:thm:positive-negative-parts}
Niech $(a_n)$ będzie ciągiem rzeczywistym. Szereg $\sum a_n$ jest zbieżny bezwzględnie wtedy i tylko wtedy, gdy oba szeregi $\sum a_n^+$ i $\sum a_n^-$ są zbieżne. Jeżeli $\sum a_n$ jest zbieżny warunkowo, to obie sumy częściowe tych szeregów dążą do $+\infty$.
\end{twierdzenie}
\begin{proof}
Jeżeli $\sum|a_n|$ jest zbieżny, z $0\le a_n^\pm\le |a_n|$ oraz kryterium porównawczego wynika zbieżność obu szeregów. Odwrotnie, jeśli oba są zbieżne, ich suma $\sum(a_n^++a_n^-)=\sum|a_n|$ jest zbieżna z liniowości.

Załóżmy teraz zbieżność warunkową. Niech $P_N=\sum_{n=1}^N a_n^+$, $Q_N=\sum_{n=1}^N a_n^-$ oraz $S_N=\sum_{n=1}^N a_n$. Mamy $S_N=P_N-Q_N$. Gdyby $(P_N)$ miało skończoną granicę, to $Q_N=P_N-S_N$ również miałoby skończoną granicę, wbrew brakowi zbieżności bezwzględnej. Gdyby natomiast $(Q_N)$ miało skończoną granicę, to $P_N=S_N+Q_N$ również miałoby skończoną granicę, z tym samym skutkiem. Oba ciągi są niemalejące, więc na mocy kryterium zbieżności szeregów nieujemnych z tomu I, rozdział 09, dążą do $+\infty$.
\end{proof}

\section{Twierdzenie Riemanna o przestawieniach}
\begin{twierdzenie}[Riemanna]
\label{ma1:thm:riemann-rearrangement}
Niech $(a_n)_{n\ge1}$ będzie ciągiem rzeczywistym takim, że szereg $\sum_{n=1}^\infty a_n$ jest zbieżny warunkowo (czyli zbieżny, lecz nie bezwzględnie). Dla każdej liczby $s\in\R$ istnieje bijekcja $\pi:\N_+\to\N_+$, dla której $\sum_{k=1}^\infty a_{\pi(k)}=s$. Istnieją również przestawienia, których sumy częściowe dążą odpowiednio do $+\infty$ i do $-\infty$, oraz przestawienie, którego sumy częściowe nie mają granicy nawet w $\overline{\R}$.
\end{twierdzenie}
\begin{proof}
Z twierdzenia \ref{ma1:thm:positive-negative-parts} wynika, że zbiór indeksów wyrazów dodatnich i zbiór indeksów wyrazów ujemnych są nieskończone: gdyby którykolwiek był skończony, odpowiedni ciąg sum częściowych części dodatniej lub ujemnej byłby od pewnego indeksu stały, wbrew rozbieżności do $+\infty$. Wyliczmy je rosnąco jako $(p_j)_{j\ge1}$ i $(q_j)_{j\ge1}$. Piszemy $u_j=a_{p_j}>0$ oraz $v_j=-a_{q_j}>0$. Każda suma częściowa części dodatniej szeregu jest sumą pierwszych kilku $u_j$ (lub sumą pustą), a dla każdego $J\ge1$ zachodzi $\sum_{j=1}^J u_j=\sum_{n=1}^{p_J}a_n^+$. Ponieważ $p_J\ge J$, prawa strona dąży do $+\infty$. Analogicznie $\sum_{j=1}^J v_j=\sum_{n=1}^{q_J}a_n^-\to+\infty$. Piszemy więc $\sum_{j=1}^\infty u_j=\sum_{j=1}^\infty v_j=+\infty$ w sensie rozbieżności do nieskończoności. Ponadto $a_n\to0$ na mocy warunku koniecznego zbieżności. Ponieważ $p_j,q_j\ge j$, mamy $a_{p_j}\to0$ i $a_{q_j}\to0$ jako podciągi ciągu zbieżnego do zera. Rachunek granic daje więc $u_j,v_j\to0$. Zera wyliczymy osobno rosnąco według indeksów; ich wstawianie nie zmienia wartości sum częściowych.

Ustalmy najpierw $s\in\R$. Zaczynamy od sumy $0$. Jeżeli $0>s$, dodajemy po kolei nieużyte wyrazy $-v_j$ do pierwszej chwili, w której suma stanie się mniejsza od $s$. Taki skończony krok istnieje, ponieważ $\sum_{j=1}^J v_j\to+\infty$ przy $J\to\infty$. W przeciwnym razie zaczynamy bez tego wstępnego kroku. Następnie dodajemy co najmniej jeden kolejny nieużyty wyraz $u_j$, kończąc blok w pierwszej chwili, gdy suma przekroczy $s$. Od tej chwili dodajemy kolejne nieużyte wyrazy $-v_j$, kończąc blok w pierwszej chwili, gdy suma spadnie poniżej $s$. Powtarzamy oba kroki w nieskończoność. Każdy blok jest skończony: sumy pozostałych dodatnich wyrazów i modułów pozostałych ujemnych wyrazów są nieograniczone po pominięciu skończonego początku. Każdy z dwóch rodzajów bloków zawiera co najmniej jeden wyraz.

Oznaczmy przez $U_k$ sumę częściową na końcu $k$-tego dodatniego bloku, a przez $V_k$ sumę na końcu następującego po nim ujemnego bloku. Niech $i_k,j_k$ będą numerami ostatnich wykorzystanych wyrazów odpowiednio w ciągach $(u_j)$ i $(v_j)$. Bezpośrednio przed ostatnim dodaniem wyrazu dodatniego suma jest nie większa od $s$, a bezpośrednio przed ostatnim dodaniem wyrazu ujemnego jest nie mniejsza od $s$. Wynika to z wyboru pierwszego przekroczenia progu (suma początkowa pierwszego dodatniego bloku jest również nie większa od $s$). Otrzymujemy
\[
 0<U_k-s\le u_{i_k},\qquad 0<s-V_k\le v_{j_k}.
\]
Ponieważ każdy blok zawiera co najmniej jeden nowy wyraz danego znaku, $i_k,j_k\ge k$. Dla dowolnego $\varepsilon>0$ wybieramy $K$ tak, aby $u_j,v_j<\varepsilon$ dla $j\ge K$. Dla $k\ge K$ oba przekroczenia progu są mniejsze od $\varepsilon$, zatem $U_k\to s$ i $V_k\to s$. Wszystkie sumy częściowe wewnątrz kolejnych bloków dodatnich leżą pomiędzy $V_{k-1}$ a $U_k$ (dla pierwszego bloku używamy sumy po kroku wstępnym). Wewnątrz bloków ujemnych leżą pomiędzy $V_k$ a $U_k$. Dla dowolnego $\varepsilon>0$ wszystkie te krańce, poza skończenie wieloma początkowymi blokami, należą do $(s-\varepsilon,s+\varepsilon)$. Każdy blok jest skończony, więc również cały ciąg sum częściowych zbiega do $s$.

Ponieważ w każdym cyklu wykorzystujemy co najmniej jeden nowy wyraz każdego znaku, a wyrazy danego znaku wybieramy według kolejności ich indeksów, każdy indeks wyrazu niezerowego zostanie użyty dokładnie raz. Po każdym pełnym cyklu wstawiamy jeden kolejny nieużyty indeks wyrazu zerowego, jeżeli taki jeszcze istnieje. Każdy indeks zera również zostanie użyty dokładnie raz. Powstałe wyliczenie wszystkich dodatnich liczb naturalnych bez powtórzeń określa bijekcję $\pi$; wstawienie zer nie zmienia wniosku o granicy.

Dla granicy $+\infty$ w $k$-tym cyklu dodajemy najpierw co najmniej jeden kolejny dodatni wyraz, a następnie tyle kolejnych dodatnich wyrazów, ile potrzeba, aby suma przekroczyła $k+v_k$. Dodajemy wtedy dokładnie jeden ujemny wyraz $-v_k$. Po tym kroku suma jest większa od $k$. W każdym cyklu suma najpierw rośnie, a potem spada do wartości większej od $k$. Dla $k\ge2$ wszystkie sumy wewnątrz cyklu są więc większe od $k-1$, bo początkowa suma jest większa od $k-1$. Po każdym cyklu można ponownie wstawić jedno zero. Wykorzystujemy wszystkie indeksy każdego znaku i każdy indeks zera, jak wyżej. Dla dowolnego $M\in\R$ wybieramy całkowite $K\ge2$ takie, że $K>M+1$; wszystkie sumy po pierwszych $K-1$ cyklach są większe od $M$. Otrzymujemy granicę $+\infty$. Przez zastosowanie konstrukcji do $(-a_n)$ uzyskujemy przestawienie o granicy $-\infty$.

Wreszcie możemy w $k$-tym cyklu dodawać kolejno wyrazy dodatnie do pierwszego przekroczenia poziomu $k$, a następnie wyrazy ujemne do pierwszego spadku poniżej $-k$. Oba bloki są skończone z rozbieżności odpowiednich pozostałych sum. Są niepuste, ponieważ pierwszy cykl rozpoczynamy od $0<1$, każdy następny od sumy mniejszej od $-(k-1)<k$, a każdy blok ujemny od sumy większej od $k>-k$. Dodajemy zera według tej samej reguły. Znowu otrzymujemy bijekcję, ale sumy częściowe na końcach dodatnich bloków dążą do $+\infty$, a na końcach ujemnych do $-\infty$. Granice dolna i górna tego ciągu sum wynoszą odpowiednio $-\infty$ i $+\infty$, ponieważ istnieją podciągi dążące do obu tych wartości. Zatem ciąg nie jest zbieżny w $\overline{\R}$.
\end{proof}

\begin{przyklad}
\label{ma1:ex:harmonic-rearrangement}
Niech $S=\sum_{n=1}^\infty(-1)^{n-1}/n$. Przestawmy wyrazy szeregu, umieszczając po dwa kolejne dodatnie wyrazy, a następnie jeden kolejny ujemny wyraz:
\[
 1+\frac13-\frac12+\frac15+\frac17-\frac14+\cdots.
\]
W $k$-tym bloku występują wyrazy o indeksach $4k-3,4k-1,2k$ szeregu wyjściowego. Każdy nieparzysty indeks ma dokładnie jedną z postaci $4k-3$ albo $4k-1$ dla $k\ge1$, a każdy parzysty indeks ma jednoznacznie postać $2k$ dla $k\ge1$. Zbiory tych trzech rodzajów indeksów są parami rozłączne, więc jest to przestawienie.
Oznaczmy sumy częściowe szeregu wyjściowego przez $S_N$, a przestawionego przez $T_N$. Dla $m\ge1$ zachodzi
\[
 T_{3m}-S_{4m}
 =\sum_{k=m+1}^{2m}\frac1{2k}
 =\frac12\sum_{k=m+1}^{2m}\frac1k.
\]
Istotnie, w obu sumach występuje dokładnie $2m$ pierwszych wyrazów dodatnich, natomiast $T_{3m}$ zawiera tylko $m$ pierwszych wyrazów ujemnych, a $S_{4m}$ zawiera ich $2m$. Różnica jest więc sumą modułów brakujących wyrazów ujemnych. Nie obliczamy tej granicy za pomocą logarytmu. Zauważmy natomiast, że
\[
 S_{2m}=\sum_{k=1}^{2m}\frac1k
 -2\sum_{k=1}^m\frac1{2k}
 =\sum_{k=m+1}^{2m}\frac1k.
\]
Stąd $T_{3m}=S_{4m}+\tfrac12 S_{2m}\to\tfrac32 S$. Ponieważ moduły wyrazów w kolejnych blokach dążą do zera, długość bloku jest stała i wynosi $3$, więc wniosek \ref{ma2:cor:bounded-blocks} daje $T_N\to\tfrac32 S$. Z kryterium Leibniza mamy $S\ge S_2=1/2$, zatem $\tfrac32 S\ne S$.
\end{przyklad}


\section{Kryterium sum po zbiorach skończonych}
Dla skończonego zbioru $F\subseteq\N_+$ oznaczmy
$S_F=\sum_{n\in F}x_n$, przy czym $S_\varnothing=0$.
Suma skończona nie zależy od kolejności składników.
Piszemy $S_F\to S$ względem inkluzji, jeśli dla każdego
$\varepsilon>0$ istnieje skończony $F_0$ taki, że
$\|S_F-S\|<\varepsilon$ dla wszystkich skończonych
$F\supseteq F_0$. Jest to zbieżność sieci sum skończonych.

\begin{twierdzenie}[Kryterium bezwarunkowej zbieżności]
\label{ma2:thm:unconditional-finite}
W przestrzeni Banacha następujące warunki są równoważne:
\begin{enumerate}[label=(\roman*)]
\item każdy przestawiony szereg jest zbieżny (bez
zakładania równości jego sum z sumą pierwotną);
\item szereg jest bezwarunkowo zbieżny;
\item istnieje $S\in E$, dla którego $S_F\to S$
względem inkluzji;
\item dla każdego $\varepsilon>0$ istnieje całkowite
$N\ge1$ takie, że dla każdego skończonego
$F\subseteq\{N+1,N+2,\ldots\}$ zachodzi
$\|S_F\|<\varepsilon$;
\item każdy podszereg (wybrany przez ściśle rosnący
ciąg dodatnich indeksów) jest zbieżny.
\end{enumerate}
Sumy z (ii) i (iii) są równe sumie szeregu wyjściowego.
\end{twierdzenie}
\begin{proof}
Implikacja (ii)$\Rightarrow$(i) wynika z definicji.

(i)$\Rightarrow$(iv): załóżmy, że (iv) jest fałszywe.
Istnieje $\varepsilon>0$ takie, że dla każdego $N\ge1$
można znaleźć skończony zbiór $F\subseteq\{N+1,\ldots\}$
z $\|S_F\|\ge\varepsilon$. Budujemy przestawienie etapami.
Przed etapem $k$ zbiór $D$ wszystkich wypisanych indeksów
jest skończony. Wybieramy $N=\max(D\cup\{k\})$,
a następnie odpowiedni niepusty zbiór $F_k$ poza tym
początkiem. Wypisujemy wszystkie jego elementy rosnąco,
po czym wypisujemy najmniejszy jeszcze niewypisany indeks.
Każdy etap jest skończony i żaden indeks nie powtarza się.
Po $k$ etapach wszystkie indeksy od $1$ do $k$ są wypisane:
indukcyjnie indeks $k$ albo już został wypisany, albo
jest najmniejszym brakującym indeksem i zostaje dopisany.
Powstaje więc bijekcja $\pi:\N_+\to\N_+$.
Sumy częściowe tuż przed i po bloku $F_k$ różnią się
o $S_{F_k}$, którego norma jest co najmniej $\varepsilon$.
Ponieważ oba odpowiadające im indeksy sum częściowych
dążą do nieskończoności, warunek Cauchy'ego nie zachodzi.
Przeczy to (i).

(iv)$\Rightarrow$(iii): dla $\varepsilon>0$ dobieramy $N$
z (iv). Jeśli $q\ge p\ge N+1$, to
$\|\sum_{n=p}^q x_n\|<\varepsilon$, ponieważ
$\{p,\ldots,q\}\subseteq\{N+1,N+2,\ldots\}$. Kryterium Cauchy'ego
(twierdzenie \ref{ma2:thm:series-cauchy}) i zupełność
dają istnienie $S=\lim_M S_{\{1,\ldots,M\}}$.
Teraz dobieramy $N$ tak, aby dla wszystkich skończonych
$F\subseteq\{N+1,\ldots\}$ było $\|S_F\|<\varepsilon/2$.
Dla każdego $M>N$ mamy
$\|S_{\{1,\ldots,M\}}-S_{\{1,\ldots,N\}}\|<\varepsilon/2$.
Przejście do granicy daje
$\|S-S_{\{1,\ldots,N\}}\|\le\varepsilon/2$.
Dla skończonego $G\supseteq\{1,\ldots,N\}$ otrzymujemy
\[
 \|S_G-S\|\le\|S_{G\setminus\{1,\ldots,N\}}\|
   +\|S_{\{1,\ldots,N\}}-S\|<\varepsilon.
\]
To jest (iii). Nie potrzebujemy ogólnego twierdzenia
o zupełności sieci, a jedynie zupełności ciągowej.

(iii)$\Rightarrow$(ii): niech $\pi$ będzie bijekcją.
Dla $\varepsilon>0$ wybieramy skończony $F_0$ z (iii).
Jeśli $F_0\ne\varnothing$, to dla
$k\ge\max_{n\in F_0}\pi^{-1}(n)$ zbiór
$\{\pi(1),\ldots,\pi(k)\}$ zawiera $F_0$; jeśli
$F_0=\varnothing$, inkluzja zachodzi dla każdego $k$.
Oznacza to, że odpowiednia suma częściowa jest odległa
od $S$ o mniej niż $\varepsilon$.

(iv)$\Rightarrow$(v): dla rosnących indeksów $n_j$
mamy $n_j\ge j$. Po wyborze $N$ z (iv), każda suma
$\sum_{j=p}^q x_{n_j}$ przy $q\ge p>N$ jest sumą po
skończonym zbiorze poza pierwszymi $N$ indeksami.
Zatem kryterium Cauchy'ego i zupełność dają zbieżność.

(v)$\Rightarrow$(iv): jeśli (iv) jest fałszywe, dla
pewnego $\varepsilon>0$ wybieramy rekurencyjnie niepuste
skończone zbiory $F_k$ z $\|S_{F_k}\|\ge\varepsilon$
tak, aby każdy indeks w $F_{k+1}$ był większy od
$\max F_k$. Wystarczy zastosować negację (iv) z
$N=\max F_k$ (w pierwszym kroku z $N=1$).
Wyliczenie elementów kolejnych bloków w porządku
rosnącym daje ściśle rosnący ciąg indeksów.
Jeśli $m_k=\sum_{j=1}^k |F_j|$ i $m_0=0$, to $m_k\ge k$
i różnica sum częściowych podszeregu o indeksach
$m_k$ i $m_{k-1}$ wynosi $S_{F_k}$. Dla dowolnego
proponowanego progu warunku Cauchy'ego wybieramy $k$
tak duże, by $m_{k-1}$ przekraczało próg. Norma różnicy
jest co najmniej $\varepsilon$, więc warunek Cauchy'ego
nie zachodzi. Sprzeczność dowodzi (iv).

W (iii) sumy po $\{1,\ldots,N\}$ zbiegają do $S$,
bo każdy skończony $F_0$ jest zawarty w takim zbiorze
dla dostatecznie dużego $N$. W (ii) wystarczy wziąć
bijekcję identycznościową.
\end{proof}

\begin{uwaga}
Zbieżność wszystkich przestawień implikuje ich wspólną
sumę również w niezupełnej przestrzeni unormowanej:
dowód (i)$\Rightarrow$(iv) nie korzysta z zupełności,
a w dowodzie (iv)$\Rightarrow$(iii) granica zwykłych
sum częściowych istnieje wtedy z (i). Sam warunek
(iv) bez założenia zupełności nie gwarantuje istnienia
sumy w $E$.
\end{uwaga}

\begin{twierdzenie}[Ograniczone mnożniki]
\label{ma2:thm:unconditional-multipliers}
Jeżeli $\sum x_n$ jest bezwarunkowo zbieżny w przestrzeni
Banacha $E$, a $(a_n)$ jest ograniczonym ciągiem skalarów,
to $\sum a_nx_n$ jest bezwarunkowo zbieżny.
\end{twierdzenie}
\begin{proof}
Najpierw dla skończonego zbioru $F$ i rzeczywistych
współczynników $0\le t_n\le1$ wykażemy oszacowanie
\[
 \left\|\sum_{n\in F}t_n x_n\right\|
 \le \max_{G\subseteq F}\|S_G\|.
\]
Dla $F=\varnothing$ obie strony są zerowe. Dla
$F=\{n_1,\ldots,n_m\}$ porządkujemy indeksy tak, by
$t_{n_1}\ge\cdots\ge t_{n_m}$. Dla wygody oznaczmy
$b_k=t_{n_k}$ dla $1\le k\le m$ i $b_{m+1}=0$.
Wówczas
\[
 \sum_{j=1}^m t_{n_j}x_{n_j}
 =\sum_{k=1}^m(b_k-b_{k+1})
               \sum_{j=1}^k x_{n_j}.
\]
Współczynnik przy $x_{n_j}$ po prawej to suma
różnic od $k=j$ do $m$, czyli $b_j=t_{n_j}$.
Wszystkie różnice są nieujemne, a ich suma wynosi
$t_{n_1}\le1$. Nierówność trójkąta daje oszacowanie.

Niech $M=\sup_n|a_n|$. Gdy $M=0$, wszystkie
wyrazy są zerowe. Dla $M>0$ każdy $a_n/M$
w przypadku rzeczywistym jest różnicą dwóch liczb
z $[0,1]$, gdyż $|a_n/M|\le1$. W przypadku zespolonym zapisujemy go jako
$(t_n^+-t_n^-)+\mathrm i(s_n^+-s_n^-)$,
gdzie $t_n^\pm,s_n^\pm\in[0,1]$ są częściami dodatnimi
i ujemnymi odpowiednich współrzędnych, których moduły
nie przekraczają $|a_n/M|\le1$.
Z powyższego oszacowania wynika w obu przypadkach
\[
 \left\|\sum_{n\in F}a_nx_n\right\|
 \le 4M\max_{G\subseteq F}\|S_G\|.
\]
Dla $\varepsilon>0$ stosujemy warunek (iv) twierdzenia
\ref{ma2:thm:unconditional-finite} z dokładnością
$\varepsilon/(4M)$. Jeśli $F$ leży poza wybranym
początkiem, każdy jego podzbiór również tam leży,
więc skończone maksimum jest mniejsze od
$\varepsilon/(4M)$. Otrzymujemy warunek (iv)
dla szeregu $\sum a_nx_n$, a zatem jego bezwarunkową
zbieżność.
\end{proof}

\section{Wymiar skończony i nieskończony}
\begin{twierdzenie}
\label{ma2:thm:finite-unconditional-absolute}
W skończenie wymiarowej przestrzeni unormowanej
zbieżność bezwarunkowa jest równoważna zbieżności
bezwzględnej.
\end{twierdzenie}
\begin{proof}
Przestrzeń $E$ jest zupełna z twierdzenia
\ref{ma2:thm:finite}, więc zbieżność bezwzględna
implikuje bezwarunkową na mocy twierdzenia
\ref{ma2:thm:absolute-rearrangement}.
Dla wymiaru $0$ teza zachodzi, ponieważ wszystkie
wyrazy są zerowe. Niech zatem $d=\dim E\ge1$
i niech $(e_j)_{j=1}^d$ będzie bazą.
Zapiszmy $x_n=\sum_{j=1}^d a_{nj}e_j$.
Jeśli $\sum x_n$ jest bezwarunkowo zbieżny, to
z twierdzenia \ref{ma2:thm:series-coordinates}
każdy szereg współrzędnych jest zbieżny przy
każdym przestawieniu.
Dla rzeczywistego szeregu skalarnego, który jest
zbieżny, lecz nie bezwzględnie, twierdzenie
\ref{ma1:thm:riemann-rearrangement} daje rozbieżne
przestawienie. Zatem w przypadku rzeczywistym
każdy szereg współrzędnych jest bezwzględnie zbieżny.
W przypadku zespolonym dla każdej współrzędnej
piszemy $a_{nj}=u_{nj}+\mathrm i v_{nj}$.
Ponieważ $|\operatorname{Re}z|,|\operatorname{Im}z|\le |z|$,
zbieżność dowolnego przestawienia szeregu zespolonego
implikuje zbieżność odpowiadających mu przestawień obu
szeregów rzeczywistych. Z poprzedniego argumentu
$\sum_n|u_{nj}|<\infty$ i $\sum_n|v_{nj}|<\infty$,
a z $|a_{nj}|\le|u_{nj}|+|v_{nj}|$ wynika
$\sum_n|a_{nj}|<\infty$. Dla każdego $N$ mamy
\[
 \sum_{n=1}^N\|x_n\|
 \le \sum_{j=1}^d\|e_j\|\sum_{n=1}^N|a_{nj}|
 \le \sum_{j=1}^d\|e_j\|\sum_{n=1}^\infty|a_{nj}|<\infty.
\]
W pierwszej nierówności zastosowaliśmy jednorodność
i nierówność trójkąta do każdego $x_n$, a następnie
zamieniliśmy kolejność dwóch sum skończonych.
Prawa strona jest niezależna od $N$, więc niemalejące
sumy częściowe szeregu norm są zbieżne.
\end{proof}
Twierdzenie Riemanna w pełnym zakresie dowolnej
zadanej sumy dotyczy szeregów rzeczywistych.
Nie wynika z niego możliwość otrzymania dowolnej
sumy zespolonej z warunkowo zbieżnego szeregu:
przykładowo wszystkie sumy częściowe przestawień
rzeczywistego szeregu leżą na osi rzeczywistej.

\begin{przyklad}[Bezwarunkowa, ale nie bezwzględna zbieżność]
\label{ma2:ex:unconditional-l2}
W przestrzeni Banacha $\ell^2$ niech $e_n$ oznacza
standardowy wektor jednostkowy i niech $x_n=e_n/n$.
Mamy $\|x_n\|_2=1/n$, więc szereg norm jest rozbieżny.
Natomiast dla skończonego $F\subseteq\{N+1,N+2,\ldots\}$
\[
 \left\|\sum_{n\in F}\frac{e_n}{n}\right\|_2^2
 =\sum_{n\in F}\frac1{n^2}\le\sum_{n>N}\frac1{n^2}.
\]
Ostatni ogon dąży do zera, ponieważ szereg
$\sum n^{-2}$ jest zbieżny (tom I, rozdział 09).
Dla $\varepsilon>0$ wybieramy $N$ z ogonem
mniejszym niż $\varepsilon^2$, co daje (iv)
z twierdzenia \ref{ma2:thm:unconditional-finite}.
Szereg jest więc bezwarunkowo zbieżny, lecz nie
bezwzględnie zbieżny.
Jego suma to $x=(1/j)_{j\ge1}\in\ell^2$:
dla zwykłych sum częściowych
$\|x-S_N\|_2^2=\sum_{j>N}1/j^2\to0$.
\end{przyklad}

\begin{uwaga}
Zbieżność bezwzględna implikuje bezwarunkową w przestrzeni
Banacha, a bezwarunkowa implikuje zwykłą. Pierwsza
implikacja nie odwraca się w ogólności, jak pokazuje
przykład \ref{ma2:ex:unconditional-l2}; druga nie odwraca
się nawet w $\R$, na mocy twierdzenia Riemanna i
przykładu \ref{ma1:ex:alternating-harmonic}.
Nie twierdzimy, że każde nieskończenie wymiarowe $E$
ma taki kontrprzykład, gdyż ogólny wynik
Dvoretzky'ego--Rogersa wymaga osobnego dowodu.
\end{uwaga}

% Materiał skalarny: AM I-11, zachowane etykiety wyników.
\chapter{Iloczyny Cauchy'ego}
\label{ma1:ch:cauchy-product}
\noindent Wykorzystujemy wyniki rozdziałów
\ref{ma2:ch:banach-series} i \ref{ma2:ch:unconditional-series}.

Mnożenie dwóch skończonych sum prowadzi do sumy iloczynów wszystkich par składników. Dla szeregów nieskończonych pojawia się problem: w jakiej kolejności należy dodawać te iloczyny i kiedy otrzymana suma jest iloczynem sum wyjściowych? Najpierw rozstrzygniemy to dla szeregów rzeczywistych i zespolonych, a następnie dla ciągłych odwzorowań dwuliniowych między przestrzeniami unormowanymi. Nie korzystamy z teorii funkcji wykładniczej ani szeregów potęgowych.

\section{Definicja i geometria indeksów}
\begin{definicja}
\label{ma1:def:cauchy-product}
Niech $K=\R$ lub $K=\C$, a $(a_n)_{n\ge0}$ i $(b_n)_{n\ge0}$ będą ciągami w $K$. Iloczynem Cauchy'ego szeregów $\sum_{n=0}^\infty a_n$ i $\sum_{n=0}^\infty b_n$ nazywamy szereg $\sum_{n=0}^\infty c_n$, gdzie dla każdego całkowitego $n\ge0$
\[
 c_n=\sum_{j=0}^n a_jb_{n-j}.
\]
Definicja nie wymaga zbieżności żadnego z trzech szeregów.
\end{definicja}
W całym rozdziale przez $\N_0$ oznaczamy zbiór $\{0,1,2,\ldots\}$. Indeksowanie od zera upraszcza wzory: wyraz $c_n$ zawiera dokładnie te iloczyny $a_jb_k$, dla których $j+k=n$. Dla szeregów indeksowanych od jedynki można najpierw przeindeksować wyrazy, przyjmując $\widetilde a_n=a_{n+1}$ i $\widetilde b_n=b_{n+1}$ dla $n\in\N_0$. Takie przeindeksowanie nie zmienia sum, ponieważ $\sum_{n=0}^N\widetilde a_n=\sum_{j=1}^{N+1}a_j$ i analogicznie dla drugiego szeregu.

\begin{lemat}[Tożsamości skończone]
\label{ma1:lem:cauchy-finite-identities}
Niech $(a_n),(b_n)$ będą ciągami w $K=\R$ lub $K=\C$, indeksowanymi od zera, i niech $(c_n)$ będzie określony jak w definicji \ref{ma1:def:cauchy-product}. Dla każdego całkowitego $N\ge0$ oznaczmy
$A_N=\sum_{j=0}^N a_j$, $B_N=\sum_{k=0}^N b_k$ i $C_N=\sum_{n=0}^N c_n$. Wówczas
\[
 C_N=\sum_{\substack{j,k\ge0\\j+k\le N}}a_jb_k
 =\sum_{j=0}^N a_jB_{N-j},
 \qquad
 A_NB_N=\sum_{j=0}^N\sum_{k=0}^N a_jb_k.
\]
We wszystkich sumach indeksy są całkowite.
\end{lemat}
\begin{proof}
Rozwijamy $C_N=\sum_{n=0}^N\sum_{j=0}^n a_jb_{n-j}$. Przypisanie $(n,j)\mapsto(j,n-j)$ jest bijekcją zbioru par całkowitych $0\le j\le n\le N$ na zbiór par nieujemnych $j,k$ z $j+k\le N$. Przypisaniem odwrotnym jest $(j,k)\mapsto(j+k,j)$. Dla ustalonego $j$ dopuszczalne wartości $k$ to $0,\ldots,N-j$, co daje drugą równość. Ostatnia równość wynika z rozdzielności mnożenia względem dodawania. Wszystkie przekształcenia dotyczą sum skończonych.
\end{proof}

\begin{figure}[htbp]
\centering
\begin{tikzpicture}[x=.8cm,y=.8cm,>=Stealth]
\fill[blue!10] (-.22,-.22) rectangle (4.22,4.22);
\fill[teal!22] (-.22,-.22)--(4.22,-.22)--(4.22,.22)--(.22,4.22)--(-.22,4.22)--cycle;
\draw[->] (-.4,0)--(5.3,0) node[right] {$j$};
\draw[->] (0,-.4)--(0,5.3) node[above] {$k$};
\foreach \j in {0,...,4}{
 \foreach \k in {0,...,4}{
  \pgfmathtruncatemacro{\d}{\j+\k}
  \ifnum\d>4
   \fill[blue!65!black] (\j,\k) circle (2pt);
  \else
   \fill[teal!75!black] (\j,\k) circle (2pt);
  \fi
 }
}
\draw[teal!75!black,thick] (0,4)--(4,0);
\draw[densely dashed,blue!65!black] (0,4)--(4,4)--(4,0);
\node[left] at (0,4) {$N$};
\node[below] at (4,0) {$N$};
\node[below left] at (0,0) {$0$};
\node[anchor=west,align=left] at (5.5,3.6) {\small Trójkąt: $j+k\le N$\\\small suma $C_N$};
\node[anchor=west,align=left] at (5.5,1.6) {\small Poza trójkątem: $j+k>N$\\\small część kwadratu $j,k\le N$};
\fill[teal!75!black] (5.25,3.65) circle (2pt);
\fill[blue!65!black] (5.25,1.65) circle (2pt);
\end{tikzpicture}
\caption{Dla $N=4$ suma częściowa iloczynu Cauchy'ego obejmuje zielone punkty trójkąta, a iloczyn sum częściowych wszystkie zielone i niebieskie punkty kwadratu. Odcinki i tło służą jedynie wyróżnieniu skończonych zbiorów indeksów.}
\label{ma1:fig:cauchy-index-sets}
\end{figure}
\FloatBarrier

\section{Iloczyn szeregów bezwzględnie zbieżnych}
\begin{twierdzenie}
\label{ma1:thm:absolute-cauchy-product}
Niech $(a_n),(b_n)$ będą ciągami w $\R$ lub $\C$, indeksowanymi od zera, i niech oba szeregi $\sum_{n=0}^\infty a_n$, $\sum_{n=0}^\infty b_n$ będą bezwzględnie zbieżne do odpowiednio $A,B$. Jeżeli $c_n=\sum_{j=0}^n a_jb_{n-j}$ dla $n\ge0$, to $\sum_{n=0}^\infty c_n$ jest bezwzględnie zbieżny, ma sumę $AB$ oraz
\[
 \sum_{n=0}^\infty |c_n|
 \le \left(\sum_{n=0}^\infty |a_n|\right)
      \left(\sum_{n=0}^\infty |b_n|\right).
\]
\end{twierdzenie}
\begin{proof}
Oznaczmy $\alpha=\sum_{n=0}^\infty|a_n|$, $\beta=\sum_{n=0}^\infty|b_n|$. Dla każdego całkowitego $N\ge0$ z nierówności trójkąta, lematu \ref{ma1:lem:cauchy-finite-identities} zastosowanego do modułów oraz nieujemności składników mamy
\[
 \sum_{n=0}^N|c_n|\le
 \sum_{\substack{j,k\ge0\\j+k\le N}}|a_j||b_k|
 \le \sum_{j=0}^N\sum_{k=0}^N |a_j||b_k|
 \le\alpha\beta.
\]
Sumy częściowe szeregu modułów są niemalejące i ograniczone z góry przez $\alpha\beta$. Z kryterium zbieżności szeregów nieujemnych (tom I, rozdział 09) szereg modułów jest zbieżny, a jego suma spełnia podane oszacowanie.

Niech $A_N,B_N,C_N$ oznaczają sumy częściowe z lematu \ref{ma1:lem:cauchy-finite-identities}. Z tożsamości skończonych wynika $A_NB_N-C_N=\sum_{0\le j,k\le N,\ j+k>N}a_jb_k$. Dla $m=\lfloor N/2\rfloor$, jeżeli $0\le j,k\le N$ i $j+k>N$, to $j>m$ lub $k>m$: w przeciwnym razie $j+k\le2m\le N$. Zatem
\[
 |A_NB_N-C_N|
 \le \beta\sum_{j>m}|a_j|+
      \alpha\sum_{k>m}|b_k|.
\]
Istotnie, sumę modułów po zbiorze brakujących par ograniczamy przez sumę po dwóch zbiorach: $m<j\le N$, $0\le k\le N$ oraz $0\le j\le N$, $m<k\le N$. Wspólne pary można policzyć dwukrotnie, gdyż składniki są nieujemne. Sumy po tych zbiorach rozkładają się na iloczyny i są nie większe od podanych oszacowań. Gdy $N\to\infty$, mamy $m\to\infty$, a obie reszty szeregów modułów dążą do zera, ponieważ $\sum_{j>m}|a_j|=\alpha-\sum_{j=0}^{m}|a_j|\to0$, a dla drugiego szeregu zachodzi analogiczna równość. Oszacowanie daje $|A_NB_N-C_N|\to0$. Ponieważ $A_NB_N\to AB$, z nierówności $|C_N-AB|\le |C_N-A_NB_N|+|A_NB_N-AB|$ otrzymujemy $C_N\to AB$.
\end{proof}

\begin{przyklad}
\label{ma1:ex:geometric-product}
Niech $q\in\C$ i $|q|<1$. Biorąc $a_n=b_n=q^n$ dla $n\ge0$, otrzymujemy $c_n=\sum_{j=0}^n q^jq^{n-j}=(n+1)q^n$ (dla $q=0$ stosujemy konwencję $q^0=1$). Oba szeregi geometryczne są bezwzględnie zbieżne, ponieważ $|q^n|=|q|^n$ i $|q|<1$. Z twierdzenia \ref{ma2:thm:vector-geometric}, zastosowanego do $v=1$, każdy ma sumę $1/(1-q)$, więc
\[
 \sum_{n=0}^\infty (n+1)q^n=\frac1{(1-q)^2}.
\]
Dla $q=1/2$ suma wynosi $4$. Wzór wynika z mnożenia szeregów, bez różniczkowania i bez funkcji potęgowych o wykładnikach rzeczywistych.
\end{przyklad}

\section{Twierdzenie Mertensa}
\begin{twierdzenie}[Mertensa]
\label{ma1:thm:mertens-product}
Niech $(a_n),(b_n)$ będą ciągami w $\R$ lub $\C$, indeksowanymi od zera. Załóżmy, że $\sum_{n=0}^\infty a_n=A$, $\sum_{n=0}^\infty b_n=B$ i co najmniej jeden z tych szeregów jest zbieżny bezwzględnie. Wtedy iloczyn Cauchy'ego, czyli szereg o wyrazach $c_n=\sum_{j=0}^n a_jb_{n-j}$, jest zbieżny do $AB$.
\end{twierdzenie}
\begin{proof}
Możemy założyć, że szereg $\sum a_n$ jest bezwzględnie zbieżny. Istotnie, dla każdego $n\in\N_0$ podstawienie $k=n-j$ w skończonej sumie oraz przemienność mnożenia dają $\sum_{j=0}^n a_jb_{n-j}=\sum_{k=0}^n b_ka_{n-k}$, więc zamiana ról obu ciągów nie zmienia $c_n$. Oznaczmy $\alpha=\sum_{n=0}^\infty|a_n|$ oraz dla każdego całkowitego $N\ge0$ przyjmijmy $A_N=\sum_{j=0}^N a_j$, $B_N=\sum_{k=0}^N b_k$, $C_N=\sum_{n=0}^N c_n$. Z $B_n\to B$ wynika ograniczoność ciągu $(B_n-B)$, więc istnieje $M\ge1$ takie, że $|B_n-B|\le M$ dla każdego $n\ge0$.

Z lematu \ref{ma1:lem:cauchy-finite-identities} otrzymujemy
\[
 C_N-A_N B=\sum_{j=0}^N a_j(B_{N-j}-B).
\]
Ustalmy $\varepsilon>0$. Wybierzmy $J\in\N_0$ z $M\sum_{j>J}|a_j|<\varepsilon/2$, korzystając ze zbieżności szeregu modułów i ustalonego skończonego $M$. Ponieważ $B_n\to B$, istnieje $L\ge0$ takie, że $|B_n-B|<\varepsilon/(2(1+\alpha))$ dla $n\ge L$. Dla $N\ge J+L$ oraz $0\le j\le J$ mamy $N-j\ge L$, zatem
\[
\begin{aligned}
 |C_N-A_N B|
 &\le\sum_{j=0}^J |a_j||B_{N-j}-B|
      +\sum_{j=J+1}^N |a_j||B_{N-j}-B|\\
 &\le\frac{\varepsilon\alpha}{2(1+\alpha)}
        +M\sum_{j>J}|a_j|
 <\varepsilon.
\end{aligned}
\]
Druga suma jest zerowa, jeśli $N=J$. Wobec tego $C_N-A_NB\to0$. Ponieważ $A_N\to A$, otrzymujemy $C_N\to AB$.
\end{proof}

\begin{przyklad}
Niech $a_n=2^{-n}$ oraz $b_n=(-1)^n/(n+1)$ dla $n\ge0$. Pierwszy szereg ma sumę $2$ i jest zbieżny bezwzględnie, drugi jest zbieżny warunkowo z kryterium Leibniza. Jeżeli $S=\sum_{n=0}^\infty(-1)^n/(n+1)$, twierdzenie Mertensa daje
\[
 \sum_{n=0}^\infty\left(\sum_{j=0}^n\frac{(-1)^{n-j}}{2^j(n-j+1)}\right)=2S.
\]
Iloczyn ten jest zbieżny warunkowo. Istotnie, z definicji $c_0=b_0$ i dla $n\ge1$ zachodzi $c_n=b_n+\frac12c_{n-1}$ (wydzielamy składnik $j=0$ i w pozostałej sumie zmieniamy indeks na $j-1$). Gdyby $\sum_{n=0}^\infty|c_n|$ był zbieżny, to szereg $\sum_{n=1}^\infty(|c_n|+\frac12|c_{n-1}|)$ również byłby zbieżny na mocy liniowości i niezmienniczości zbieżności po pominięciu skończonego początku. Z $|b_n|\le |c_n|+\frac12|c_{n-1}|$ dla $n\ge1$ i kryterium porównawczego wynikałaby zbieżność $\sum_{n=0}^\infty|b_n|$, sprzecznie z rozbieżnością szeregu harmonicznego.
\end{przyklad}

\section{Iloczyn dwóch szeregów zbieżnych warunkowo}
\begin{przyklad}
\label{ma1:ex:divergent-cauchy-product}
Niech $a_n=b_n=(-1)^n/\sqrt{n+1}$ dla całkowitych $n\ge0$. Oba szeregi są zbieżne z kryterium Leibniza: $1/\sqrt{n+1}$ jest ciągiem malejącym, a dla $\varepsilon>0$ nierówność $n+1>1/\varepsilon^2$ daje $1/\sqrt{n+1}<\varepsilon$. Nie są zbieżne bezwzględnie, ponieważ $1/\sqrt{n+1}\ge1/(n+1)$, a szereg harmoniczny jest rozbieżny.
Jednak
\[
 c_n=(-1)^n\sum_{j=0}^n
       \frac1{\sqrt{(j+1)(n-j+1)}}.
\]
Dla $0\le j\le n$ mamy $j+1\le n+1$ i $n-j+1\le n+1$, zatem $\sqrt{(j+1)(n-j+1)}\le n+1$. Wszystkie składniki sumy bez znaku są dodatnie, a jest ich $n+1$. Multiplikatywność nieujemnych pierwiastków kwadratowych wynika z jednoznaczności pierwiastka: iloczyn dwóch nieujemnych pierwiastków jest nieujemny, a jego kwadrat jest iloczynem liczb podpierwiastkowych. Wobec tego
\[
 |c_n|\ge(n+1)\frac1{n+1}=1.
\]
Wyrazy iloczynu Cauchy'ego nie dążą do zera. Iloczyn ten jest rozbieżny na mocy twierdzenia \ref{ma2:thm:series-cauchy}.
\end{przyklad}
Pytanie o istnienie sumy iloczynu Cauchy'ego trzeba odróżnić od pytania o jej wartość. Następne twierdzenie pokazuje, że jeśli wszystkie trzy szeregi są zbieżne, to suma iloczynu musi być iloczynem sum, nawet bez zbieżności bezwzględnej.

\section{Twierdzenie Abela o iloczynie Cauchy'ego}
\begin{lemat}
\label{ma1:lem:triangular-means}
Niech $K=\R$ lub $K=\C$, a $(x_n)_{n\ge0},(y_n)_{n\ge0}$ będą ciągami w $K$ zbieżnymi odpowiednio do $x,y\in K$. Wówczas
\[
 \lim_{N\to\infty}\frac1{N+1}\sum_{j=0}^N x_jy_{N-j}=xy.
\]
\end{lemat}
\begin{proof}
Ciąg $(y_n)$ jest ograniczony, więc wybieramy $M\ge1$ z $|y_n|\le M$ dla każdego $n$ oraz $|x|\le M$. Dla dowolnych całkowitych $j,k\ge0$ rozkład $x_jy_k-xy=(x_j-x)y_k+x(y_k-y)$ i nierówność trójkąta dają
$|x_jy_k-xy|\le M(|x_j-x|+|y_k-y|)$.
Ustalmy $\varepsilon>0$. Istnieje całkowite $J\ge1$ takie, że $|x_j-x|<\varepsilon/(4M)$ i $|y_j-y|<\varepsilon/(4M)$ dla każdego $j\ge J$. Jeżeli $N\ge2J$, to dla $J\le j\le N-J$ zachodzi $|x_jy_{N-j}-xy|<\varepsilon/2$.

Z ograniczoności obu ciągów istnieje rzeczywiste $D>0$ takie, że $|x_jy_k-xy|\le D$ dla wszystkich $j,k\ge0$. Na przykład dla ograniczeń $|x_j|\le P$, $|y_k|\le Q$ można wziąć $D=PQ+|xy|+1$. Dla $N\ge2J$ w sumie po $0\le j\le N$ dokładnie $2J$ indeksów nie spełnia $J\le j\le N-J$: po $J$ indeksów przy każdym końcu. W konsekwencji
\[
 \left|\frac1{N+1}\sum_{j=0}^N x_jy_{N-j}-xy\right|
 \le \frac{\varepsilon}{2}+\frac{2JD}{N+1}.
\]
Dla każdego całkowitego $N\ge2J$ spełniającego ponadto $N+1>4JD/\varepsilon$ drugi składnik jest mniejszy od $\varepsilon/2$. Taki próg istnieje z własności Archimedesa, co dowodzi tezy.
\end{proof}
\begin{twierdzenie}[Abela o iloczynie Cauchy'ego]
\label{ma1:thm:abel-cauchy-product}
Niech $(a_n),(b_n)$ będą ciągami liczb rzeczywistych lub zespolonych, indeksowanymi od zera. Jeżeli szeregi $\sum_{n=0}^\infty a_n$, $\sum_{n=0}^\infty b_n$ oraz ich iloczyn Cauchy'ego $\sum_{n=0}^\infty c_n$, gdzie $c_n=\sum_{j=0}^n a_jb_{n-j}$, są zbieżne odpowiednio do $A,B,C$, to $C=AB$.
\end{twierdzenie}
\begin{proof}
Oznaczmy sumy częściowe przez $A_N,B_N,C_N$ jak w lemacie \ref{ma1:lem:cauchy-finite-identities}. Dla każdego całkowitego $N\ge0$ zachodzi tożsamość
\[
 \sum_{m=0}^N C_m=\sum_{j=0}^N A_j B_{N-j}.
\]
Aby ją uzasadnić, ustalmy parę nieujemnych indeksów całkowitych $r,s$ z $r+s\le N$. Składnik $a_rb_s$ po lewej występuje w sumach $C_m$ dokładnie dla $r+s\le m\le N$, a więc $N-r-s+1$ razy. Po prawej występuje dokładnie dla $r\le j\le N-s$, również $N-r-s+1$ razy. Jeżeli $r+s>N$, nie występuje po żadnej ze stron. Ponieważ rozpatrujemy skończone sumy, równość liczności wystąpień dla każdej pary dowodzi tożsamości.

Dzielimy obie strony przez $N+1$. Z lematu \ref{ma1:lem:triangular-means} zastosowanego do $x_j=A_j$, $y_j=B_j$ wynika, że prawa strona po podzieleniu dąży do $AB$. Stosując ten sam lemat do $x_j=C_j$ oraz stałego ciągu $y_j=1$, otrzymujemy, że lewa strona po podzieleniu dąży do $C$. Z jednoznaczności granicy mamy $C=AB$.
\end{proof}
Zaletą tego dowodu jest niezależność od teorii szeregów potęgowych.

\begin{przyklad}
\label{ma1:ex:convergent-conditional-product}
Niech $a_n=b_n=(-1)^n/(n+1)$ dla $n\in\N_0$. Oba szeregi są warunkowo zbieżne do tej samej sumy $S$ (przykład \ref{ma1:ex:alternating-harmonic}). Tym razem ich iloczyn Cauchy'ego jest zbieżny. Dla każdego $n\ge0$ mamy
\[
 c_n=(-1)^n\sum_{j=0}^n\frac1{(j+1)(n-j+1)}
     =(-1)^n\frac{2H_{n+1}}{n+2},
 \qquad H_m=\sum_{j=1}^m\frac1j\quad(m\ge1).
\]
Druga równość wynika z tożsamości
$\frac1{(j+1)(n-j+1)}=\frac1{n+2}(\frac1{j+1}+\frac1{n-j+1})$ i zmiany indeksu w drugiej sumie.
Ciąg $H_m/(m+1)$ jest nierosnący, ponieważ
\[
 \frac{H_m}{m+1}-\frac{H_{m+1}}{m+2}
 =\frac{H_m-1}{(m+1)(m+2)}\ge0.
\]
Ponadto dąży do zera: dla $\varepsilon>0$ wybieramy $J\ge1$ z $1/(J+1)<\varepsilon/2$. Dla $m\ge J$ mamy $H_m\le H_J+(m-J)/(J+1)$, stąd $0\le H_m/(m+1)\le H_J/(m+1)+1/(J+1)<\varepsilon$ dla dostatecznie dużego $m$.
Z kryterium Leibniza szereg $\sum_{n=0}^\infty c_n$ jest zbieżny, a z twierdzenia \ref{ma1:thm:abel-cauchy-product} jego suma wynosi $S^2$. Nie jest bezwzględnie zbieżny, gdyż $2H_{n+1}/(n+2)\ge2/(n+2)$, a szereg po prawej jest rozbieżny przez porównanie z szeregiem harmonicznym po pominięciu pierwszego wyrazu. Wraz z przykładem \ref{ma1:ex:divergent-cauchy-product} pokazuje to, że sama warunkowa zbieżność obu czynników nie rozstrzyga zbieżności iloczynu.
\end{przyklad}
 Wykorzystana własność średnich arytmetycznych pojawi się ponownie przy omawianiu metod sumowania szeregów.

\section{Sumowanie iloczynów po wszystkich parach indeksów}
\begin{twierdzenie}
\label{ma1:thm:pair-enumeration}
Niech $(a_n)_{n\ge0},(b_n)_{n\ge0}$ będą ciągami w $\R$ lub $\C$ takimi, że oba szeregi są zbieżne bezwzględnie, odpowiednio do $A,B$. Dla dowolnej bijekcji $\pi:\N_+\to\N_0\times\N_0$, gdzie $\pi(l)=(j_l,k_l)$, szereg
\[
 \sum_{l=1}^\infty a_{j_l}b_{k_l}
\]
jest zbieżny bezwzględnie i ma sumę $AB$.
\end{twierdzenie}
\begin{proof}
Istnieje wyliczenie wszystkich par w kolejności rosnącej sumy współrzędnych: dla kolejnych $n=0,1,\ldots$ wypisujemy $(0,n),(1,n-1),\ldots,(n,0)$. Każda para $(j,k)$ należy do dokładnie jednego bloku, o numerze $n=j+k$, i występuje w nim dokładnie raz. Blok o numerze $n$ ma $n+1$ elementów, a przed nim jest $\sum_{i=0}^{n-1}(i+1)=n(n+1)/2$ elementów (suma pusta dla $n=0$). Wobec tego pozycja pary $(j,k)$, dla $n=j+k$, jest równa $n(n+1)/2+j+1$. Bloki wypełniają kolejne rozłączne przedziały dodatnich indeksów, a ich długości mają nieograniczone sumy częściowe. Otrzymujemy bijekcję $\pi_0:\N_+\to\N_0\times\N_0$.
Oznaczmy $\alpha=\sum_{n\ge0}|a_n|$ i $\beta=\sum_{n\ge0}|b_n|$. Dla dowolnego skończonego początkowego fragmentu tego wyliczenia istnieje całkowite $M\ge0$ takie, że wszystkie jego pary mają obie współrzędne nie większe od $M$. Suma modułów odpowiadających tym parom jest nie większa od
$(\sum_{j=0}^M|a_j|)(\sum_{k=0}^M|b_k|)\le\alpha\beta$. Zatem szereg wyrazów $a_jb_k$ wyliczonych przez $\pi_0$ jest zbieżny bezwzględnie, a więc zbieżny.
Grupując jego kolejne bloki według $j+k=n$, otrzymujemy iloczyn Cauchy'ego. Twierdzenie \ref{ma2:thm:series-grouping} zapewnia zachowanie sumy, a twierdzenie \ref{ma1:thm:absolute-cauchy-product} identyfikuje ją z $AB$. Dowolna bijekcja $\pi$ wyznacza przestawienie wyliczenia $\pi_0$, mianowicie bijekcję $\pi_0^{-1}\circ\pi$ zbioru $\N_+$ na siebie. Z twierdzenia \ref{ma2:thm:absolute-rearrangement} każde takie przestawienie jest zbieżne bezwzględnie do tej samej sumy.
\end{proof}

\section{Wielokrotny iloczyn Cauchy'ego}
\begin{wniosek}
\label{ma1:cor:finite-cauchy-products}
Niech $r\in\N_+$ i niech dla każdego $i\in\{1,\ldots,r\}$ szereg $\sum_{n=0}^\infty a_n^{(i)}$ w $\R$ lub $\C$ będzie zbieżny bezwzględnie do $A_i$. Dla każdego $n\in\N_0$ połóżmy
\[
 d_n=\sum_{\substack{n_1,\ldots,n_r\in\N_0\\
                         n_1+\cdots+n_r=n}}
                    \prod_{i=1}^r a_{n_i}^{(i)}.
\]
Każda suma definiująca $d_n$ jest skończona, ponieważ wszystkie $n_i$ należą do $\{0,\ldots,n\}$. Wówczas szereg $\sum_{n=0}^\infty d_n$ jest zbieżny bezwzględnie i ma sumę $\prod_{i=1}^r A_i$. Ponadto
\[
 \sum_{n=0}^\infty |d_n|
 \le \prod_{i=1}^r\left(\sum_{n=0}^\infty |a_n^{(i)}|\right).
\]
\end{wniosek}
\begin{proof}
Dla $r=1$ teza jest bezpośrednia. Załóżmy ją dla $r$ szeregów i rozważmy $r+1$ szeregów. Oznaczmy współczynniki iloczynu pierwszych $r$ przez $d_j$. Współczynnik iloczynu Cauchy'ego tego szeregu i ostatniego szeregu jest równy
\[
 \sum_{j=0}^n d_j a_{n-j}^{(r+1)}
 =\sum_{\substack{n_1,\ldots,n_{r+1}\in\N_0\\
                        n_1+\cdots+n_{r+1}=n}}
             \prod_{i=1}^{r+1} a_{n_i}^{(i)}.
\]
Istotnie, każda krotka po prawej występuje dokładnie raz po lewej, dla $j=n_1+\cdots+n_r$. Wszystkie te sumy są skończone, bo każdy indeks należy do $\{0,\ldots,n\}$. Z założenia indukcyjnego szereg $\sum d_j$ jest bezwzględnie zbieżny. Stosując twierdzenie \ref{ma1:thm:absolute-cauchy-product}, otrzymujemy bezwzględną zbieżność nowego szeregu oraz iloczyn sum i oszacowanie sumy modułów przez iloczyn poprzedniego oszacowania i $\sum_{n\ge0}|a_n^{(r+1)}|$. Indukcja kończy dowód.
\end{proof}
\begin{uwaga}
Iloczyn Cauchy'ego nie oznacza mnożenia wyraz po wyrazie, czyli szeregu $\sum a_nb_n$. Pozwala on natomiast przejść od iloczynu dwóch szeregów do sumy współczynników wyznaczanych przez skończone sumy. Będzie to podstawą dowodu prawa dodawania argumentów funkcji wykładniczej. Poniżej rozszerzymy te wyniki na odwzorowania dwuliniowe.
\end{uwaga}

\section{Iloczyn Cauchy'ego względem odwzorowania dwuliniowego}
Niech $E,F,G$ będą przestrzeniami unormowanymi nad tym samym
ciałem $\K=\R$ lub $\K=\C$. Nie zakładamy istnienia
mnożenia wektorów w żadnej z tych przestrzeni.

\begin{lemat}[Ciągłość odwzorowania dwuliniowego]
\label{ma2:lem:bilinear-bound}
Dla odwzorowania dwuliniowego $T:E\times F\to G$
następujące warunki są równoważne: ciągłość $T$ względem
normy $\|(x,y)\|=\max(\|x\|,\|y\|)$ oraz istnienie
stałej $K\ge0$ takiej, że dla wszystkich $x\in E$, $y\in F$
\[
 \|T(x,y)\|\le K\|x\|\|y\|.
\]
\end{lemat}
\begin{proof}
Jeżeli $T$ jest ciągłe, to z $T(0,0)=0$ istnieje
$\delta>0$ takie, że $\|T(u,v)\|<1$ dla
$\max(\|u\|,\|v\|)<\delta$. Dla $x\ne0$, $y\ne0$
podstawiamy $u=\delta x/(2\|x\|)$ i $v=\delta y/(2\|y\|)$.
Obie normy są równe $\delta/2$, więc dwuliniowość daje
$\|T(x,y)\|<4\|x\|\|y\|/\delta^2$.
Jeśli którykolwiek wektor jest zerowy, $T(x,y)=0$.
Można zatem przyjąć $K=4/\delta^2$.
Odwrotnie, z dwuliniowości
\[
 T(x',y')-T(x,y)=T(x'-x,y')+T(x,y'-y).
\]
Jeśli $(x',y')\to(x,y)$, to $\|y'\|\le\|y\|+1$ od
pewnego miejsca. Norma różnicy nie przekracza
$K(\|x'-x\|(\|y\|+1)+\|x\|\|y'-y\|)$, co dąży do zera.
Kryterium ciągowe ciągłości w przestrzeniach metrycznych
daje ciągłość w każdym punkcie.
\end{proof}

\begin{definicja}
\label{ma2:def:bilinear-cauchy}
Dla ciągłego odwzorowania dwuliniowego $T:E\times F\to G$
oraz ciągów $(x_n)_{n\ge0}$ w $E$ i $(y_n)_{n\ge0}$ w $F$
określamy iloczyn Cauchy'ego względem $T$ przez
$z_n=\sum_{j=0}^n T(x_j,y_{n-j})$. Jest to ciąg w $G$.
Oznaczamy
\[
 X_N=\sum_{j=0}^N x_j,\qquad Y_N=\sum_{k=0}^N y_k,\qquad
 Z_N=\sum_{n=0}^N z_n.
\]
\end{definicja}
Dwuliniowość oraz bijekcja indeksów z lematu
\ref{ma1:lem:cauchy-finite-identities} dają
\[
\begin{aligned}
 Z_N&=\sum_{j+k\le N}T(x_j,y_k)
 =\sum_{j=0}^N T(x_j,Y_{N-j})
 =\sum_{k=0}^N T(X_{N-k},y_k),\\
 T(X_N,Y_N)&=\sum_{j=0}^N\sum_{k=0}^N T(x_j,y_k).
\end{aligned}
\]
W pierwszej sumie $j,k\in\N_0$; wszystkie sumy są skończone.
Druga postać $Z_N$ wynika z zebrania składników dla
ustalonego $j$ i liniowości w drugim argumencie $T$;
trzecia z analogicznego zebrania składników dla ustalonego
$k$ i liniowości w pierwszym argumencie. Rozwinięcie
$T(X_N,Y_N)$ wykorzystuje kolejno liniowość w obu argumentach.
Nie zamieniamy argumentów $T$, gdyż przestrzenie $E,F$
mogą być różne i przemienność nie jest założona.

\begin{twierdzenie}[Mertensa w wersji dwuliniowej]
\label{ma2:thm:bilinear-mertens}
Niech $T:E\times F\to G$ będzie ciągłe dwuliniowe.
Załóżmy, że $\sum x_n=x\in E$ i $\sum y_n=y\in F$.
Jeżeli co najmniej jeden z szeregów jest bezwzględnie
zbieżny, to $\sum z_n=T(x,y)$ w $G$.
Nie wymaga to zupełności żadnej z przestrzeni.
\end{twierdzenie}
\begin{proof}
Wybieramy $K\ge1$ spełniające oszacowanie z lematu
\ref{ma2:lem:bilinear-bound}. Najpierw załóżmy
$\alpha=\sum_{j\ge0}\|x_j\|<\infty$.
Z $Y_m\to y$ wynika istnienie $M\ge1$ z
$\|Y_m-y\|\le M$ dla każdego $m\ge0$.
Z tożsamości skończonych otrzymujemy
$Z_N-T(X_N,y)=\sum_{j=0}^N T(x_j,Y_{N-j}-y)$.
Dla $\varepsilon>0$ wybieramy $J\ge0$ tak, by
$KM\sum_{j>J}\|x_j\|<\varepsilon/2$.
Następnie wybieramy $L\ge0$ tak, by
$\|Y_m-y\|<\varepsilon/(2K(1+\alpha))$ dla $m\ge L$.
Dla $N\ge J+L$ mamy $N-j\ge L$ przy $j\le J$, więc
\[
\begin{aligned}
 \|Z_N-T(X_N,y)\|
 &\le K\sum_{j=0}^J\|x_j\|\|Y_{N-j}-y\|
    +K\sum_{j=J+1}^N\|x_j\|\|Y_{N-j}-y\|\\
 &\le \frac{\varepsilon\alpha}{2(1+\alpha)}
      +KM\sum_{j>J}\|x_j\|<\varepsilon.
\end{aligned}
\]
Druga suma może być pusta. Ponadto
$\|T(X_N,y)-T(x,y)\|\le K\|X_N-x\|\|y\|\to0$.
Nierówność trójkąta daje $Z_N\to T(x,y)$.

Jeżeli zamiast tego $\sum\|y_k\|<\infty$, stosujemy
już udowodniony przypadek do odwzorowania
$\widetilde T:F\times E\to G$ określonego przez
$\widetilde T(v,u)=T(u,v)$. Jest ono dwuliniowe
i spełnia to samo oszacowanie z $K$, więc jest ciągłe.
Jego współczynnik iloczynu szeregu $(y_n)$ i $(x_n)$
wynosi $\sum_{k=0}^n T(x_{n-k},y_k)=z_n$ po zmianie
indeksu $j=n-k$. Otrzymana suma jest równa
$\widetilde T(y,x)=T(x,y)$, co kończy dowód.
\end{proof}

\begin{wniosek}[Dwa szeregi bezwzględnie zbieżne]
\label{ma2:cor:bilinear-absolute}
Niech $T:E\times F\to G$ będzie ciągłe dwuliniowe,
a $K\ge0$ niech spełnia $\|T(u,v)\|\le K\|u\|\|v\|$
dla wszystkich $u\in E$, $v\in F$. Jeżeli oba szeregi
mają sumy $x,y$ i oba szeregi norm są zbieżne, to iloczyn Cauchy'ego
jest bezwzględnie zbieżny do $T(x,y)$ i
\[
 \sum_{n=0}^{\infty}\|z_n\|
 \le K\left(\sum_{j=0}^\infty\|x_j\|\right)
      \left(\sum_{k=0}^\infty\|y_k\|\right).
\]
W szczególności gdy $E,F$ są przestrzeniami Banacha,
wystarczy założyć zbieżność obu szeregów norm.
Zupełność $G$ nie jest potrzebna, ponieważ granica
iloczynu jest konkretnym elementem $T(x,y)\in G$.
\end{wniosek}
\begin{proof}
Z oszacowania dwuliniowego i nierówności trójkąta
dla każdego $N\ge0$ mamy
\[
 \sum_{n=0}^N\|z_n\|
 \le K\sum_{j+k\le N}\|x_j\|\|y_k\|
 \le K\left(\sum_{j=0}^N\|x_j\|\right)
      \left(\sum_{k=0}^N\|y_k\|\right).
\]
Ostatni iloczyn jest ograniczony przez iloczyn pełnych
sum norm. Sumy częściowe szeregu norm są niemalejące,
więc mają skończoną granicę spełniającą oszacowanie.
Twierdzenie \ref{ma2:thm:bilinear-mertens} daje
zbieżność szeregu wektorowego do $T(x,y)$.
Gdy $E,F$ są zupełne, istnienie $x,y$ wynika
z twierdzenia \ref{ma2:thm:absolute-series} po przeindeksowaniu.
\end{proof}

\begin{lemat}[Średnie po przekątnych]
\label{ma2:lem:bilinear-diagonal-means}
Jeśli $u_j\to u$ w $E$, $v_j\to v$ w $F$ i $T$ jest
ciągłe dwuliniowe, to
\[
 \frac1{N+1}\sum_{j=0}^N T(u_j,v_{N-j})\longrightarrow T(u,v).
\]
\end{lemat}
\begin{proof}
Wybierzmy $K\ge1$ z lematu \ref{ma2:lem:bilinear-bound}
oraz $M\ge1$ z $\|v_k\|,\|u\|\le M$ dla wszystkich $k$.
Dla wszystkich $j,k$ dwuliniowość daje
\[
 \|T(u_j,v_k)-T(u,v)\|
 \le KM(\|u_j-u\|+\|v_k-v\|).
\]
Dla $\varepsilon>0$ wybieramy $J\ge1$ z
$\|u_j-u\|,\|v_j-v\|<\varepsilon/(4KM)$ dla $j\ge J$.
Jeżeli $N\ge2J$ i $J\le j\le N-J$, norma różnicy
jest mniejsza od $\varepsilon/2$.
Z ograniczoności obu ciągów istnieją $P,Q\ge1$
z $\|u_j\|\le P$, $\|v_k\|\le Q$. Dla wszystkich par
norma różnicy nie przekracza $D=KPQ+\|T(u,v)\|$.
Poza środkowym przedziałem jest dokładnie $2J$ indeksów,
więc norma różnicy średniej od $T(u,v)$ jest nie większa
niż $\varepsilon/2+2JD/(N+1)$. Drugi składnik
jest mniejszy od $\varepsilon/2$ dla dostatecznie dużego $N$.
\end{proof}

\begin{twierdzenie}[Abela w wersji dwuliniowej]
\label{ma2:thm:bilinear-abel}
Niech $T:E\times F\to G$ będzie ciągłe dwuliniowe,
a $(z_n)$ niech będzie ciągiem współczynników iloczynu
Cauchy'ego z definicji \ref{ma2:def:bilinear-cauchy}.
Jeżeli $\sum x_n=x$, $\sum y_n=y$ oraz $\sum z_n=z$
są zbieżne w odpowiednich przestrzeniach, to $z=T(x,y)$.
Nie potrzeba zupełności ani zbieżności bezwzględnej.
\end{twierdzenie}
\begin{proof}
Dla $N\ge0$ zachodzi
$\sum_{m=0}^N Z_m=\sum_{j=0}^N T(X_j,Y_{N-j})$.
Istotnie, dla każdej pary $r,s\ge0$ z $r+s\le N$
składnik $T(x_r,y_s)$ pojawia się po lewej w $Z_m$
dla $r+s\le m\le N$, po prawej dla $r\le j\le N-s$,
czyli po obu stronach dokładnie $N-r-s+1$ razy.
Pozostałe pary nie występują. Wszystkie sumy są skończone.
Dzielimy przez $N+1$. Prawa strona dąży do $T(x,y)$
na mocy lematu \ref{ma2:lem:bilinear-diagonal-means}.
Lewa strona dąży do $z$: stosujemy ten sam lemat
do odwzorowania $(w,a)\mapsto aw$ z $G\times\K$ do $G$,
które jest ciągłe dwuliniowe z oszacowaniem $K=1$,
oraz ciągów $Z_j\to z$ i $a_j=1$. Jednoznaczność granicy
daje tezę.
\end{proof}

\begin{przyklad}[Mnożenie macierzy]
\label{ma2:ex:matrix-cauchy-product}
W przestrzeni macierzy kwadratowych $d\times d$, $d\ge1$,
z normą operatorową indukowaną normą euklidesową,
odwzorowanie $T(U,V)=UV$ jest dwuliniowe i ciągłe:
dla każdego wektora $\xi$ mamy
$\|UV\xi\|\le\|U\|\|V\xi\|\le\|U\|\|V\|\|\xi\|$,
więc po wzięciu supremum po $\|\xi\|\le1$ otrzymujemy
$\|UV\|\le\|U\|\|V\|$. Jeśli szeregi macierzy $\sum_{j\ge0}A_j=A$ oraz
$\sum_{j\ge0}B_j=B$ są zbieżne i co najmniej jeden
jest bezwzględnie zbieżny,
twierdzenie \ref{ma2:thm:bilinear-mertens} daje
\[
 \sum_{n=0}^\infty\sum_{j=0}^n A_j B_{n-j}=AB
\]
dla wyrazów szeregów $A_j,B_j$ (nie ich sum częściowych).
W ogólności nie wolno zamienić kolejności mnożenia:
dla $d=2$ i macierzy jednostek $E_{12},E_{21}$ mamy
$E_{12}E_{21}=E_{11}\ne E_{22}=E_{21}E_{12}$, jak wynika
bezpośrednio z definicji mnożenia macierzy.
\end{przyklad}

\begin{wniosek}[Dowolne wyliczenie par w przypadku dwuliniowym]
\label{ma2:cor:bilinear-pair-enumeration}
Niech $T:E\times F\to G$ będzie ciągłe dwuliniowe.
Jeżeli szeregi $\sum x_n=x$ i $\sum y_n=y$ są bezwzględnie
zbieżne oraz $G$ jest przestrzenią Banacha, to dla każdej
bijekcji $\pi:\N_+\to\N_0\times\N_0$,
$\pi(l)=(j_l,k_l)$, szereg
$\sum_{l\ge1}T(x_{j_l},y_{k_l})$ jest bezwzględnie
zbieżny do $T(x,y)$.
\end{wniosek}
\begin{proof}
Wybierzmy $K\ge0$ z lematu \ref{ma2:lem:bilinear-bound}.
Dla dowolnego skończonego początkowego fragmentu wyliczenia
wszystkie pary należą do pewnego kwadratu
$\{0,\ldots,M\}^2$. Z iniektywności wyliczenia i nieujemności
norm otrzymujemy oszacowanie sumy norm tego fragmentu przez
\[
 K\left(\sum_{j=0}^\infty\|x_j\|\right)
  \left(\sum_{k=0}^\infty\|y_k\|\right)<\infty.
\]
Zatem każde takie wyliczenie daje szereg bezwzględnie
zbieżny, a z zupełności $G$ także zbieżny.
Dla wyliczenia po kolejnych przekątnych $j+k=n$,
opisanego w dowodzie twierdzenia
\ref{ma1:thm:pair-enumeration}, grupowanie kolejnych
bloków zachowuje sumę na mocy twierdzenia
\ref{ma2:thm:series-grouping}. Szereg zgrupowany jest
iloczynem Cauchy'ego względem $T$, a jego suma jest
równa $T(x,y)$ z wniosku \ref{ma2:cor:bilinear-absolute}.
Każde inne wyliczenie par jest przestawieniem
pierwszego; twierdzenie \ref{ma2:thm:absolute-rearrangement}
zapewnia zachowanie sumy.
\end{proof}

W następnym rozdziale przejdziemy do iloczynów nieskończonych,
czyli granic iloczynów częściowych, a później do rodzin
sumowalnych i szeregów wielokrotnych.

% Materiał podstawowy przejęty z AM I-12 z zachowaniem etykiet.
\chapter{Iloczyny nieskończone}
\label{ma1:ch:infinite-products}

W definicji szeregu przechodziliśmy do granicy sum częściowych. Analogicznie można badać granice iloczynów częściowych. Istnieje jednak ważna różnica: dodanie skończenie wielu składników nie może zmienić zbieżności szeregu, natomiast jeden czynnik równy zeru sprawia, że wszystkie dalsze iloczyny częściowe są zerowe. Odróżnimy zatem zbieżność do niezerowej granicy od zanikania iloczynów częściowych. Najpierw przeprowadzimy dowody bez odwoływania się do logarytmu, a następnie wykorzystamy funkcje elementarne z tomu I.

\section{Definicja i własności elementarne}
\begin{definicja}
\label{ma1:def:infinite-product}
Niech $K=\R$ lub $K=\C$, a $(u_n)_{n\ge1}$ będzie ciągiem w $K$. Iloczyn nieskończony $\prod_{n=1}^\infty u_n$ rozpatrujemy wraz z ciągiem iloczynów częściowych
\[
 P_0=1,\qquad P_N=\prod_{n=1}^N u_n\quad(N\in\N_+).
\]
Jeżeli wszystkie $u_n\ne0$, to iloczyn nazywamy zbieżnym, gdy istnieje granica $P_N\to P\in K\setminus\{0\}$. Wówczas piszemy $\prod_{n=1}^\infty u_n=P$. Jeżeli wszystkie czynniki są niezerowe, lecz nie istnieje skończona niezerowa granica $(P_N)$, iloczyn nazywamy rozbieżnym, również gdy $P_N\to0$. Jednoznaczność wartości iloczynu zbieżnego wynika z jednoznaczności granicy w przestrzeni metrycznej (tom I).
\end{definicja}
Przyjmujemy w tym rozdziale konwencję zbieżności do niezerowej granicy. Gdy występuje czynnik zerowy, wszystkie iloczyny częściowe od jego indeksu są zerowe; tę sytuację rozpatrujemy osobno. Można również stosować konwencję dopuszczającą skończenie wiele czynników zerowych i wymagającą zbieżności niezerowego ogona. W każdym przypadku trzeba odróżnić zwykłą granicę ciągu $(P_N)$ od zbieżności iloczynu w powyższym sensie.

\begin{twierdzenie}
\label{ma1:thm:product-elementary}
Niech $K=\R$ lub $K=\C$, a $(u_n)_{n\ge1},(v_n)_{n\ge1}$ będą ciągami w $K\setminus\{0\}$.
\begin{enumerate}[label=(\roman*)]
\item Dla dowolnego $m\in\N_+$, przy konwencji iloczynu pustego równego $1$, iloczyn $\prod_{n\ge1}u_n$ jest zbieżny wtedy i tylko wtedy, gdy $\prod_{n\ge m}u_n$ jest zbieżny. W przypadku zbieżności
$\prod_{n\ge1}u_n=(\prod_{n=1}^{m-1}u_n)(\prod_{n\ge m}u_n)$.
\item Jeżeli $\prod u_n=U$ i $\prod v_n=V$ są zbieżne, to $\prod(u_nv_n)=UV$ oraz $\prod u_n^{-1}=U^{-1}$.
\item Zbieżność $\prod u_n$ implikuje $u_n\to1$.
\end{enumerate}
\end{twierdzenie}
\begin{proof}
Oznaczmy $P_0=1$ oraz $P_N=\prod_{n=1}^N u_n$ dla $N\ge1$. Dla (i) iloczyn częściowy od $m$ do $N$ wynosi $P_N/P_{m-1}$, gdzie $P_{m-1}\ne0$. Jeśli $P_N\to P\ne0$, iloczyny od $m$ dążą do $P/P_{m-1}\ne0$. Odwrotnie, jeśli iloczyny od $m$ dążą do $T\ne0$, to $P_N=P_{m-1}\prod_{n=m}^N u_n\to P_{m-1}T\ne0$. Otrzymujemy również wzór na wartość iloczynu. Dla (ii) oznaczmy $Q_N=\prod_{n=1}^N v_n$. Iloczyny częściowe iloczynów o czynnikach $u_nv_n$ oraz $u_n^{-1}$ są odpowiednio $P_NQ_N$ i $1/P_N$. Rachunek granic daje granice $UV\ne0$ i $1/U\ne0$. Dla (iii) mamy $u_n=P_n/P_{n-1}$. Oba ciągi w ilorazie dążą do tego samego niezerowego $P$, zatem $u_n\to P/P=1$. W szczególności dla czynników rzeczywistych od pewnego indeksu $u_n>0$, ponieważ od pewnego indeksu $|u_n-1|<1/2$. Znak niezerowej granicy rzeczywistej jest znakiem iloczynu skończonego początku obejmującego wszystkie ujemne czynniki: dalsze iloczyny częściowe mają ten sam znak, a niezerowa granica zachowuje ten znak.
\end{proof}
\begin{przyklad}
\label{ma1:ex:telescoping-products}
Dla $n\ge2$ rozważmy trzy ciągi czynników (definicję iloczynu stosujemy po przeindeksowaniu). Z bezpośredniego skracania iloczynów skończonych dla $N\ge2$ otrzymujemy
\[
\begin{aligned}
 \prod_{n=2}^N(1-1/n)&=1/N,\\
 \prod_{n=2}^N(1+1/n)&=(N+1)/2,\\
 \prod_{n=2}^N(1-1/n^2)&=
 \prod_{n=2}^N\frac{n-1}{n}\frac{n+1}{n}
 =\frac{N+1}{2N}.
\end{aligned}
\]
Istotnie, w pierwszym iloczynie skracamy sąsiednie czynniki $(n-1)/n$, w drugim $(n+1)/n$, a trzeci jest iloczynem dwóch poprzednich, gdyż $(1-1/n)(1+1/n)=1-1/n^2$. W każdym przypadku czynniki dążą do $1$. W pierwszym przypadku iloczyny częściowe dążą do zera, w drugim do $+\infty$, a tylko w trzecim iloczyn jest zbieżny w sensie definicji \ref{ma1:def:infinite-product}, do $1/2$. Warunek konieczny zbieżności nie jest wystarczający.
\end{przyklad}

\section{Kryterium Cauchy'ego dla iloczynów}
\begin{twierdzenie}
\label{ma1:thm:product-cauchy}
Niech $(u_n)_{n\ge1}$ będzie ciągiem niezerowych liczb rzeczywistych lub zespolonych. Iloczyn $\prod_{n\ge1}u_n$ jest zbieżny wtedy i tylko wtedy, gdy
\[
 \forall\varepsilon\in\R_{>0}\ \exists N\in\N_+\
 \forall q,p\in\N_+\quad
 \left(q\ge p\ge N\ \Rightarrow\
 \left|\prod_{n=p}^{q}u_n-1\right|<\varepsilon\right).
\]
\end{twierdzenie}
\begin{proof}
Niech $P_0=1$, $P_n=\prod_{j=1}^n u_j$ dla $n\ge1$. Jeżeli $P_n\to P\ne0$, to od pewnego indeksu $|P_n-P|<|P|/2$, więc $|P_n|\ge|P|/2$. Dla $\varepsilon>0$ wybieramy $M$ tak, aby dla wszystkich $n\ge M$ zachodziły zarówno ta nierówność, jak i $|P_n-P|<\varepsilon|P|/4$. Gdy $q\ge p\ge M+1$, mamy
\[
 \left|\prod_{n=p}^q u_n-1\right|
 =\frac{|P_q-P_{p-1}|}{|P_{p-1}|}
 \le\frac{|P_q-P|+|P_{p-1}-P|}{|P|/2}
 <\varepsilon.
\]

Odwrotnie, stosujemy warunek z tezy najpierw dla $\varepsilon=1/2$. Istnieje $N_0\ge1$ takie, że dla każdego $q\ge N_0$ mamy $|\prod_{n=N_0}^q u_n-1|<1/2$. Ponieważ $P_{N_0-1}\ne0$ i $\prod_{n=N_0}^q u_n=P_q/P_{N_0-1}$, z nierówności trójkąta oraz nierówności odwrotnej otrzymujemy $1/2<|P_q/P_{N_0-1}|<3/2$. Stąd
\[
 \frac12|P_{N_0-1}|<|P_q|<\frac32|P_{N_0-1}|.
\]
Ciąg iloczynów częściowych jest więc ograniczony, również po uwzględnieniu skończonego początku; możemy przyjąć $B=\max(\{1,\frac32|P_{N_0-1}|\}\cup\{|P_n|:0\le n<N_0\})$. Wówczas $|P_n|\le B$ dla wszystkich $n\ge0$, także gdy $N_0=1$.

Dla dowolnego $\varepsilon>0$ wybieramy $N$ z warunku tezy dla dokładności $\varepsilon/B$. Dla $q>p\ge N$ otrzymujemy
\[
 |P_q-P_p|=|P_p|\left|\prod_{n=p+1}^q u_n-1\right|
 <B\frac{\varepsilon}{B}=\varepsilon.
\]
Przypadek $p=q$ jest trywialny, a dla $q<p$ zamieniamy indeksy. Z zupełności $\R$ i $\C$ (twierdzenie \ref{ma2:thm:finite}) ciąg $(P_n)$ ma granicę $P$. Dolne oszacowanie $|P_q|>|P_{N_0-1}|/2$ dla $q\ge N_0$ i zbieżność modułów dają $|P|\ge |P_{N_0-1}|/2>0$. Zatem iloczyn jest zbieżny.
\end{proof}

\section{Oszacowania skończonych iloczynów}
\begin{lemat}
\label{ma1:lem:product-finite-estimate}
Niech $m\in\N_+$ i niech $z_1,\ldots,z_m\in K$, gdzie $K=\R$ lub $K=\C$. Oznaczmy $s=\sum_{j=1}^m|z_j|$. Wówczas
\[
 \left|\prod_{j=1}^m(1+z_j)-1\right|
 \le \prod_{j=1}^m(1+|z_j|)-1.
\]
Jeżeli $s<1$, to
\[
 \prod_{j=1}^m(1+|z_j|)\le\frac1{1-s},
 \qquad
 \left|\prod_{j=1}^m(1+z_j)-1\right|
 \le\frac{s}{1-s}.
\]
\end{lemat}
\begin{proof}
Niech $I=\{1,\ldots,m\}$. Rozdzielność daje
\[
 \prod_{j=1}^m(1+z_j)=\sum_{E\subseteq I}\prod_{j\in E}z_j.
\]
Istotnie, z każdego czynnika wybieramy dokładnie jeden składnik: $z_j$ dla $j\in E$ lub $1$ dla $j\notin E$. Każdy podzbiór występuje dokładnie raz. Składnik odpowiadający zbiorowi pustemu wynosi $1$. Po odjęciu go i zastosowaniu nierówności trójkąta do pozostałych składników dostajemy pierwsze oszacowanie, gdyż suma ich modułów to dokładnie $\prod(1+|z_j|)-1$.

Niech $t_j=|z_j|\ge0$ oraz $s_k=\sum_{j=1}^k t_j$ dla $0\le k\le m$, gdzie $s_0=0$. Jeśli $s_m<1$, to $0\le s_k<1$ dla każdego $k$. Pokażemy indukcyjnie, że $\prod_{j=1}^k(1+t_j)\le1/(1-s_k)$. Dla $k=0$ zachodzi równość. W kroku indukcyjnym korzystamy z
\[
 (1-s_k)(1+t_k)
 =(1-s_{k-1}-t_k)(1+t_k)
 =1-s_{k-1}-s_{k-1}t_k-t_k^2
 \le1-s_{k-1}.
\]
Ponieważ oba mianowniki są dodatnie, wynika stąd $(1+t_k)/(1-s_{k-1})\le1/(1-s_k)$. Z założenia indukcyjnego, mnożąc przez dodatnie $1+t_k$, otrzymujemy
$\prod_{j=1}^k(1+t_j)\le(1+t_k)/(1-s_{k-1})\le1/(1-s_k)$, co kończy krok indukcyjny. Ostatnie oszacowanie wynika z pierwszego oraz $1/(1-s)-1=s/(1-s)$.
\end{proof}

\begin{lemat}
\label{ma1:lem:positive-product-bound}
Niech $m\in\N_+$ i niech $t_1,\ldots,t_m$ będą nieujemnymi liczbami rzeczywistymi. Wówczas
$\prod_{j=1}^m(1+t_j)\ge1+\sum_{j=1}^m t_j$.
Jeżeli ponadto $t_j\le1$ dla każdego $j$, to
$\prod_{j=1}^m(1-t_j)\ge1-\sum_{j=1}^m t_j$.
\end{lemat}
\begin{proof}
Pierwsza nierówność wynika z rozwinięcia iloczynu: składniki stopnia co najmniej dwa są nieujemne. Dla drugiej przeprowadzamy indukcję. Gdy $m=1$, zachodzi równość. Z założenia indukcyjnego, po pomnożeniu przez nieujemne $1-t_{m+1}$, mamy
\[
 \prod_{j=1}^{m+1}(1-t_j)
 \ge (1-\sum_{j=1}^m t_j)(1-t_{m+1})
 =1-\sum_{j=1}^{m+1}t_j+t_{m+1}\sum_{j=1}^m t_j
 \ge1-\sum_{j=1}^{m+1}t_j.
\]
\end{proof}

\section{Zbieżność bezwzględna}
\begin{definicja}
\label{ma1:def:absolute-product}
Niech $(a_n)_{n\ge1}$ będzie ciągiem liczb rzeczywistych lub zespolonych takim, że $1+a_n\ne0$ dla każdego $n$. Iloczyn $\prod_{n\ge1}(1+a_n)$ nazywamy bezwzględnie zbieżnym, jeżeli szereg $\sum_{n\ge1}|a_n|$ jest zbieżny.
\end{definicja}
\begin{twierdzenie}
\label{ma1:thm:absolute-product}
Niech $(a_n)_{n\ge1}$ będzie ciągiem w $\R$ lub $\C$, takim że $1+a_n\ne0$ dla każdego $n\ge1$. Jeśli szereg $\sum_{n=1}^\infty|a_n|$ jest zbieżny, to $\prod(1+a_n)$ jest zbieżny do niezerowej granicy $P$. Ponadto, oznaczając $P_N=\prod_{n=1}^N(1+a_n)$ i $P_0=1$, dla każdego całkowitego $N\ge0$ takiego, że $s_N=\sum_{n>N}|a_n|<1$, otrzymujemy
\[
 \left|\frac{P}{P_N}-1\right|\le\frac{s_N}{1-s_N}.
\]
\end{twierdzenie}
\begin{proof}
Ustalmy $\varepsilon>0$ i wybierzmy rzeczywiste $\delta$ z $0<\delta<\min(1/2,\varepsilon/2)$. Ponieważ szereg modułów jest zbieżny, istnieje $N_0$ z $\sum_{n\ge N_0}|a_n|<\delta$. Dla $q\ge p\ge N_0$ mamy $s=\sum_{n=p}^q|a_n|<\delta$. Z lematu \ref{ma1:lem:product-finite-estimate} wynika
\[
 \left|\prod_{n=p}^q(1+a_n)-1\right|
 \le\frac{s}{1-s}\le2s<\varepsilon.
\]
Kryterium \ref{ma1:thm:product-cauchy} daje zbieżność iloczynu do niezerowego $P$.

Dla ustalonego $N$ z $s_N<1$ oraz $q>N$ mamy $\sum_{n=N+1}^q|a_n|\le s_N$, a funkcja $s\mapsto s/(1-s)$ jest rosnąca na $[0,1)$, gdyż dla $0\le s\le t<1$ nierówność $s/(1-s)\le t/(1-t)$ po pomnożeniu przez dodatnie mianowniki jest równoważna $s\le t$. Z lematu otrzymujemy $|P_q/P_N-1|\le s_N/(1-s_N)$. Ponieważ $P_N\ne0$ i $P_q\to P$, rachunek granic oraz zbieżność modułów dają $|P_q/P_N-1|\to|P/P_N-1|$. Nieostra nierówność zachowuje się w granicy, co daje tezę.
\end{proof}
\begin{przyklad}
Dla $a_n=2^{-n}$ mamy $s_N=2^{-N}$, więc dla $N\ge1$ iloczyn $\prod_{n=1}^\infty(1+2^{-n})$ spełnia
$|P/P_N-1|\le1/(2^N-1)$. Dla $N=10$ jest to $1/1023<1/1000$. Możemy więc oszacować błąd względny względem obliczonego iloczynu częściowego bez znajomości dokładnej wartości $P$.
\end{przyklad}

\section{Czynniki o ustalonym znaku odchylenia od jedynki}
\begin{twierdzenie}
\label{ma1:thm:positive-product-series}
Niech $(t_n)_{n\ge1}$ będzie ciągiem nieujemnych liczb rzeczywistych.
\begin{enumerate}[label=(\roman*)]
\item Iloczyn $\prod_{n\ge1}(1+t_n)$ jest zbieżny wtedy i tylko wtedy, gdy szereg $\sum_{n\ge1}t_n$ jest zbieżny. Jeśli szereg jest rozbieżny, iloczyny częściowe dążą do $+\infty$.
\item Jeżeli dodatkowo $t_n<1$ dla wszystkich $n$, to iloczyn $\prod_{n\ge1}(1-t_n)$ jest zbieżny wtedy i tylko wtedy, gdy szereg $\sum_{n\ge1}t_n$ jest zbieżny. Jeśli szereg jest rozbieżny, iloczyny częściowe dążą do zera.
\end{enumerate}
\end{twierdzenie}
\begin{proof}
Zbieżność szeregu $\sum t_n$ daje w obu przypadkach zbieżność bezwzględną iloczynu, więc jego zbieżność wynika z twierdzenia \ref{ma1:thm:absolute-product}.

Załóżmy rozbieżność szeregu $\sum t_n$. Ponieważ $t_n\ge0$, jego sumy częściowe dążą do $+\infty$. Z lematu \ref{ma1:lem:positive-product-bound} mamy
$\prod_{n=1}^N(1+t_n)\ge1+\sum_{n=1}^N t_n$.
Prawa strona dąży do $+\infty$, co daje tezę w (i).

Dla (ii) zauważmy, że $1/(1-t_n)\ge1+t_n$: po pomnożeniu przez dodatnie $1-t_n$ otrzymujemy równoważną nierówność $1\ge1-t_n^2$. Stąd
\[
 0<\prod_{n=1}^N(1-t_n)
 \le\frac1{\prod_{n=1}^N(1+t_n)}
 \le\frac1{1+\sum_{n=1}^N t_n}.
\]
Dla każdego $\varepsilon>0$ z rozbieżności sum częściowych do $+\infty$ wynika, że od pewnego indeksu $1+\sum_{n=1}^N t_n>1/\varepsilon$. Ostatnia wielkość jest więc mniejsza od $\varepsilon$, a z podwójnego oszacowania iloczyny częściowe dążą do zera. Iloczyn nie jest zatem zbieżny w sensie przyjętej definicji.
\end{proof}
\begin{uwaga}
Twierdzenie daje równoważność zbieżności bezwzględnej iloczynu $\prod(1+a_n)$ oraz zbieżności iloczynu nieujemnych czynników $\prod(1+|a_n|)$, pod warunkiem że wszystkie czynniki wyjściowego iloczynu są niezerowe. Nie utożsamiamy tej własności ze zbieżnością iloczynu $\prod |1+a_n|$; na przykład dla $a_n=-2$ dla wszystkich $n$ mamy $\prod |1+a_n|=1$, podczas gdy $\prod(1+a_n)$ ma naprzemiennie iloczyny częściowe $-1$ i $1$, a $\sum|a_n|$ jest rozbieżny.
\end{uwaga}

\section{Zmiana kolejności czynników}
\begin{twierdzenie}
\label{ma1:thm:product-rearrangement}
Niech $(a_n)$ będzie ciągiem liczb rzeczywistych lub zespolonych, $1+a_n\ne0$ dla każdego $n$, i niech $\sum|a_n|$ będzie zbieżny. Dla każdej bijekcji $\pi:\N_+\to\N_+$ iloczyn $\prod_{k\ge1}(1+a_{\pi(k)})$ jest bezwzględnie zbieżny i ma tę samą wartość co $\prod_{n\ge1}(1+a_n)$.
\end{twierdzenie}
\begin{proof}
Z twierdzenia \ref{ma2:thm:absolute-rearrangement} szereg $\sum_{k\ge1}|a_{\pi(k)}|$ jest zbieżny. Z twierdzenia \ref{ma1:thm:absolute-product} oba iloczyny mają więc niezerowe granice; oznaczmy granicę iloczynu przestawionego przez $Q$, a granicę iloczynu wyjściowego przez $P$. Niech $P_N=\prod_{n=1}^N(1+a_n)$.
Niech $s_N=\sum_{n>N}|a_n|$. Wybierzmy $N$ tak duże, aby $s_N<1$. Z surjektywności $\pi$ istnieje $K_0=\max\{\pi^{-1}(1),\ldots,\pi^{-1}(N)\}$. Dla $K\ge K_0$ zbiór $F_K=\{\pi(1),\ldots,\pi(K)\}$ zawiera wszystkie indeksy od $1$ do $N$. Z przemienności skończonego mnożenia i iniektywności $\pi$ wynika
\[
 \prod_{k=1}^K(1+a_{\pi(k)})
 =P_N\prod_{n\in F_K\setminus\{1,\ldots,N\}}(1+a_n).
\]
Suma modułów $a_n$ po drugim zbiorze jest nie większa od $s_N$. Lemat \ref{ma1:lem:product-finite-estimate} daje zatem
\[
 \left|\frac{\prod_{k=1}^K(1+a_{\pi(k)})}{P_N}-1\right|
 \le\frac{s_N}{1-s_N}.
\]
Dla pustego zbioru drugi iloczyn wynosi $1$, więc oszacowanie również zachodzi. Przechodząc do granicy po $K$, otrzymujemy $|Q/P_N-1|\le s_N/(1-s_N)$. Teraz $P_N\to P\ne0$ i $s_N\to0$, więc $|Q/P-1|=0$, czyli $Q=P$.
\end{proof}

\section{Zbieżność warunkowa i związki z szeregami}
\begin{definicja}
\label{ma1:def:conditional-product}
Niech $(u_n)_{n\ge1}$ będzie ciągiem niezerowych liczb rzeczywistych lub zespolonych. Iloczyn $\prod_{n=1}^\infty u_n$ nazywamy zbieżnym warunkowo, jeżeli jest zbieżny do niezerowej granicy, lecz szereg $\sum_{n=1}^\infty|u_n-1|$ jest rozbieżny.
\end{definicja}
Nie można w ogólności zastąpić szeregu modułów w twierdzeniu \ref{ma1:thm:absolute-product} zwykłym szeregiem $\sum a_n$. Poniższe przykłady rozróżniają obie własności.

\begin{przyklad}
\label{ma1:ex:conditional-product}
Dla $k\ge1$ przyjmijmy $u_{2k-1}=1+1/k$ i $u_{2k}=k/(k+1)$. Wszystkie czynniki są dodatnie, a każdy iloczyn dwóch kolejnych czynników w parze wynosi $1$. Zatem
\[
 P_{2k}=1,\qquad P_{2k-1}=1+1/k.
\]
Oba podciągi dążą do $1$. Dla dowolnej dokładności $\varepsilon>0$ istnieje wspólny indeks $K$ taki, że oba rodzaje iloczynów częściowych są w odległości mniejszej niż $\varepsilon$ od $1$ dla $k\ge K$. Każdy $N\ge2K$ jest jednym z takich indeksów, więc $P_N\to1$. Jeśli $a_n=u_n-1$, to $a_{2k-1}=1/k$ oraz $a_{2k}=-1/(k+1)$. Szereg $\sum|a_n|$ jest rozbieżny, gdyż jego sumy częściowe o parzystych indeksach są nie mniejsze od $\sum_{k=1}^m1/k$ dla odpowiedniego $m$. Iloczyn jest zbieżny, lecz nie bezwzględnie, czyli jest zbieżny warunkowo. Zauważmy, że $\sum a_n$ jest w tym przykładzie również zbieżny: jego sumy częściowe o indeksie $2k$ wynoszą $\sum_{j=1}^k(1/j-1/(j+1))=1-1/(k+1)$, a sumy o indeksie $2k-1$ wynoszą $1-1/k+1/k=1$. Stąd suma szeregu wynosi $1$.
\end{przyklad}

\begin{przyklad}
\label{ma1:ex:conditional-product-rearranged}
Przestawmy czynniki z przykładu \ref{ma1:ex:conditional-product}, biorąc po dwa kolejne czynniki większe od $1$, a następnie jeden kolejny czynnik mniejszy od $1$. W $k$-tym bloku występują czynniki o indeksach $4k-3,4k-1,2k$. Każdy dodatni indeks parzysty ma jednoznacznie postać $2k$, a każdy dodatni indeks nieparzysty ma jednoznacznie jedną z postaci $4k-3$, $4k-1$, gdzie $k\ge1$. Te trzy rodziny są rozłączne, więc rzeczywiście otrzymujemy bijekcję indeksów.

Oznaczmy iloczyny częściowe przestawionego iloczynu przez $Q_N$. Po $m$ pełnych blokach, dla $m\ge1$, mamy
\[
 Q_{3m}=\prod_{j=1}^{2m}\frac{j+1}{j}
          \prod_{j=1}^{m}\frac{j}{j+1}
        =\frac{2m+1}{m+1}\longrightarrow2.
\]
Ponadto $Q_{3m+1}=Q_{3m}(1+1/(2m+1))$ i $Q_{3m+2}=Q_{3m+1}(1+1/(2m+2))$. Rachunek granic daje zbieżność obu tych podciągów do $2$. Dla dowolnej dokładności $\varepsilon>0$ wybieramy wspólny próg $M$ dla wszystkich trzech podciągów. Każde $N\ge3M$ ma postać $3m+r$, gdzie $m\ge M$ i $r\in\{0,1,2\}$, zatem $|Q_N-2|<\varepsilon$. Otrzymujemy iloczyn zbieżny do $2$, podczas gdy iloczyn wyjściowy był równy $1$. Założenia bezwzględnej zbieżności w twierdzeniu \ref{ma1:thm:product-rearrangement} nie można więc zastąpić samą zbieżnością.
\end{przyklad}

\begin{przyklad}
\label{ma1:ex:series-converges-product-zero}
Dla $k\ge1$ połóżmy $a_{2k-1}=1/\sqrt{k+1}$ i $a_{2k}=-1/\sqrt{k+1}$. Wszystkie czynniki $1+a_n$ są dodatnie. Parzyste sumy częściowe szeregu $\sum a_n$ wynoszą $0$, a nieparzyste $1/\sqrt{k+1}\to0$. Szereg jest zbieżny. Jednak
\[
 (1+a_{2k-1})(1+a_{2k})=1-\frac1{k+1}
 =\frac{k}{k+1},
 \qquad P_{2m}=\frac1{m+1}.
\]
Ponadto $P_{2m-1}=P_{2m-2}(1+1/\sqrt{m+1})=(1+1/\sqrt{m+1})/m\to0$. Wszystkie iloczyny częściowe dążą do zera, więc iloczyn jest rozbieżny w sensie definicji \ref{ma1:def:infinite-product}.
\end{przyklad}

\begin{przyklad}
\label{ma1:ex:product-converges-series-diverges}
Niech $t_k=1/\sqrt{k}$ dla $k\ge1$ oraz $u_{2k-1}=1+t_k$, $u_{2k}=1/(1+t_k)$. Ponownie $P_{2k}=1$, $P_{2k-1}=1+t_k\to1$, więc iloczyn jest zbieżny do $1$. Natomiast dla $a_n=u_n-1$ mamy
\[
 a_{2k-1}+a_{2k}=t_k-\frac{t_k}{1+t_k}
 =\frac{t_k^2}{1+t_k}\ge\frac1{2k},
\]
gdyż $0<t_k\le1$. Parzyste sumy częściowe szeregu $\sum a_n$ dążą do $+\infty$ z rozbieżności szeregu harmonicznego. Nieparzysta suma częściowa jest większa od poprzedniej parzystej o dodatni $t_k$. Stąd cały ciąg sum częściowych dąży do $+\infty$, mimo zbieżności iloczynu.
\end{przyklad}

\begin{uwaga}
Przykłady pokazują, że warunek $\sum a_n$ zbieżny nie jest ani konieczny, ani wystarczający do zbieżności $\prod(1+a_n)$. Poniżej wykorzystamy funkcje elementarne skonstruowane już w tomie I, aby uzyskać ścisły związek między iloczynami i szeregami logarytmów.
\end{uwaga}


\section{Logarytmy czynników rzeczywistych i zespolonych}
Korzystamy z funkcji wykładniczej, logarytmu rzeczywistego
oraz funkcji trygonometrycznych skonstruowanych w tomie I
(rozdział 18, Funkcje elementarne). W szczególności
$\exp(z+w)=\exp z\exp w$, funkcja $\exp$ jest ciągła
i nigdzie niezerowa, a $\log:(0,\infty)\to\R$ jest ciągłą
odwrotnością rzeczywistej funkcji wykładniczej.
Nie korzystamy z teorii różniczkowania funkcji zespolonych
ani z twierdzeń o szeregach potęgowych z dalszych rozdziałów.

\begin{twierdzenie}[Iloczyny dodatnie]
\label{ma2:thm:positive-log-product}
Dla ciągu dodatnich liczb rzeczywistych $(u_n)$
iloczyn $\prod_{n\ge1}u_n$ jest zbieżny wtedy i tylko wtedy,
gdy szereg $\sum_{n\ge1}\log u_n$ jest zbieżny.
W przypadku zbieżności
\[
 \prod_{n=1}^{\infty}u_n
 =\exp\left(\sum_{n=1}^{\infty}\log u_n\right).
\]
Jeżeli sumy częściowe szeregu logarytmów dążą do $+\infty$
lub $-\infty$, iloczyny częściowe dążą odpowiednio do
$+\infty$ lub $0$. Odwrotne implikacje również zachodzą.
\end{twierdzenie}
\begin{proof}
Dla każdego $N$ iloczyn częściowy $P_N$ jest dodatni,
a prawo logarytmowania iloczynu, zastosowane indukcyjnie,
daje $\log P_N=\sum_{n=1}^N\log u_n$.
Równoważnie $P_N=\exp(\sum_{n=1}^N\log u_n)$.
Jeżeli $P_N\to P\ne0$, to $P>0$ i z ciągłości logarytmu
sumy częściowe logarytmów dążą do $\log P$.
Jeśli sumy częściowe dążą do $S\in\R$, to z ciągłości
funkcji wykładniczej $P_N\to\exp S>0$.
Te same tożsamości oraz granice funkcji $\exp$ w obu
nieskończonościach i logarytmu w $0^+$ oraz $+\infty$
(tom I, rozdział 18) dowodzą pozostałych równoważności.
\end{proof}

Dla liczb zespolonych użyjemy wyłącznie logarytmu
określonego w prawej półpłaszczyźnie. Oznaczmy
\[
 L(w)=\log |w|+\mathrm i\arctan
 \frac{\operatorname{Im}w}{\operatorname{Re}w}
 \qquad(\operatorname{Re}w>0).
\]
Tutaj $\arctan$ przyjmuje wartości w $(-\pi/2,\pi/2)$.
Jest to funkcja ciągła, $L(1)=0$ i $\exp L(w)=w$.
Istotnie, jeśli $w=x+\mathrm i y$, $x>0$,
$\theta=\arctan(y/x)$, to $\cos\theta>0$ oraz
$\sin\theta=(y/x)\cos\theta$. Z tożsamości
$\sin^2\theta+\cos^2\theta=1$ dostajemy
$\cos\theta=x/|w|$ i $\sin\theta=y/|w|$.
Wzór Eulera z tomu I daje
$\exp L(w)=|w|(\cos\theta+\mathrm i\sin\theta)=w$.
Dla rzeczywistych $w>0$ mamy $L(w)=\log w$.

\begin{twierdzenie}[Kryterium logarytmiczne w przypadku zespolonym]
\label{ma2:thm:complex-log-product}
Niech $(u_n)$ będzie ciągiem niezerowych liczb zespolonych
takim, że $u_n\to1$. Wybierzmy $m\ge1$ tak, aby
$\operatorname{Re}u_n>0$ dla $n\ge m$.
Iloczyn $\prod_{n\ge1}u_n$ jest zbieżny wtedy i tylko wtedy,
gdy szereg $\sum_{n\ge m}L(u_n)$ jest zbieżny w $\C$.
W przypadku zbieżności
\[
 P=\left(\prod_{n=1}^{m-1}u_n\right)
       \exp\left(\sum_{n=m}^{\infty}L(u_n)\right).
\]
\end{twierdzenie}
\begin{proof}
Dla $N\ge m$ prawo dodawania argumentów funkcji
wykładniczej daje
\[
 \exp S_N=R_N,\qquad
 S_N=\sum_{n=m}^N L(u_n),\qquad R_N=\prod_{n=m}^N u_n .
\]
Zbieżność $S_N$ implikuje więc zbieżność $R_N$ do
niezerowej granicy. Twierdzenie
\ref{ma1:thm:product-elementary} pozwala dołączyć
skończony początek i daje wzór na $P$.

Odwrotnie, niech $R_N\to R\ne0$. Dla dostatecznie
dużego $N$ mamy $\operatorname{Re}(R_N/R)>0$,
a z ciągłości $L$ otrzymujemy $b_N=L(R_N/R)\to0$.
Wybierzmy $c\in\C$ z $\exp c=R$; istnienie takiej liczby
wynika z postaci trygonometrycznej niezerowej liczby
zespolonej (tom I, rozdział 18): jeśli
$R=|R|(\cos\theta+\mathrm i\sin\theta)$, przyjmujemy
$c=\log|R|+\mathrm i\theta$.
Wtedy $\exp(S_N-c-b_N)=1$.
Z prawa dodawania argumentów i wzoru Eulera wynika
$\exp z=1$ dokładnie wtedy, gdy $z\in2\pi\mathrm i\Z$:
dla $z=x+\mathrm i y$ moduł wykładniczej jest $\exp x$,
więc $x=0$, a $\cos y=1$ i $\sin y=0$ zachodzą
dokładnie dla całkowitych wielokrotności $2\pi$.
Istnieją zatem całkowite $k_N$ z
\[
 S_N=c+b_N+2\pi\mathrm i k_N .
\]
Dla kolejnych dostatecznie dużych indeksów mamy
\[
 2\pi\mathrm i(k_{N+1}-k_N)
 =L(u_{N+1})-(b_{N+1}-b_N)\longrightarrow0.
\]
Istotnie, $L(u_{N+1})\to L(1)=0$ oraz $b_N\to0$.
Lewa strona ma moduł równy zero lub co najmniej $2\pi$,
więc od pewnego indeksu $k_{N+1}=k_N$. Indukcyjnie
$k_N$ jest wtedy stały. Stąd $S_N$ ma skończoną granicę,
co kończy dowód.
\end{proof}

\section{Zbieżność bezwarunkowa iloczynów}
\begin{twierdzenie}[Kryterium iloczynów po zbiorach skończonych]
\label{ma2:thm:product-finite-tails}
Niech $(u_n)$ będzie ciągiem niezerowych skalarów.
Następujące warunki są równoważne:
\begin{enumerate}[label=(\roman*)]
\item każde przestawienie iloczynu jest zbieżne do
niezerowej granicy (nie zakładamy równości granic);
\item dla każdego $\varepsilon>0$ istnieje $N\ge1$ takie,
że dla każdego skończonego $H\subseteq\{N+1,N+2,\ldots\}$
zachodzi
\[
 \left|\prod_{n\in H}u_n-1\right|<\varepsilon;
\]
\item istnieje $P\ne0$ takie, że dla każdego $\varepsilon>0$
istnieje skończony zbiór $F_0\subseteq\N_+$, dla którego
$|\prod_{n\in F}u_n-P|<\varepsilon$ dla wszystkich
skończonych $F\supseteq F_0$.
\end{enumerate}
W przypadku spełnienia tych warunków wszystkie
przestawienia mają tę samą wartość $P$.
\end{twierdzenie}
\begin{proof}
Iloczyn po zbiorze pustym przyjmujemy równy $1$.
Gdyby (ii) nie zachodziło, istniałoby $\varepsilon>0$
takie, że poza każdym skończonym początkiem można wybrać
niepusty skończony zbiór $H$ z
$|\prod_{n\in H}u_n-1|\ge\varepsilon$.
Konstruujemy wyliczenie etapami jak w dowodzie twierdzenia
\ref{ma2:thm:unconditional-finite}: w każdym etapie
wybieramy taki zbiór poza maksimum dotychczas użytych
indeksów i numerem etapu, wypisujemy jego elementy
rosnąco, następnie wypisujemy najmniejszy dotychczas
niewypisany indeks. Indeksy nie powtarzają się, a po
$k$ etapach wypisano wszystkie indeksy $1,\ldots,k$
(indukcyjnie najmniejszy brakujący indeks jest wtedy
nie większy niż $k$). Otrzymujemy bijekcję.
Iloczyny po kolejnych wybranych blokach nie spełniają
warunku Cauchy'ego z twierdzenia
\ref{ma1:thm:product-cauchy}, ponieważ ich początkowe
pozycje w nowym wyliczeniu dążą do nieskończoności.
Iloczyn przestawiony jest rozbieżny, wbrew (i).
Zatem (i) implikuje (ii).

Załóżmy (ii). W szczególności dla kolejnych przedziałów
indeksów warunek Cauchy'ego daje zbieżność zwykłego
iloczynu do pewnego $P\ne0$. Oznaczmy jego iloczyny
częściowe przez $P_N$. Istnieje $B\ge1$ takie, że
$|P_N|\le B$ dla wszystkich $N$.
Dla $\varepsilon>0$ wybieramy $N$ tak duże, aby
$|P_N-P|<\varepsilon/2$ oraz warunek (ii) był spełniony
z dokładnością $\varepsilon/(2B)$.
Jeśli $F\supseteq\{1,\ldots,N\}$ jest skończony, to
\[
 \left|\prod_{n\in F}u_n-P\right|
 \le |P_N|\left|\prod_{n\in F\setminus\{1,\ldots,N\}}u_n-1\right|
      +|P_N-P|<\varepsilon.
\]
Otrzymujemy (iii).
Wreszcie dla dowolnej bijekcji $\pi$ każdy skończony
$F_0$ z (iii) jest zawarty w $\{\pi(1),\ldots,\pi(k)\}$
dla dostatecznie dużych $k$. Z (iii) iloczyny częściowe
przestawienia dążą do tego samego niezerowego $P$.
Dowodzi to (i) i równości wartości.
\end{proof}

\begin{uwaga}
Bezwzględna zbieżność iloczynu implikuje (ii):
jeżeli $s_N=\sum_{n>N}|u_n-1|<1$, to lemat
\ref{ma1:lem:product-finite-estimate} ogranicza
odchylenie od $1$ przez $s_N/(1-s_N)$, niezależnie
od wybranego skończonego zbioru indeksów.
Poniżej wykażemy również implikację odwrotną.
\end{uwaga}

\begin{lemat}
\label{ma2:lem:log-comparable}
Istnieje $\rho\in(0,1/2]$ takie, że dla $a\in\C$
z $|a|\le\rho$ zachodzi
\[
 \tfrac12|a|\le |L(1+a)|\le 2|a|.
\]
\end{lemat}
\begin{proof}
Dla $|a|\le1/2$ mamy $\operatorname{Re}(1+a)>0$.
Z ciągłości $L$ w $1$ i $L(1)=0$ wybieramy
$\rho\in(0,1/2]$ tak, aby $|L(1+a)|\le1/2$ dla
$|a|\le\rho$. Dla $b=L(1+a)$ zachodzi
$\exp b-1=a$. Z definicji wykładniczej
w tomie I, dla $|b|\le1/2$,
\[
 |\exp b-1-b|
 \le\sum_{k=2}^\infty\frac{|b|^k}{k!}
 \le |b|^2\sum_{k=2}^\infty (1/2)^{k-2}/2
 =|b|^2.
\]
Nierówność trójkąta i nierówność odwrotna dają
$\tfrac12|b|\le|a|\le\tfrac32|b|$, gdyż $|b|\le1/2$.
Stąd $\tfrac12|a|\le |b|\le2|a|$.
\end{proof}

\begin{twierdzenie}[Zbieżność bezwarunkowa a bezwzględna]
\label{ma2:thm:product-unconditional-absolute}
Dla niezerowych skalarów rzeczywistych lub zespolonych
$u_n$ iloczyn $\prod u_n$ jest bezwzględnie zbieżny wtedy
i tylko wtedy, gdy każde jego przestawienie jest zbieżne
do niezerowej granicy.
\end{twierdzenie}
\begin{proof}
Pierwsza implikacja wynika z twierdzenia
\ref{ma1:thm:product-rearrangement}.
Odwrotnie, załóżmy zbieżność wszystkich przestawień.
Zbieżność iloczynu w kolejności wyjściowej daje
$u_n\to1$. Wybierzmy $m$ tak duże, aby
$|u_n-1|\le\rho$ dla $n\ge m$, gdzie $\rho$ jest
z lematu \ref{ma2:lem:log-comparable}.
Każde przestawienie ogona jest zbieżne: poprzedzamy
je czynnikami $u_1,\ldots,u_{m-1}$, otrzymując
przestawienie całego iloczynu, po czym stosujemy
twierdzenie \ref{ma1:thm:product-elementary}.
Kryterium \ref{ma2:thm:complex-log-product} zastosowane
do każdego przestawienia ogona (po przeindeksowaniu)
daje zbieżność każdego przestawienia szeregu
$\sum_{n\ge m}L(u_n)$. W przypadku rzeczywistym możemy
traktować liczby rzeczywiste jako zespolone, bez zmiany
modułu i zbieżności.
Z twierdzenia \ref{ma2:thm:unconditional-finite} wynika
zbieżność bezwarunkowa szeregu logarytmów, a z twierdzenia
\ref{ma2:thm:finite-unconditional-absolute} jego
zbieżność bezwzględna, ponieważ $\C$ jest przestrzenią
skończenie wymiarową. Z lematu
\ref{ma2:lem:log-comparable} wynika
$|u_n-1|\le2|L(u_n)|$ dla $n\ge m$.
Kryterium porównawcze daje zbieżność
$\sum_{n\ge m}|u_n-1|$. Dołączenie skończonego początku
nie zmienia zbieżności szeregu norm.
\end{proof}

W następnym rozdziale przejdziemy do rodzin sumowalnych
i szeregów wielokrotnych.

% Rozwinięcie AM II-03 i AM II-04; dyskretne odpowiedniki twierdzeń z TM-12.
\chapter{Rodziny sumowalne i szeregi wielokrotne}
\label{ma2:ch:summable-families}

W rozdziale \ref{ma2:ch:unconditional-series} badaliśmy
sumy po skończonych zbiorach indeksów szeregu. Teraz
zbiór indeksów nie musi być przeliczalny ani uporządkowany.
Rozwiniemy teorię bez odwoływania się do całki Lebesgue'a.
W przypadku przeliczalnych zbiorów indeksów otrzymamy
dyskretne odpowiedniki twierdzeń Tonellego i Fubiniego
(teoria miary, rozdział 12, etykiety
\texttt{thm:tonelli-12} i \texttt{thm:fubini-12}).
Dowody ogólnych twierdzeń miarowych pozostają w tamtym skrypcie.

\section{Sumy rodzin nieujemnych}

\begin{definicja}
\label{ma2:def:positive-family}
Dla dowolnego zbioru $I$ oraz rodziny
$(a_i)_{i\in I}$ o wartościach w $[0,+\infty]$ określamy
\[
 \sum_{i\in I}a_i
 :=\sup\left\{\sum_{i\in F}a_i:
                 F\subseteq I,\ F\text{ skończony}\right\}.
\]
Suma pusta wynosi $0$. Jeżeli skończona suma zawiera
składnik $+\infty$, przyjmuje wartość $+\infty$.
\end{definicja}
Jeżeli $I=\N_+$, definicja zgadza się z granicą
niemalejących sum częściowych. Istotnie, każdy skończony
zbiór indeksów jest zawarty w pewnym początkowym
przedziale, a każdy taki przedział jest dopuszczalny
w supremum. Jest to rozszerzenie definicji z AM I,
rozdział 09, sekcja o szeregach nieujemnych.

\begin{lemat}[Monotoniczność i skończony podział]
\label{ma2:lem:positive-family-partition}
Jeżeli $a_i\le b_i$ dla wszystkich $i\in I$, to
$\sum_{i\in I}a_i\le\sum_{i\in I}b_i$.
Jeżeli $I=A\cup B$, $A\cap B=\varnothing$, to
\[
 \sum_{i\in I}a_i=\sum_{i\in A}a_i+\sum_{i\in B}a_i.
\]
\end{lemat}
\begin{proof}
Pierwsza teza wynika z porównania każdej skończonej sumy.
Dla drugiej oznaczmy sumy po $A,B$ przez $\alpha,\beta$.
Rozkładając każdy skończony $F\subseteq I$ na rozłączne
części $F\cap A$ i $F\cap B$, otrzymujemy
$\sum_{i\in F}a_i\le\alpha+\beta$. Bierzemy supremum po $F$.
Jeżeli $\alpha,\beta$ są skończone, to dla każdego
$\varepsilon>0$ wybieramy skończone $F_A\subseteq A$,
$F_B\subseteq B$ o sumach większych niż
$\alpha-\varepsilon/2$ i $\beta-\varepsilon/2$.
Ich suma zbiorów jest skończona i rozłączna, więc
$\sum_{i\in I}a_i>\alpha+\beta-\varepsilon$.
Dowolność $\varepsilon$ daje nierówność odwrotną.
Jeżeli któraś z liczb $\alpha,\beta$ jest nieskończona,
to sumy po odpowiednich skończonych podzbiorach
są dowolnie duże (lub nieskończone), a zatem
$\sum_{i\in I}a_i=+\infty$.
\end{proof}

\begin{twierdzenie}[Przeliczalność nośnika i oszacowanie ogona]
\label{ma2:thm:positive-family-support}
Jeżeli $\alpha=\sum_{i\in I}a_i<\infty$, to zbiór
$\{i\in I:a_i>0\}$ jest co najwyżej przeliczalny.
Ponadto dla każdego $\varepsilon>0$ istnieje skończony
$F_0\subseteq I$ taki, że
$\sum_{i\in I\setminus F_0}a_i<\varepsilon$.
\end{twierdzenie}
\begin{proof}
Dla $m\ge1$ zbiór $A_m=\{i:a_i\ge1/m\}$ jest skończony:
gdyby zawierał $r$ różnych elementów przy $r>m\alpha$,
suma po nich byłaby co najmniej $r/m>\alpha$.
Każde dodatnie $a_i$ należy do jednego z tych progów,
więc $\{i:a_i>0\}=\bigcup_{m\ge1}A_m$.
Jest to przeliczalna suma zbiorów skończonych, a więc
zbiór co najwyżej przeliczalny.
Z definicji supremum wybieramy skończony $F_0$
z $\sum_{i\in F_0}a_i>\alpha-\varepsilon$.
Lemat \ref{ma2:lem:positive-family-partition} daje
$\sum_{i\in I\setminus F_0}a_i
=\alpha-\sum_{i\in F_0}a_i<\varepsilon$.
\end{proof}

\section{Rodziny w przestrzeniach unormowanych}

\begin{definicja}
\label{ma2:def:summable-family}
Niech $E$ będzie przestrzenią unormowaną nad $\K$,
$I$ dowolnym zbiorem, a $(x_i)_{i\in I}$ rodziną w $E$.
Dla skończonego $F\subseteq I$ piszemy
$S_F=\sum_{i\in F}x_i$. Rodzinę nazywamy sumowalną
do $S\in E$, gdy dla każdego $\varepsilon>0$ istnieje
skończony $F_0\subseteq I$ taki, że
$\|S_F-S\|<\varepsilon$ dla każdego skończonego
$F\supseteq F_0$. Piszemy wówczas $\sum_{i\in I}x_i=S$.
Jest to zbieżność sieci sum skończonych skierowanej
przez inkluzję: dla dwóch skończonych zbiorów ich suma
jest wspólnym większym elementem.
Rodzinę nazywamy bezwzględnie sumowalną, gdy
$\sum_{i\in I}\|x_i\|<\infty$.
\end{definicja}
Suma rodziny, jeśli istnieje, jest jednoznaczna.
Jeżeli $S,T$ spełniają definicję, to dla $\varepsilon>0$
bierzemy sumę dwóch zbiorów progowych z dokładnością
$\varepsilon/2$. Dla tego zbioru $F$ zachodzi
$\|S-T\|\le\|S-S_F\|+\|S_F-T\|<\varepsilon$, więc $S=T$.
Przypadek pustego $I$ daje sumę $0$.
Z definicji wynika liniowość sumowania: dla dwóch
sumowalnych rodzin na tym samym zbiorze indeksów
i skalarów $\alpha,\beta$ wybieramy wspólny próg
dla dokładności $\varepsilon/(1+|\alpha|+|\beta|)$.
Nierówność trójkąta dla skończonych sum daje
$\sum_i(\alpha x_i+\beta y_i)=\alpha\sum_i x_i+\beta\sum_i y_i$.

\begin{twierdzenie}[Kryterium Cauchy'ego dla rodzin]
\label{ma2:thm:family-cauchy}
Jeżeli $(x_i)_{i\in I}$ jest sumowalna, to
\begin{equation}
\label{ma2:eq:family-tail}
 \forall\varepsilon>0\ \exists F_0\subseteq I
 \text{ skończony}\quad
 \left\|\sum_{i\in H}x_i\right\|<\varepsilon
\end{equation}
dla każdego skończonego $H\subseteq I\setminus F_0$.
Jeżeli $E$ jest przestrzenią Banacha, warunek ten
jest również wystarczający.
\end{twierdzenie}
\begin{proof}
Dla rodziny sumowalnej do $S$ wybieramy próg $F_0$
dla dokładności $\varepsilon/2$. Dla $H$ jak w tezie
mamy $S_H=S_{F_0\cup H}-S_{F_0}$, zatem
$\|S_H\|<\varepsilon$.

Odwrotnie, dla każdego $n\ge1$ wybieramy skończony
$D_n\subseteq I$ spełniający warunek dla $2^{-n}$.
Niech $F_n=\bigcup_{j=1}^nD_j$. Zbiory $F_n$ są
skończone i rosnące. Dla $m\ge n$ zbiór
$F_m\setminus F_n$ jest rozłączny z $D_n$, więc
$\|S_{F_m}-S_{F_n}\|<2^{-n}$. Ciąg $(S_{F_n})$
jest ciągiem Cauchy'ego, zatem w przestrzeni Banacha
ma granicę $S$. Przechodząc do granicy po $m$ w
nierówności, dostajemy $\|S-S_{F_n}\|\le2^{-n}$.
Jeśli $F\supseteq F_n$ jest skończony, to również
$\|S_F-S_{F_n}\|<2^{-n}$, gdyż $F\setminus F_n$
nie przecina $D_n$. Stąd $\|S_F-S\|<2^{1-n}$.
Dla każdego $\varepsilon>0$ wybór $n$ z $2^{1-n}<\varepsilon$
daje zbieżność całej sieci, nie tylko wybranego ciągu.
\end{proof}

\begin{wniosek}
\label{ma2:cor:family-support}
Rodzina spełniająca warunek \eqref{ma2:eq:family-tail}
ma co najwyżej przeliczalny nośnik
$J=\{i\in I:x_i\ne0\}$.
Jeżeli rodzina jest sumowalna, to jej suma jest równa
sumie ograniczenia do $J$.
\end{wniosek}
\begin{proof}
Dla $m\ge1$ wybieramy $F_m$ dla $\varepsilon=1/m$.
Stosując warunek do singletonu poza $F_m$, widzimy,
że $\{i:\|x_i\|\ge1/m\}\subseteq F_m$. Zatem $J$
jest przeliczalną sumą zbiorów skończonych.
Jeżeli $F_0$ jest progiem sumowalności w $I$, to
$F_0\cap J$ jest progiem w $J$: dla skończonego
$G\subseteq J$ zawierającego $F_0\cap J$ mamy
$S_G=S_{G\cup F_0}$, gdyż dodane indeksy mają
wyrazy zerowe. Odwrotnie, próg w $J$ działa także
w $I$, ponieważ $S_F=S_{F\cap J}$.
\end{proof}

\begin{twierdzenie}[Związek z szeregami]
\label{ma2:thm:family-enumeration}
Niech $E$ będzie przestrzenią Banacha, a $J$
nieskończonym przeliczalnym nośnikiem rodziny.
Dla dowolnej bijekcji $\pi:\N_+\to J$ rodzina jest
sumowalna wtedy i tylko wtedy, gdy szereg
$\sum_{n\ge1}x_{\pi(n)}$ jest bezwarunkowo zbieżny.
Wtedy obie sumy są równe i nie zależą od $\pi$.
Dla skończonego nośnika suma jest zwykłą sumą skończoną.
\end{twierdzenie}
\begin{proof}
Z wniosku \ref{ma2:cor:family-support} wystarczy
rozpatrzyć rodzinę na $J$ (również gdy sumowalność
ograniczenia jest założeniem). Bijekcja $\pi$
przenosi skończone podzbiory $\N_+$ na wszystkie
skończone podzbiory $J$, zachowując inkluzję.
Warunek z definicji \ref{ma2:def:summable-family}
jest więc dokładnie warunkiem zbieżności sieci sum
skończonych szeregu. Równoważność i równość sum
wynikają z twierdzenia \ref{ma2:thm:unconditional-finite}.
Dla skończonego nośnika można obrać cały $J$ jako
zbiór progowy dla dowolnej dokładności.
\end{proof}

\begin{twierdzenie}[Sumowalność bezwzględna]
\label{ma2:thm:absolute-family}
W przestrzeni Banacha każda rodzina bezwzględnie
sumowalna jest sumowalna. Jeżeli $S$ jest jej sumą,
to dla każdego skończonego $F\subseteq I$
\[
 \|S-S_F\|\le\sum_{i\in I\setminus F}\|x_i\|.
\]
\end{twierdzenie}
\begin{proof}
Twierdzenie \ref{ma2:thm:positive-family-support}
zastosowane do $a_i=\|x_i\|$ daje dla każdego
$\varepsilon>0$ skończony $F_0$ z sumą norm poza
nim mniejszą od $\varepsilon$. Nierówność trójkąta
daje warunek \eqref{ma2:eq:family-tail}, więc
twierdzenie \ref{ma2:thm:family-cauchy} zapewnia
istnienie $S$. Dla dowolnego $\delta>0$ wybieramy
skończony $G$ zawierający $F$ i próg sumowalności
dla $\delta$. Wówczas
\[
 \|S-S_F\|\le\|S-S_G\|+\|S_{G\setminus F}\|
 <\delta+\sum_{i\in I\setminus F}\|x_i\|.
\]
Dowolność $\delta$ kończy dowód.
\end{proof}

\begin{uwaga}
Zbieżność bezwarunkowa nie musi oznaczać sumowalności
bezwzględnej w nieskończenie wymiarowej przestrzeni
Banacha: przykład \ref{ma2:ex:unconditional-l2}
daje rodzinę $(e_n/n)$ w $\ell^2$.
W skończonym wymiarze obie własności są równoważne,
na mocy wniosku \ref{ma2:cor:family-support},
twierdzenia \ref{ma2:thm:family-enumeration} i twierdzenia
\ref{ma2:thm:finite-unconditional-absolute}.
\end{uwaga}

\section{Sumowanie po rozłącznych częściach}

\begin{twierdzenie}[Dyskretne twierdzenie Tonellego]
\label{ma2:thm:family-tonelli}
Niech $I=\bigcup_{j\in J}I_j$, gdzie $(I_j)_{j\in J}$
jest rodziną parami rozłącznych zbiorów; dopuszczamy
zbiory puste. Dla $a_i\in[0,+\infty]$ zachodzi
\[
 \sum_{i\in I}a_i=\sum_{j\in J}\sum_{i\in I_j}a_i
\]
w $[0,+\infty]$.
\end{twierdzenie}
\begin{proof}
Oznaczmy lewą stronę przez $A$, a wewnętrzne sumy
przez $b_j$. Każdy skończony $F\subseteq I$
przecina tylko skończenie wiele części $I_j$;
oznaczmy ich indeksy przez $J_F$.
Z rozłączności mamy
$\sum_{i\in F}a_i=\sum_{j\in J_F}\sum_{i\in F\cap I_j}a_i
\le\sum_{j\in J}b_j$. Biorąc supremum po $F$,
otrzymujemy pierwszą nierówność.
Jeżeli $A=+\infty$, równość już wynika.
Niech więc $A<\infty$. Z monotoniczności
$b_j\le A$ dla każdego $j$. Dla dowolnego niepustego
skończonego $K\subseteq J$ i $\varepsilon>0$ wybieramy
skończone $F_j\subseteq I_j$ z
$\sum_{i\in F_j}a_i>b_j-\varepsilon/|K|$.
Suma zbiorów $F_j$ dla $j\in K$ jest skończona
i rozłączna, więc
$\sum_{j\in K}b_j<A+\varepsilon$.
Dowolność $\varepsilon$ daje $\sum_{j\in K}b_j\le A$;
dla pustego $K$ nierówność również zachodzi.
Biorąc supremum po $K$, otrzymujemy tezę.
\end{proof}

\begin{twierdzenie}[Grupowanie rodziny sumowalnej]
\label{ma2:thm:family-grouping}
Niech $E$ będzie przestrzenią Banacha, $(x_i)_{i\in I}$
rodziną sumowalną do $S$, a $(I_j)_{j\in J}$
podziałem $I$ na parami rozłączne części, z dopuszczeniem
części pustych. Każde ograniczenie do $I_j$ jest
sumowalne, a rodzina $y_j=\sum_{i\in I_j}x_i$
jest sumowalna i $\sum_{j\in J}y_j=S$.
Jeżeli rodzina wyjściowa jest bezwzględnie sumowalna,
to
\[
 \sum_{j\in J}\|y_j\|\le\sum_{i\in I}\|x_i\|.
\]
\end{twierdzenie}
\begin{proof}
Warunek \eqref{ma2:eq:family-tail} dla dowolnego
podzbioru $B\subseteq I$ wynika z warunku dla $I$:
wystarczy przeciąć zbiór progowy z $B$. Twierdzenie
\ref{ma2:thm:family-cauchy} daje sumowalność
każdego ograniczenia. Dla skończonego $K\subseteq J$
niech $B_K=\bigcup_{j\in K}I_j$.
Wykażemy $\sum_{i\in B_K}x_i=\sum_{j\in K}y_j$.
Dla $K=\varnothing$ jest to równość zer.
Dla $K\ne\varnothing$ i $\varepsilon>0$ wybieramy
w każdej części próg $D_j\subseteq I_j$ dla
dokładności $\varepsilon/|K|$. Dla dowolnego
skończonego $F\subseteq B_K$ zawierającego
$\bigcup_{j\in K}D_j$ zachodzi
\[
 \left\|\sum_{i\in F}x_i-\sum_{j\in K}y_j\right\|
 \le \sum_{j\in K}\left\|\sum_{i\in F\cap I_j}x_i-y_j\right\|
 <\varepsilon.
\]
Jednoznaczność sumy daje żądaną równość.

Ustalmy $\varepsilon>0$ i wybierzmy skończony $F_0$
z definicji sumowalności $S$ dla dokładności
$\varepsilon/2$. Niech $K_0$ będzie skończonym zbiorem
indeksów części przecinających $F_0$. Dla każdego
skończonego $K\supseteq K_0$ mamy $F_0\subseteq B_K$.
Niech $F\subseteq B_K$ będzie skończony, zawiera $F_0$
oraz próg sumowalności ograniczenia do $B_K$ dla
$\varepsilon/2$. Wtedy
\[
 \left\|\sum_{j\in K}y_j-S\right\|
 \le \left\|\sum_{i\in B_K}x_i-S_F\right\|+\|S_F-S\|
 <\varepsilon.
\]
To dowodzi sumowalności rodziny $(y_j)$ do $S$.
Jeżeli ponadto suma norm rodziny wyjściowej jest
skończona, z twierdzenia \ref{ma2:thm:absolute-family}
dla $F=\varnothing$ mamy
$\|y_j\|\le\sum_{i\in I_j}\|x_i\|$.
Monotoniczność i twierdzenie \ref{ma2:thm:family-tonelli}
dają ostatnią nierówność.
\end{proof}

\section{Szeregi podwójne i wielokrotne}

\begin{wniosek}[Dyskretne twierdzenia Tonellego i Fubiniego]
\label{ma2:cor:double-fubini}
Niech $I,J$ będą dowolnymi zbiorami.
Dla $a_{ij}\in[0,+\infty]$ zachodzi
\[
 \sum_{(i,j)\in I\times J}a_{ij}
 =\sum_{i\in I}\sum_{j\in J}a_{ij}
 =\sum_{j\in J}\sum_{i\in I}a_{ij}.
\]
Dla rodziny sumowalnej $(x_{ij})$ w przestrzeni
Banacha zachodzą te same równości w tej przestrzeni,
a każda z wewnętrznych i zewnętrznych rodzin
jest sumowalna. W szczególności wystarcza
$\sum_{(i,j)\in I\times J}\|x_{ij}\|<\infty$.
\end{wniosek}
\begin{proof}
Stosujemy twierdzenie \ref{ma2:thm:family-tonelli}
lub \ref{ma2:thm:family-grouping} do dwóch podziałów
$I\times J$: na części $\{i\}\times J$ oraz na części
$I\times\{j\}$. W przypadku sumy norm istnienie
sumy zapewnia twierdzenie \ref{ma2:thm:absolute-family}.
Każda z części jest w naturalnej bijekcji z odpowiednim
zbiorem indeksów, co nie zmienia skończonych sum.
\end{proof}

Dla $I=J=\N_+$ możemy używać klasycznej notacji
szeregów iterowanych. Jeżeli rodzina jest sumowalna
do $S$, jej sumy prostokątne
\[
 R_{mn}=\sum_{i=1}^m\sum_{j=1}^n x_{ij}
\]
dążą do $S$ w sensie Pringsheima: dla każdego
$\varepsilon>0$ istnieje $N$, że
$\|R_{mn}-S\|<\varepsilon$ dla $m,n\ge N$.
Istotnie, każdy skończony zbiór progowy jest zawarty
w pewnym kwadracie $\{1,\ldots,N\}^2$.
Zależności między tymi sposobami sumowania nie są
odwracalne bez dodatkowych założeń.

\begin{przyklad}[Różne sumy iterowane]
\label{ma2:ex:double-different}
Dla $i,j\ge1$ połóżmy $a_{ij}=1$ przy $j=i$,
$a_{ij}=-1$ przy $j=i+1$ i $a_{ij}=0$ w pozostałych
przypadkach. Każdy wiersz ma dwa niezerowe wyrazy,
których suma wynosi $0$. Pierwsza kolumna ma jeden
niezerowy wyraz $1$, a każda następna ma dwa
wyrazy $1,-1$. Zatem
\[
 \sum_{i=1}^{\infty}\sum_{j=1}^{\infty}a_{ij}=0,
 \qquad
 \sum_{j=1}^{\infty}\sum_{i=1}^{\infty}a_{ij}=1.
\]
Rodzina nie jest bezwzględnie sumowalna, ponieważ
suma modułów po $N$ punktach przekątnej wynosi $N$.
Nie jest też sumowalna z wniosku
\ref{ma2:cor:double-fubini}.
Sumy prostokątne są równe $1$ dla $n\le m$ i $0$
dla $n>m$: liczba dodatnich składników to
$\min(m,n)$, a ujemnych $\min(m,n-1)$.
Sumy kwadratowe są więc stałe, ale granica w sensie
Pringsheima nie istnieje.
\end{przyklad}

\begin{figure}[htbp]
\centering
\begin{tikzpicture}[x=.85cm,y=-.85cm]
\draw[step=1,gray!45] (0,0) grid (6,6);
\foreach \i in {1,...,6}{
 \node[above] at (\i-.5,0) {\small \i};
 \node[left] at (0,\i-.5) {\small \i};
 \node[teal!70!black] at (\i-.5,\i-.5) {$+1$};
}
\foreach \i in {1,...,5}{
 \node[red!70!black] at (\i+.5,\i-.5) {$-1$};
}
\node[above] at (6.65,0) {$j$};
\node[left] at (0,6.65) {$i$};
\draw[very thick,blue!60!black] (0,0) rectangle (4,4);
\draw[very thick,orange!85!black,dashed] (0,0) rectangle (5,4);
\end{tikzpicture}
\caption{Rodzina z przykładu \ref{ma2:ex:double-different}.
Niewypełnione pola oznaczają zera. Niebieski kwadrat
ma sumę $1$, a pomarańczowy prostokąt ma sumę $0$.}
\label{ma2:fig:double-different}
\end{figure}
\FloatBarrier

\begin{przyklad}[Granica prostokątna bez sum wierszy]
\label{ma2:ex:rectangular-not-iterated}
Niech $a_{1j}=(-1)^j$, $a_{2j}=-(-1)^j$,
a $a_{ij}=0$ dla $i\ge3$. Dla wszystkich $m\ge2$
i $n\ge1$ mamy $R_{mn}=0$, więc granica prostokątna
istnieje i jest równa $0$. Jednak sumy częściowe
pierwszego wiersza przyjmują naprzemiennie wartości
$-1,0$, a drugiego $1,0$. Żaden z tych szeregów
nie jest zbieżny. Zatem istnienie granicy prostokątnej
nie wystarcza nawet do określenia sumy iterowanej
po wierszach.
\end{przyklad}

\begin{wniosek}[Więcej niż dwa indeksy]
\label{ma2:cor:multiple-family}
Niech $r\ge1$ i $I_1,\ldots,I_r$ będą dowolnymi
zbiorami. Dla nieujemnej rodziny na $I_1\times\cdots\times I_r$
lub rodziny sumowalnej w przestrzeni Banacha
suma całej rodziny jest równa każdej z $r!$
sum iterowanych, z dowolną kolejnością indeksów.
W przypadku nieujemnym dopuszczamy wartość $+\infty$.
\end{wniosek}
\begin{proof}
Dla $r=1$ teza jest definicją. Załóżmy tezę dla $r-1$
i wybierzmy dowolny pierwszy indeks zewnętrzny.
Podzielmy iloczyn kartezjański na włókna dla
ustalonej wartości tego indeksu. Twierdzenie
\ref{ma2:thm:family-tonelli} lub
\ref{ma2:thm:family-grouping} daje sumowalność
odpowiednich rodzin wewnętrznych (w przypadku
nieujemnym są określone przez supremum) oraz
równość sumy z sumą po wybranym indeksie zewnętrznym.
Do każdego włókna stosujemy założenie indukcyjne
dla wybranej kolejności pozostałych $r-1$ indeksów.
Ponieważ wybór kolejności był dowolny, indukcja
kończy dowód.
\end{proof}

\section{Zastosowanie do iloczynów szeregów}

Jeśli $T:E\times F\to G$ jest ciągłym odwzorowaniem
dwuliniowym, $E,F,G$ są przestrzeniami Banacha
oraz $\sum_i\|x_i\|=\alpha<\infty$,
$\sum_j\|y_j\|=\beta<\infty$, to dla stałej $K$
z lematu \ref{ma2:lem:bilinear-bound} mamy
\[
 \sum_{(i,j)\in I\times J}\|T(x_i,y_j)\|
 \le K\alpha\beta.
\]
Istotnie, dla każdego skończonego zbioru par jego
rzuty na $I,J$ są skończone, a rozszerzenie sumy norm
na iloczyn tych rzutów daje oszacowanie przez
$K\alpha\beta$. Bierzemy supremum.
Twierdzenie \ref{ma2:thm:absolute-family} daje
sumowalność rodziny. Dla ustalonego $i$ ciągłe
odwzorowanie liniowe $v\mapsto T(x_i,v)$ zachowuje
granice sieci sum skończonych: z oszacowania
dwuliniowego norma różnicy jest nie większa niż
$K\|x_i\|$ razy norma różnicy odpowiedniej sumy
skończonej i jej granicy. Analogicznie dla drugiego
argumentu. Z wniosku \ref{ma2:cor:double-fubini}
otrzymujemy
\[
 \sum_{(i,j)\in I\times J}T(x_i,y_j)
 =T\left(\sum_{i\in I}x_i,\sum_{j\in J}y_j\right).
\]
Dla przeliczalnych rodzin jest to zgodne z wnioskiem
\ref{ma2:cor:bilinear-pair-enumeration} z rozdziału
o iloczynach Cauchy'ego, który nie wymagał zupełności
przestrzeni przy istniejących sumach.

W kolejnym rozdziale przejdziemy do algebr Banacha
i szeregu Neumanna.

% Rozwinięcie AF-15 i AF-31, bez teorii widma i rachunku holomorficznego.
\chapter{Algebry Banacha i szereg Neumanna}
\label{ma2:ch:banach-algebras}

Wykorzystamy szeregi w przestrzeniach Banacha oraz
twierdzenia o ciągłych odwzorowaniach dwuliniowych
z rozdziałów \ref{ma2:ch:banach-series}
i \ref{ma1:ch:cauchy-product}.
Podstawy teorii algebr Banacha pochodzą z analizy
funkcjonalnej (rozdziały 15 i 31). Nie potrzebujemy
tutaj teorii widma ani analizy zespolonej.

\section{Algebry unormowane}

\begin{definicja}[Algebra unormowana i algebra Banacha]
\label{ma2:def:normed-algebra}
\label{def:banach-algebra31}
Algebrą unormowaną nad $\K\in\{\R,\C\}$ nazywamy
przestrzeń unormowaną $A$ wyposażoną w łączne
mnożenie dwuliniowe spełniające
\[
 \|ab\|\le\|a\|\|b\|\qquad(a,b\in A).
\]
Jeżeli przestrzeń $A$ jest zupełna, nazywamy ją
algebrą Banacha. Jeżeli istnieje element $1_A$
taki, że $1_Aa=a1_A=a$ dla każdego $a\in A$,
nazywamy go jednostką. W dalszej części rozdziału
zakładamy istnienie jednostki i $\|1_A\|=1$,
chyba że zaznaczono inaczej. Nie zakładamy
przemienności mnożenia.
\end{definicja}
Jest to rozszerzenie definicji zespolonej algebry
Banacha z AF, rozdział 31, podrozdział o elementach
odwracalnych. Mnożenie jest ciągłe na mocy lematu
\ref{ma2:lem:bilinear-bound} ze stałą $K=1$.
Dla każdego $a\in A$ określamy $a^0=1_A$ oraz
$a^{n+1}=a^na$. Dla ustalonego $m\ge0$ równość $a^ma^0=a^m$
jest własnością jednostki. Jeżeli $a^ma^n=a^{m+n}$,
to z łączności
$a^ma^{n+1}=(a^ma^n)a=a^{m+n+1}$.
Zatem $a^ma^n=a^{m+n}$ dla wszystkich $m,n\ge0$.
Ponadto $\|a^0\|=1$, a krok
$\|a^{n+1}\|\le\|a^n\|\|a\|$ dowodzi indukcyjnie
$\|a^n\|\le\|a\|^n$. Stosujemy konwencję
$t^0=1$ również dla $t=0$.

\begin{przyklad}[Podstawowe algebry Banacha]
\label{ex:banach-algebras31}
Niech $E$ będzie niezerową przestrzenią Banacha.
Przestrzeń $\mathcal L(E,E)$ z normą operatorową
i składaniem operatorów jest algebrą Banacha.
Zupełność wynika ze stwierdzenia
\ref{prop:operator-norm12}. Łączność i dwuliniowość
wynikają z własności składania operatorów liniowych.
Dla każdego $x\in E$ mamy
$\|STx\|\le\|S\|\|T\|\|x\|$, skąd
$\|ST\|\le\|S\|\|T\|$. Jednostką jest operator
identycznościowy, którego norma wynosi $1$.

Jeżeli $K$ jest niepustą zwartą przestrzenią metryczną,
to $C(K,\K)$ z normą supremum i mnożeniem punktowym
jest przemienną algebrą Banacha. Zupełność wynika
z twierdzenia \ref{ma2:thm:continuous}, a nierówność
$\|fg\|_\infty\le\|f\|_\infty\|g\|_\infty$ z oszacowania
punktowego. Jednostką jest funkcja stale równa $1$.

Przestrzeń $\ell^1(\N_0,\K)$ ze splotem
\[
 (a*b)_n=\sum_{j=0}^n a_jb_{n-j}
\]
jest przemienną algebrą Banacha. Zupełność po
przeindeksowaniu wynika z twierdzenia
\ref{thm:seq-banach12}. Oszacowanie
$\|a*b\|_1\le\|a\|_1\|b\|_1$
jest dokładnie oszacowaniem z twierdzenia
\ref{ma1:thm:absolute-cauchy-product}.
Dwuliniowość wynika z rozdzielności sum skończonych,
przemienność z zamiany $j$ na $n-j$, a łączność z tego,
że dla każdego $n$ obie współrzędne
$((a*b)*c)_n$ i $(a*(b*c))_n$ są równe
$\sum_{i+j+k=n}a_i b_j c_k$ (indeksy nieujemne).
Wynika to z rozwinięcia skończonych sum
jak w dowodzie wniosku \ref{ma1:cor:finite-cauchy-products}.
Jednostką jest $(1,0,0,\ldots)$ o normie $1$.
\end{przyklad}
W pierwszym przypadku dla $E=\K^d$, $d\ge1$,
otrzymujemy algebrę macierzy z normą operatorową.

\section{Elementy odwracalne}

\begin{definicja}
\label{ma2:def:algebra-invertible}
Element $a\in A$ jest odwracalny, gdy istnieje $b\in A$
takie, że $ab=ba=1_A$. Zbiór takich elementów
oznaczamy przez $A^\times$.
\end{definicja}
Jeżeli $ba=1_A$ i $ac=1_A$, to
$b=b(ac)=(ba)c=c$. W szczególności obustronna
odwrotność jest jednoznaczna; piszemy $a^{-1}$.
Zbiór $A^\times$ jest grupą: jednostka jest odwracalna,
odwrotność elementu odwracalnego ma odwrotność $a$,
a z łączności wynika
$(ab)(b^{-1}a^{-1})=1_A=(b^{-1}a^{-1})(ab)$.
Zatem $(ab)^{-1}=b^{-1}a^{-1}$.

\begin{przyklad}[Odwrotność jednostronna nie wystarcza]
\label{ma2:ex:one-sided-inverse}
W $\ell^2(\N_+)$ określamy przesunięcia
\[
 R(x_1,x_2,\ldots)=(0,x_1,x_2,\ldots),\qquad
 L(x_1,x_2,\ldots)=(x_2,x_3,\ldots).
\]
Są liniowe, $\|Rx\|_2=\|x\|_2$ oraz
$\|Lx\|_2^2=\|x\|_2^2-|x_1|^2\le\|x\|_2^2$,
więc należą do algebry ograniczonych operatorów.
Z definicji $LR=I$, ale $RL(e_1)=0\ne e_1$.
Gdyby $R$ był odwracalny, mnożąc $LR=I$ z prawej
przez $R^{-1}$, otrzymalibyśmy $L=R^{-1}$.
Gdyby $L$ był odwracalny, mnożąc $LR=I$ z lewej
przez $L^{-1}$, otrzymalibyśmy $R=L^{-1}$.
W obu przypadkach wynikałoby $RL=I$, sprzeczność.
\end{przyklad}

\section{Szereg Neumanna}

\begin{twierdzenie}[Szereg Neumanna i oszacowanie reszty]
\label{ma2:thm:neumann}
\label{lem:neumann31}
\label{lem:neumann15}
Niech $A$ będzie algebrą Banacha z jednostką
o normie $1$. Jeśli $u\in A$ i $q=\|u\|<1$, to
\[
 (1_A-u)^{-1}=\sum_{n=0}^\infty u^n,\qquad
 \|(1_A-u)^{-1}\|\le\frac1{1-q}.
\]
Dla $N\ge0$ mamy ponadto
\[
 \left\|(1_A-u)^{-1}-\sum_{n=0}^N u^n\right\|
 \le\frac{q^{N+1}}{1-q}.
\]
\end{twierdzenie}
\begin{proof}
Twierdzenie w przypadku zespolonej algebry Banacha
i jego wersja operatorowa zostały udowodnione
w AF (odpowiednio rozdziały 31 i 15).
Podajemy krótki argument wspólny dla obu ciał,
z oszacowaniem reszty. Z $\|u^n\|\le q^n$ oraz
twierdzenia \ref{ma2:thm:absolute-series}
(po przeindeksowaniu) szereg ma sumę $v\in A$.
Dla $S_N=\sum_{n=0}^N u^n$ rozdzielność daje
$(1_A-u)S_N=S_N(1_A-u)=1_A-u^{N+1}$.
Z $\|u^{N+1}\|\le q^{N+1}\to0$ i ciągłości
mnożenia otrzymujemy $(1_A-u)v=v(1_A-u)=1_A$.
Z nierówności trójkąta $\|S_N\|\le\sum_{n=0}^N q^n$
i ciągłości normy otrzymujemy $\|v\|\le1/(1-q)$.
Oszacowanie reszty wynika z twierdzenia
\ref{ma2:thm:absolute-series}, ponieważ
$\sum_{n>N}q^n=q^{N+1}/(1-q)$ na mocy wzoru
na szereg geometryczny (AM I, rozdział 09).
Dla $q=0$ mamy $u=0$, a wszystkie reszty są zerowe.
\end{proof}

\begin{wniosek}[Słabsze kryterium potęgowe]
\label{ma2:cor:neumann-power}
Dla $u\in A$ równoważne są następujące warunki:
\begin{enumerate}[label=(\roman*)]
\item szereg $\sum_{n\ge0}u^n$ jest zbieżny;
\item szereg ten jest bezwzględnie zbieżny;
\item istnieje $m\ge1$ takie, że $\|u^m\|<1$.
\end{enumerate}
Jeżeli zachodzą, suma szeregu jest równa
$(1_A-u)^{-1}$.
\end{wniosek}
\begin{proof}
(i) implikuje (iii), ponieważ wyrazy zbieżnego
szeregu dążą do zera (twierdzenie
\ref{ma2:thm:series-cauchy}).
Niech zachodzi (iii), $q_m=\|u^m\|<1$.
Każdy $n\ge0$ ma jednoznaczny rozkład
$n=km+r$ z $k\ge0$ i $0\le r<m$.
Z podmultiplikatywności
$\|u^{km+r}\|\le q_m^k\|u^r\|$.
Dla każdej skończonej sumy częściowej sumę norm
ograniczamy przez
$(\sum_{r=0}^{m-1}\|u^r\|)\sum_{k\ge0}q_m^k$.
Jest to skończona liczba, więc zachodzi (ii).
(ii) implikuje (i) z twierdzenia
\ref{ma2:thm:absolute-series}.
Ponieważ $u^{N+1}\to0$, tożsamość teleskopowa
z dowodu twierdzenia \ref{ma2:thm:neumann}
daje obustronną odwrotność.
\end{proof}

\begin{przyklad}
\label{ma2:ex:neumann-nilpotent}
Dla macierzy $U=\begin{pmatrix}0&2\\0&0\end{pmatrix}$
mamy $U^2=0$ oraz $\|U\|=2$ w normie operatorowej
indukowanej przez normę euklidesową (dla
$(x,y)$ obraz ma normę $2|y|$).
Mimo $\|U\|>1$ szereg Neumanna kończy się po
dwóch składnikach i $(I-U)^{-1}=I+U$.
Natomiast dla $u=2$ w $\R$ element $1-u=-1$
jest odwracalny, ale szereg $\sum 2^n$ jest rozbieżny.
Sama odwracalność $1-u$ nie wystarcza do zbieżności
szeregu Neumanna.
\end{przyklad}

\section{Zaburzenia elementów odwracalnych}

\begin{twierdzenie}[Oszacowania odwrotności]
\label{ma2:thm:inverse-perturbation}
Niech $a\in A^\times$ i $h\in A$.
Jeżeli $t=\|a^{-1}h\|<1$, to $a+h\in A^\times$,
\[
 (a+h)^{-1}=\sum_{n=0}^\infty(-a^{-1}h)^n a^{-1},
 \qquad
 \|(a+h)^{-1}\|\le\frac{\|a^{-1}\|}{1-t}.
\]
Ponadto
\[
 \|(a+h)^{-1}-a^{-1}\|\le
 \frac{t\|a^{-1}\|}{1-t}.
\]
Reszta po wyrazie rzędu $N$ jest nie większa
niż $t^{N+1}\|a^{-1}\|/(1-t)$.
\end{twierdzenie}
\begin{proof}
Mamy $a+h=a(1_A+a^{-1}h)$.
Z twierdzenia \ref{ma2:thm:neumann} zastosowanego
do $u=-a^{-1}h$ drugi czynnik jest odwracalny.
Reguła odwrotności iloczynu daje rozwinięcie
z zachowaniem kolejności czynników. Ciągłość
mnożenia z prawej strony przez $a^{-1}$
pozwala pomnożyć granicę szeregu. Wszystkie
oszacowania wynikają z twierdzenia
\ref{ma2:thm:neumann} i nierówności
$\|va^{-1}\|\le\|v\|\|a^{-1}\|$;
w oszacowaniu różnicy bierzemy $N=0$.
\end{proof}
W szczególności, ponieważ $a^{-1}\ne0$
(gdyż $aa^{-1}=1_A$ i $\|1_A\|=1$), kula otwarta
o środku $a$ i promieniu $1/\|a^{-1}\|$
mieści się w $A^\times$,
gdyż $\|a^{-1}h\|\le\|a^{-1}\|\|h\|$.
Zatem $A^\times$ jest otwarty. Jeżeli
$\|h\|\le1/(2\|a^{-1}\|)$, to
$\|(a+h)^{-1}-a^{-1}\|\le2\|a^{-1}\|^2\|h\|$.
Otrzymujemy ciągłość odwracania, zgodnie
z wynikiem AF-31 (lemat o szeregu Neumanna).

\begin{wniosek}[Pierwsze przybliżenie odwrotności]
\label{ma2:cor:inverse-linear}
Przy $t=\|a^{-1}h\|<1$ zachodzi
\[
 \|(a+h)^{-1}-a^{-1}+a^{-1}ha^{-1}\|
 \le\frac{t^2\|a^{-1}\|}{1-t}.
\]
Jeśli $\|h\|\le 1/(2\|a^{-1}\|)$, prawa strona
jest nie większa niż $2\|a^{-1}\|^3\|h\|^2$.
\end{wniosek}
\begin{proof}
Bierzemy $N=1$ w poprzednim twierdzeniu.
Drugie oszacowanie wynika z $t\le1/2$ oraz
$t\le\|a^{-1}\|\|h\|$.
\end{proof}
Jest to przygotowanie do rachunku różniczkowego:
liniowa część przyrostu odwrotności to
$h\mapsto-a^{-1}ha^{-1}$. Pojęcie pochodnej
Frécheta i pełna teoria różniczkowania
odwracania znajdują się w AF, rozdział 15.

\section{Szeregi potęgowe w algebrze Banacha}

\begin{twierdzenie}[Podstawianie elementu do szeregu]
\label{ma2:thm:algebra-power-series}
Niech $(c_n)_{n\ge0}$ będzie ciągiem skalarów.
Jeżeli dla $r\ge0$ zachodzi $\sum_{n\ge0}|c_n|r^n<\infty$,
to dla każdego $a\in A$ z $\|a\|\le r$ istnieje
$f(a)=\sum_{n\ge0}c_na^n$ oraz
\[
 \left\|f(a)-\sum_{n=0}^N c_na^n\right\|
 \le\sum_{n>N}|c_n|r^n.
\]
Jeżeli $(d_n)_{n\ge0}$ jest drugim ciągiem skalarów
i $\sum_{n\ge0}|d_n|r^n<\infty$, to definiując
$g(a)=\sum_{n\ge0}d_na^n$, dla $\|a\|\le r$
otrzymujemy bezwzględnie zbieżny szereg
\[
 f(a)g(a)=\sum_{n=0}^\infty
     \left(\sum_{j=0}^n c_jd_{n-j}\right)a^n.
\]
\end{twierdzenie}
\begin{proof}
Z $\|c_na^n\|\le |c_n|r^n$ i twierdzenia
\ref{ma2:thm:absolute-series} otrzymujemy
istnienie sumy oraz oszacowanie reszty.
Do obu szeregów w $A$ stosujemy wniosek
\ref{ma2:cor:bilinear-absolute} dla mnożenia
algebry. Skończona suma współczynników
iloczynu Cauchy'ego ma podaną postać, ponieważ
$(c_ja^j)(d_{n-j}a^{n-j})=c_jd_{n-j}a^n$.
Nie potrzebujemy przemienności całej algebry.
\end{proof}

\begin{wniosek}[Funkcja wykładnicza]
\label{ma2:cor:algebra-exp}
Dla każdego $a\in A$ istnieje
$\exp a=\sum_{n\ge0}a^n/n!$ i $\|\exp a\|\le e^{\|a\|}$.
Jeżeli $a,b$ komutują, to
$\exp(a+b)=(\exp a)(\exp b)$.
W szczególności $\exp a$ jest odwracalny i
$(\exp a)^{-1}=\exp(-a)$.
\end{wniosek}
\begin{proof}
Szereg skalarny $\sum r^n/n!$ ma sumę $e^r$
z definicji funkcji wykładniczej w AM I,
rozdział 18. Twierdzenie
\ref{ma2:thm:algebra-power-series} zapewnia
istnienie sumy; oszacowanie normy wynika
z nierówności trójkąta i przejścia do granicy.
Stosujemy wniosek \ref{ma2:cor:bilinear-absolute}
do szeregów o wyrazach $a^n/n!$ i $b^n/n!$.
Przy $ab=ba$ wzór dwumianowy daje
\[
 \sum_{j=0}^n\frac{a^jb^{n-j}}{j!(n-j)!}
 =\frac{(a+b)^n}{n!}.
\]
Wzór dwumianowy w tym przypadku wynika przez
indukcję: mnożymy rozwinięcie dla $n$ przez
$a+b$, przesuwamy $a$ przez potęgi $b$ dzięki
przemienności pary i łączymy współczynniki
regułą Pascala. Rząd $0$ daje jednostkę.
Stąd prawo dodawania argumentów dla
komutujących elementów. Para $a,-a$ komutuje,
a $\exp0=1_A$, więc otrzymujemy obustronną odwrotność.
\end{proof}

\section{Przykład obliczeniowy}

\begin{przyklad}[Rozwiązanie równania operatorowego]
\label{ma2:ex:neumann-equation}
Niech $E$ będzie niezerową przestrzenią Banacha,
$T$ ograniczonym operatorem liniowym na $E$
i $0\le\|T\|\le q<1$.
Dla $b\in E$ równanie $x-Tx=b$ ma dokładnie jedno
rozwiązanie
$x=\sum_{n\ge0}T^nb$.
Istotnie, w algebrze operatorów $I-T$ jest odwracalny
z twierdzenia \ref{ma2:thm:neumann}.
Ewaluacja na $b$ jest ciągła, bo
$\|Sb\|\le\|S\|\|b\|$, więc rozwinięcie odwrotności
daje wskazaną sumę. Dla
$x_N=\sum_{n=0}^N T^nb$ mamy
\[
 \|x-x_N\|\le\frac{q^{N+1}}{1-q}\|b\|,
 \qquad b-(I-T)x_N=T^{N+1}b.
\]
Drugą równość daje tożsamość teleskopowa.
Dla dowolnego przybliżenia $y\in E$ i jego reszty
$r=b-(I-T)y$ zachodzi ponadto
\[
 x-y=(I-T)^{-1}r,\qquad
 \|x-y\|\le\frac{\|r\|}{1-q}.
\]
Pierwsza równość wynika z $(I-T)(x-y)=r$,
a druga z oszacowania normy odwrotności.
Jest to oszacowanie błędu na podstawie obliczonej
reszty, niezależne od sposobu uzyskania $y$.
W szczególności dla $q=1/2$ i $N=9$
błąd normy nie przekracza $2^{-9}\|b\|$.
Dla $E=\{0\}$ równanie ma jedyne rozwiązanie $0$
i wszystkie oszacowania są trywialnie spełnione.
\end{przyklad}

W kolejnym rozdziale zajmiemy się zbieżnością
punktową i jednostajną ciągów funkcji.

% Rozwinięcie TO-18 i AM II-01, przygotowanie do szeregów funkcyjnych.
\chapter{Zbieżność punktowa i jednostajna}
\label{ma2:ch:uniform-convergence}

Niech $X$ będzie zbiorem, a $(Y,d)$ przestrzenią
metryczną. Gdy $X$ jest przestrzenią topologiczną,
przez $C(X,Y)$ oznaczamy zbiór wszystkich ciągłych
odwzorowań $X\to Y$. Rozróżnimy sposoby zbieżności ciągów
odwzorowań $X\to Y$. W przypadku przestrzeni
unormowanych wykorzystamy wyniki rozdziału
\ref{ma2:ch:normed}. Topologiczne ujęcie przestrzeni
odwzorowań znajduje się w topologii ogólnej, rozdział 18.

\section{Definicje i pierwsze kryteria}

\begin{definicja}
\label{ma2:def:pointwise-uniform}
Ciąg $(f_n)_{n\ge1}$ odwzorowań $X\to Y$ jest
zbieżny punktowo do $f:X\to Y$, gdy
\[
 \forall x\in X\ \forall\varepsilon>0\ \exists N\ge1\
 \forall n\ge N:\quad d(f_n(x),f(x))<\varepsilon.
\]
Jest zbieżny jednostajnie do $f$ na $X$, gdy
\[
 \forall\varepsilon>0\ \exists N\ge1\
 \forall n\ge N\ \forall x\in X:\quad
 d(f_n(x),f(x))<\varepsilon.
\]
\end{definicja}
Różnica polega na tym, że w drugim przypadku wybór
$N$ jest niezależny od $x$. Zbieżność jednostajna
implikuje zbieżność punktową. Granica jest jednoznaczna
z jednoznaczności granicy ciągu w przestrzeni metrycznej
(AM I, rozdział 06, twierdzenie o jednoznaczności
granicy, etykieta \texttt{ma1:thm:unique-limit}).
Istotnie, dwie granice punktowe muszą mieć tę samą
wartość w każdym punkcie dziedziny. Na zbiorze pustym obie własności
zachodzą automatycznie; istnieje dokładnie jedno
odwzorowanie z tego zbioru do $Y$.

\begin{twierdzenie}[Kryteria zbieżności jednostajnej]
\label{ma2:thm:uniform-criteria}
Niech $X\ne\varnothing$, a $f_n,f:X\to Y$
będą odwzorowaniami. Następujące warunki są równoważne:
\begin{enumerate}[label=(\roman*)]
\item $f_n\to f$ jednostajnie na $X$;
\item $M_n:=\sup_{x\in X}d(f_n(x),f(x))\to0$,
gdzie dopuszczamy $M_n=+\infty$;
\item dla każdego ciągu $(x_n)$ w $X$ mamy
$d(f_n(x_n),f(x_n))\to0$.
\end{enumerate}
\end{twierdzenie}
\begin{proof}
Z (i), zastosowanego z dokładnością $\varepsilon/2$,
mamy $M_n\le\varepsilon/2<\varepsilon$ od pewnego
indeksu, co daje (ii). Z (ii) wynika (i), bo
$d(f_n(x),f(x))\le M_n$ dla każdego $x$.
Z (ii) otrzymujemy również (iii): dla dowolnego
ciągu $(x_n)$ od pewnego indeksu mamy
$0\le d(f_n(x_n),f(x_n))\le M_n$, a $M_n\to0$.
Gdyby (i) nie zachodziło, istniałoby $\varepsilon>0$
takie, że dla dowolnego $N$ można wybrać $n\ge N$
i $x\in X$ z $d(f_n(x),f(x))\ge\varepsilon$.
Wybieramy rekurencyjnie ściśle rosnące indeksy
$n_k$ i punkty $x_{n_k}$ spełniające tę nierówność:
najpierw stosujemy warunek z $N=1$, a po wyborze
$n_k$ stosujemy go z $N=n_k+1$.
Uzupełniamy pozostałe wyrazy ciągu $(x_n)$ dowolnym
ustalonym punktem $X$. Wzdłuż podciągu $(n_k)$
odległości są nie mniejsze niż $\varepsilon$,
co przeczy (iii).
\end{proof}

\begin{przyklad}
\label{ma2:ex:powers-pointwise}
Na $[0,1]$ funkcje $f_n(x)=x^n$ zbiegają punktowo
do funkcji $f$ równej $0$ na $[0,1)$ i $1$ w punkcie $1$.
Dla $x<1$ wynika to ze zbieżności potęg liczby
o module mniejszym od $1$ (AM I, rozdział 06, przykład ciągu geometrycznego,
etykieta \texttt{ma1:ex:geometric-sequence}).
Dla każdego $n$ mamy $\sup_{[0,1]}|f_n-f|=1$:
na $[0,1)$ wartości $x^n$ są mniejsze od $1$,
ale dążą do $1$ przy $x\to1^-$, a w punkcie $1$
różnica wynosi $0$. Zatem zbieżność nie jest
jednostajna. Na $[0,r]$, gdzie $0\le r<1$,
supremum różnicy wynosi $r^n\to0$, więc tam
zbieżność jest jednostajna. Rysunek
\ref{ma2:fig:powers-uniform} ilustruje oba przypadki.
\end{przyklad}

\begin{figure}[htbp]
\centering
\begin{tikzpicture}[x=8cm,y=3.8cm,>=Stealth]
 \draw[->] (-.035,0)--(1.08,0) node[right]{$x$};
 \draw[->] (0,-.06)--(0,1.14) node[above]{$y$};
 \draw[gray!55,densely dotted] (0,1)--(1,1)--(1,0);
 \foreach \n/\kol in {1/blue!75!black,2/teal!80!black,4/orange!90!black,12/red!70!black}{
  \draw[\kol,thick,domain=0:1,samples=81] plot (\x,{pow(\x,\n)});
 }
 \draw[black,very thick] (0,0)--(.992,0);
 \draw[black,fill=white] (1,0) circle[radius=1.8pt];
 \fill (1,1) circle[radius=1.8pt];
 \draw[gray,dashed] (.7,0)--(.7,.7);
 \node[below] at (.7,0) {$r$};
 \node[below] at (1,0) {$1$};
 \node[left] at (0,1) {$1$};
 \node[below left] at (0,0) {$0$};
 \node[blue!75!black,anchor=west] at (.05,.87) {$n=1$};
 \node[teal!80!black,anchor=west] at (.05,.73) {$n=2$};
 \node[orange!90!black,anchor=west] at (.05,.59) {$n=4$};
 \node[red!70!black,anchor=west] at (.05,.45) {$n=12$};
\end{tikzpicture}
\caption{Funkcje $x^n$ i ich granica punktowa (czarna).
Na każdym $[0,r]$, $r<1$, błąd nie przekracza $r^n$,
lecz na całym $[0,1]$ jego supremum wynosi $1$.}
\label{ma2:fig:powers-uniform}
\end{figure}
\FloatBarrier

\begin{przyklad}[Ciągła granica punktowa nie wystarcza]
\label{ma2:ex:moving-peak}
Na $[0,1]$ połóżmy
$g_n(x)=\max\{1-n|x-1/n|,0\}$.
Funkcje są ciągłe, $g_n(0)=0$, a dla ustalonego
$x>0$ mamy $g_n(x)=0$ od chwili, gdy $2/n<x$.
Zatem granicą punktową jest funkcja stale równa $0$.
Jednak $g_n(1/n)=1$ i $0\le g_n\le1$, więc
$\sup_{[0,1]}|g_n|=1$. Zbieżność nie jest jednostajna,
mimo ciągłości granicy i zwartości dziedziny.
\end{przyklad}

\section{Jednostajne kryterium Cauchy'ego}

\begin{twierdzenie}
\label{ma2:thm:uniform-cauchy}
Niech $(f_n)$ będzie ciągiem odwzorowań $X\to Y$.
Jeśli $f_n\to f$ jednostajnie, to
\[
 \forall\varepsilon>0\ \exists N\ge1\
 \forall m,n\ge N\ \forall x\in X:\quad
 d(f_m(x),f_n(x))<\varepsilon.
\]
Jeżeli $Y$ jest zupełna, warunek ten jest także
wystarczający do zbieżności jednostajnej do pewnej
funkcji $f:X\to Y$.
\end{twierdzenie}
\begin{proof}
Dla $\varepsilon>0$ wybieramy $N$ z definicji
zbieżności jednostajnej dla dokładności $\varepsilon/2$.
Dla $m,n\ge N$ i wszystkich $x\in X$ mamy
\[
 d(f_m(x),f_n(x))
 \le d(f_m(x),f(x))+d(f(x),f_n(x))<\varepsilon.
\]
To daje warunek konieczny.
Odwrotnie, dla każdego $x$ ciąg $(f_n(x))$ jest
ciągiem Cauchy'ego, więc z zupełności $Y$ ma
jednoznaczną granicę $f(x)$.
Dla $\varepsilon>0$ wybieramy wspólny $N$
w warunku Cauchy'ego z dokładnością $\varepsilon/2$.
Dla $n\ge N$ i każdego $x$, przechodząc przy ustalonych
$n,x$ do granicy $m\to\infty$, otrzymujemy
\[
 d(f_n(x),f(x))\le\varepsilon/2<\varepsilon.
\]
Przejście do granicy uzasadnia nierówność
$|d(y,z)-d(y,w)|\le d(z,w)$.
Indeks $N$ jest niezależny od $x$, co daje tezę.
\end{proof}
Jeżeli granica punktowa $f:X\to Y$ już istnieje,
to jednostajny warunek Cauchy'ego daje zbieżność
jednostajną do $f$ także bez zupełności $Y$:
w powyższym dowodzie zupełność służyła wyłącznie
do uzyskania granic punktowych. Wystarczy przejść
do granicy po $m$ z tym samym wspólnym progiem $N$.

Jeżeli $Y$ nie jest zupełna, wybierzmy w niej ciąg
Cauchy'ego $(y_n)$ bez granicy. Funkcje stałe
$f_n(x)=y_n$ na niepustym $X$ spełniają warunek
jednostajny Cauchy'ego, ale nie mają nawet granicy
punktowej. Zupełność jest zatem istotna.

\begin{twierdzenie}[Metryka zbieżności jednostajnej]
\label{ma2:thm:uniform-metric}
Na zbiorze $Y^X$ wszystkich odwzorowań $X\to Y$
wzór
\[
 D(f,g)=\sup_{x\in X}\min\{1,d(f(x),g(x))\}
\]
określa metrykę, przy konwencji $D=0$ dla $X=\varnothing$.
Zbieżność w tej metryce jest zbieżnością jednostajną.
Jeżeli $X\ne\varnothing$, metryka $D$ jest zupełna
wtedy i tylko wtedy, gdy $Y$ jest zupełna.
\end{twierdzenie}
\begin{proof}
Wartości $D$ należą do $[0,1]$. Symetria
i nieujemność wynikają z aksjomatów metryki $d$.
Jeśli $D(f,g)=0$, to dla każdego $x$ mamy
$\min\{1,d(f(x),g(x))\}=0$, czyli $f(x)=g(x)$.
Odwrotnie, $f=g$ daje $D(f,g)=0$.
Dla $a,b\ge0$ zachodzi
$\min\{1,a+b\}\le\min\{1,a\}+\min\{1,b\}$:
jeśli któryś składnik jest co najmniej $1$, prawa
strona jest co najmniej $1$, a w przeciwnym przypadku
wynosi $a+b$. Stosując tę nierówność do nierówności
trójkąta dla $d$ i biorąc supremum, otrzymujemy
nierówność trójkąta dla $D$.
Dla $0<\varepsilon<1$, jeśli $D(f_n,f)<\varepsilon$,
to $d(f_n(x),f(x))<\varepsilon$ dla wszystkich $x$.
Odwrotnie, jednostajna zbieżność z dokładnością
$\varepsilon/2$ daje $D(f_n,f)\le\varepsilon/2$.
Dla dowolnych $\varepsilon>0$ wystarczy w tych
argumentach użyć dokładności mniejszej od
$\min\{1,\varepsilon\}$, co obejmuje również
$\varepsilon\ge1$.
Te same oszacowania zastosowane do par $f_m,f_n$
pokazują równoważność warunku Cauchy'ego w metryce
$D$ i jednostajnego warunku Cauchy'ego.
Z twierdzenia \ref{ma2:thm:uniform-cauchy}
wynika zupełność, gdy $Y$ jest zupełna.
Odwrotnie, dla ciągu Cauchy'ego $(y_n)$ w $Y$
bierzemy funkcje stałe $f_n(x)=y_n$. Są ciągiem
Cauchy'ego w $D$, zatem przy zupełności $D$ mają
granicę $f$. Ewaluacja w dowolnym $x_0\in X$
daje $y_n\to f(x_0)$, czyli zupełność $Y$.
Dla pustego $X$ przestrzeń funkcji jest jednoelementowa.
\end{proof}

W przypadku przestrzeni unormowanej $E$ i funkcji
ograniczonych na niepustym $X$ można używać normy supremum:
\[
 \|f_n-f\|_\infty=\sup_{x\in X}\|f_n(x)-f(x)\|.
\]
Zbieżność w tej normie jest zbieżnością jednostajną
(z twierdzenia \ref{ma2:thm:uniform-criteria}).
Jednostajna granica funkcji ograniczonych jest
ograniczona: wybieramy $n$ z
$\sup_x\|f_n(x)-f(x)\|\le1$, skąd
$\|f(x)\|\le\|f_n\|_\infty+1$ dla każdego $x$.
Dla $X=\varnothing$ przestrzeń funkcji ograniczonych
jest przestrzenią zerową z normą $0$.
Zupełność $B(X,E)$ dla niepustego $X$ jest równoważna
zupełności $E$ na mocy twierdzenia \ref{ma2:thm:bounded}.
Nie powtarzamy tego dowodu.

\section{Ciągłość granicy jednostajnej}

\begin{twierdzenie}
\label{ma2:thm:uniform-continuous-limit}
Niech $X$ będzie przestrzenią topologiczną,
$(Y,d)$ przestrzenią metryczną i $f_n:X\to Y$
ciągłymi odwzorowaniami. Jeżeli $f_n\to f$
jednostajnie, to $f$ jest ciągła.
Jeśli $X$ jest metryczna i wszystkie $f_n$ są
jednostajnie ciągłe, to $f$ jest jednostajnie ciągła.
\end{twierdzenie}
\begin{proof}
Pierwsza część uogólnia argument z dowodu
twierdzenia \ref{ma2:thm:continuous} na dowolną
dziedzinę topologiczną i przeciwdziedzinę metryczną.
Ustalmy $x_0\in X$ i $\varepsilon>0$.
Wybieramy $n$ tak, aby dla każdego $x\in X$
$d(f_n(x),f(x))<\varepsilon/3$.
Ciągłość $f_n$ zapewnia otwarte otoczenie $U$
punktu $x_0$ z $d(f_n(x),f_n(x_0))<\varepsilon/3$
dla $x\in U$. Dla tych $x$ mamy
\[
 d(f(x),f(x_0))
 \le d(f(x),f_n(x))+d(f_n(x),f_n(x_0))
       +d(f_n(x_0),f(x_0))<\varepsilon.
\]
To dowodzi ciągłości.
W przypadku jednostajnej ciągłości wybieramy
$n$ tak samo, po czym bierzemy $\delta>0$
z jednostajnej ciągłości $f_n$ dla
$\varepsilon/3$. Dla dowolnych $x,y\in X$
o odległości mniejszej od $\delta$ to samo
oszacowanie, z $y$ zamiast $x_0$, daje
$d(f(x),f(y))<\varepsilon$.
Wybór $\delta$ jest niezależny od $x,y$.
\end{proof}
Przykład \ref{ma2:ex:powers-pointwise} pokazuje,
że granica punktowa funkcji ciągłych może
być nieciągła.

\begin{wniosek}
\label{ma2:cor:uniform-complete-continuous}
Jeśli $X$ jest przestrzenią topologiczną, a $Y$
zupełną przestrzenią metryczną, to $C(X,Y)$
jest zupełna w metryce $D$ z twierdzenia
\ref{ma2:thm:uniform-metric}.
\end{wniosek}
\begin{proof}
Ciąg Cauchy'ego w tej metryce ma granicę
jednostajną w $Y^X$ z twierdzenia
\ref{ma2:thm:uniform-metric}. Granica jest ciągła
z twierdzenia \ref{ma2:thm:uniform-continuous-limit}.
\end{proof}

\section{Zbieżność jednostajna na zbiorach zwartych}

\begin{definicja}
\label{ma2:def:compact-uniform}
Niech $X$ będzie przestrzenią topologiczną.
Ciąg $(f_n)$ odwzorowań $X\to Y$ jest zbieżny
jednostajnie na zbiorach zwartych do $f$,
gdy dla każdego zwartego $K\subseteq X$
ograniczenia $f_n|_K$ zbiegają jednostajnie do $f|_K$.
\end{definicja}
Zbieżność jednostajna na całym $X$ implikuje
zbieżność jednostajną na zbiorach zwartych,
a ta implikuje zbieżność punktową (zbiór
jednoelementowy jest zwarty). Jeśli $X$ jest
zwarta, pierwsze dwie własności są równoważne.
Jeśli $X$ jest skończonym zbiorem, zbieżność
punktowa implikuje zbieżność jednostajną:
dla każdego punktu wybieramy próg dla ustalonej
dokładności i bierzemy maksimum skończonej
liczby progów; przypadek pusty jest automatyczny.

\begin{przyklad}[Zbieżność na zbiorach zwartych nie jest globalnie jednostajna]
\label{ex:topog-odwzorowania-gorki-uciekajace}
Przytaczamy przykład z topologii ogólnej, rozdział 18,
sekcja 2. Na $\R$ niech
$f_n(x)=\max\{1-|x-n|,0\}$.
Każda funkcja jest ciągła, $\sup_{\R}|f_n|=1$.
Dla dowolnego zwartego $K\subseteq\R$ istnieje $M$
z $K\subseteq[-M,M]$. Gdy $n>M+1$, funkcja
$f_n$ znika na $K$. Stąd zbieżność jednostajna
na zbiorach zwartych do zera, lecz nie na całej $\R$.
\end{przyklad}
Przykład \ref{ma2:ex:moving-peak} pokazuje natomiast,
że zbieżność punktowa nie implikuje zbieżności
jednostajnej na zbiorach zwartych, nawet przy
ciągłej granicy.

\begin{twierdzenie}
\label{ma2:thm:compact-uniform-continuity}
Niech $X,Y$ będą przestrzeniami metrycznymi.
Jeżeli $(f_n)$ jest ciągiem ciągłych odwzorowań
$X\to Y$ zbieżnym do $f$ jednostajnie na zbiorach
zwartych, to $f$ jest ciągła.
\end{twierdzenie}
\begin{proof}
Niech $x_j\to x$ w $X$. Zbiór
$K=\{x\}\cup\{x_j:j\ge1\}$ jest zwarty:
w dowolnym pokryciu otwartym element zawierający
$x$ zawiera wszystkie $x_j$ od pewnego indeksu,
a pozostałe punkty można pokryć skończoną liczbą
elementów pokrycia. Na $K$ mamy zbieżność
jednostajną funkcji ciągłych, więc
twierdzenie \ref{ma2:thm:uniform-continuous-limit}
daje ciągłość $f|_K$. W szczególności
$f(x_j)\to f(x)$. Ponieważ ciąg był dowolny,
kryterium ciągowe ciągłości (AM I, rozdział 15,
twierdzenie o kryterium ciągowym, etykieta
\texttt{ma1:thm:continuity-sequential}) daje ciągłość $f$.
\end{proof}
Na przestrzeniach lokalnie zwartych Hausdorffa
można również wykorzystać zwarte otoczenia
(topologia ogólna, rozdział 14). Dla każdego $x$
wybieramy zwarty zbiór $K$ zawierający otwarte
otoczenie $U$ punktu $x$. Granica $f|_K$ jest
ciągła z twierdzenia \ref{ma2:thm:uniform-continuous-limit},
a zatem $f|_U$ jest ciągła. Ponieważ $U$ jest
otwarty w $X$, daje to ciągłość $f$ w $x$.

Przypomnijmy, że topologia zwarto-otwarta na $C(X,Y)$
ma podbazę zbiorów
\[
 [K,V]=\{g\in C(X,Y):g(K)\subseteq V\},
\]
gdzie $K\subseteq X$ jest zwarty, a $V\subseteq Y$
otwarty. Topologia zbieżności jednostajnej na
zbiorach zwartych ma w punkcie $f$ bazę otoczeń
\[
 B_K(f,\varepsilon)=
 \{g\in C(X,Y):d(f(x),g(x))<\varepsilon
                    \text{ dla wszystkich }x\in K\}.
\]
Są to definicje z topologii ogólnej, rozdział 18,
sekcje 1 i 2. Dla niepustego $K$ odległość
$d(f(x),g(x))$ osiąga maksimum na $K$, ponieważ
jest ciągła; dlatego warunek punktowy z ustalonym
$\varepsilon$ jest tu równoważny ostrej nierówności
dla supremum. Dla $K=\varnothing$ otrzymujemy całe
$C(X,Y)$.

\begin{twierdzenie}[Równoważność dwóch topologii]
\label{thm:topog-odwzorowania-zwarto-otwarta-jednostajna}
Na $C(X,Y)$ dla dowolnej przestrzeni topologicznej
$X$ i metrycznej $Y$ topologia zwarto-otwarta
jest równa topologii zbieżności jednostajnej
na zbiorach zwartych. W szczególności ciąg
funkcji ciągłych zbiega w tej topologii do
$f\in C(X,Y)$ dokładnie wtedy, gdy zbiega do
$f$ jednostajnie na zbiorach zwartych.
\end{twierdzenie}
\emph{Dowód:} topologia ogólna, rozdział 18, sekcja 2,
twierdzenie o równoważności dwóch topologii
(z tą samą etykietą źródłową) i następujący po nim
wniosek o zbieżności sieci.
Dla lokalnie zwartej, $\sigma$-zwartej przestrzeni
Hausdorffa metryzacja i zupełność tej topologii
przy zupełnym $Y$ zostały udowodnione w sekcji 4
tego samego rozdziału. Nie powtarzamy tych dowodów.

\section{Twierdzenie Diniego}

\begin{twierdzenie}[Twierdzenie Diniego]
\label{ma2:thm:dini}
Niech $X$ będzie zwartą przestrzenią topologiczną,
a $(f_n)$ ciągiem ciągłych funkcji rzeczywistych
na $X$ zbieżnym punktowo do ciągłej funkcji $f$.
Jeżeli dla każdego $x\in X$ zachodzi
$f_n(x)\le f_{n+1}(x)$ dla wszystkich $n$,
albo dla każdego $x\in X$ zachodzi
$f_n(x)\ge f_{n+1}(x)$ dla wszystkich $n$,
to zbieżność jest jednostajna.
\end{twierdzenie}
\begin{proof}
W przypadku rosnącym mamy $f_n(x)\le f(x)$
dla każdego $n,x$, ponieważ ciąg rzeczywisty
jest rosnący i zbieżny do $f(x)$.
Dla $\varepsilon>0$ definiujemy otwarte zbiory
\[
 U_n=\{x\in X:f(x)-f_n(x)<\varepsilon\}.
\]
Są otwarte dzięki ciągłości różnicy,
rosnące ze względu na inkluzję i pokrywają $X$
z punktowej zbieżności. Jeżeli $X\ne\varnothing$,
zwartość daje skończone podpokrycie
$U_{n_1},\ldots,U_{n_s}$. Biorąc $N=\max_j n_j$,
otrzymujemy $X=U_N$ dzięki rosnącemu układowi.
Dla $n\ge N$ i wszystkich $x$ zachodzi więc
$0\le f(x)-f_n(x)<\varepsilon$.
Jest to zbieżność jednostajna.
W przypadku malejącym stosujemy ten sam dowód
do $-f_n$ i $-f$; przypadek pustej dziedziny
wynika z definicji.
\end{proof}

\begin{uwaga}
W istocie argument Diniego działa także dla
odwzorowań ciągłych $f_n,f:X\to(Y,d)$, gdzie $X$
jest zwarta, jeśli $f_n\to f$ punktowo i dla każdego
$x$ ciąg $d(f_n(x),f(x))$ jest nierosnący.
Zbiory $\{x:d(f_n(x),f(x))<\varepsilon\}$ są otwarte,
gdyż nierówność
\[
 |d(y,z)-d(y',z')|\le d(y,y')+d(z,z')
\]
wynikająca z dwukrotnego zastosowania nierówności
trójkąta daje ciągłość odległości względem obu
argumentów. Zbiory są rosnące dzięki założeniu
o nierosnących odległościach.
Pokrywają $X$ z punktowej zbieżności, więc argument
skończonego podpokrycia z dowodu twierdzenia
\ref{ma2:thm:dini} daje zbieżność jednostajną.
Dla funkcji rzeczywistych obejmuje to również
przypadek, w którym kierunek monotoniczności
ciągu $f_n(x)$ zależy od $x$.
\end{uwaga}

Założenia twierdzenia są istotne. Na niezwartym
$[0,1)$ funkcje $x^n$ maleją punktowo do ciągłego
zera, ale supremum $x^n$ nadal wynosi $1$.
Na $[0,1]$ ten sam ciąg ma nieciągłą granicę
(przykład \ref{ma2:ex:powers-pointwise}).
Bez monotoniczności nawet na zwartym przedziale
ciągłość granicy nie wystarcza, jak pokazuje
przykład \ref{ma2:ex:moving-peak}.

W następnym rozdziale zastosujemy kryteria
zbieżności jednostajnej do szeregów funkcyjnych.

% Rozwinięcie AM II-02, AM II-03 i AM II-08.
\chapter{Szeregi funkcyjne}
\label{ma2:ch:function-series}

Niech $X$ będzie niepustym zbiorem, a $E$ przestrzenią
unormowaną nad $\K$. Wykorzystamy teorię szeregów
wektorowych z rozdziału \ref{ma2:ch:banach-series}
oraz zbieżności jednostajnej z rozdziału
\ref{ma2:ch:uniform-convergence}. Zupełność $E$
będziemy zakładać tam, gdzie jest potrzebna
do istnienia sumy szeregu.

\section{Sumy częściowe i kryterium Cauchy'ego}

\begin{definicja}
\label{ma2:def:function-series}
Dla ciągu funkcji $u_n:X\to E$, $n\ge1$, określamy
sumy częściowe $S_0(x)=0$ oraz
$S_N(x)=\sum_{n=1}^N u_n(x)$.
Szereg funkcyjny $\sum_{n\ge1}u_n$ nazywamy
zbieżnym punktowo (odpowiednio jednostajnie)
do $S:X\to E$, jeżeli ciąg $(S_N)$ jest zbieżny
punktowo (odpowiednio jednostajnie) do $S$.
Szereg nazywamy punktowo bezwzględnie zbieżnym,
jeżeli dla każdego $x\in X$ szereg liczbowy
$\sum_{n\ge1}\|u_n(x)\|$ jest zbieżny.
\end{definicja}
W przypadku pustej dziedziny wszystkie warunki
zbieżności zachodzą automatycznie.
Zbieżność punktowa oznacza dokładnie zbieżność
każdego szeregu wektorowego $\sum_n u_n(x)$.
Jednoznaczność sumy wynika z jednoznaczności
granicy punktowej. W dalszym ciągu szeregi
indeksowane od $0$ rozumiemy analogicznie.

\begin{twierdzenie}[Jednostajne kryterium Cauchy'ego]
\label{ma2:thm:function-series-cauchy}
Niech $u_n:X\to E$ dla $n\ge1$.
Jeśli szereg $\sum_n u_n$ jest jednostajnie zbieżny,
to
\[
 \forall\varepsilon>0\ \exists N\ge1\
 \forall q\ge p\ge N\ \forall x\in X:
 \quad \left\|\sum_{n=p}^q u_n(x)\right\|<\varepsilon.
\]
Jeżeli $E$ jest przestrzenią Banacha, warunek jest
również wystarczający. Jest wystarczający także
w niezupełnej przestrzeni $E$, jeśli szereg ma już
sumę punktową w $E$.
\end{twierdzenie}
\begin{proof}
Dla $q\ge p\ge1$ zachodzi
$\sum_{n=p}^q u_n=S_q-S_{p-1}$.
Jeżeli $(S_n)$ spełnia jednostajny warunek
Cauchy'ego z progiem $M$, to wybieramy $N=M+1$.
Odwrotnie, dla $m>n\ge N$ stosujemy warunek
z tezy z $p=n+1$, $q=m$. Dla $n>m$ korzystamy
z symetrii normy różnicy, a dla $n=m$ różnica
wynosi zero. Otrzymujemy jednostajny warunek
Cauchy'ego dla $(S_n)$.
Wszystkie wnioski wynikają teraz z twierdzenia
\ref{ma2:thm:uniform-cauchy} oraz uwagi o istniejącej
granicy punktowej po jego dowodzie.
\end{proof}

\begin{wniosek}
\label{ma2:cor:uniform-series-term}
Jednostajna zbieżność szeregu implikuje
jednostajną zbieżność jego wyrazów do zera.
\end{wniosek}
\begin{proof}
W kryterium \ref{ma2:thm:function-series-cauchy}
bierzemy $q=p=n$. Otrzymujemy warunek
jednostajnej zbieżności $(u_n)$ do zera.
\end{proof}
Warunek nie jest wystarczający: na dowolnym
niepustym $X$ funkcje $u_n(x)=1/n$ zbiegają
jednostajnie do zera, ale ich szereg jest
rozbieżny (AM I, rozdział 09, przykład szeregu
harmonicznego, etykieta \texttt{ma1:ex:harmonic-series}).

Jeżeli $X$ jest przestrzenią topologiczną,
wszystkie $u_n$ są ciągłe i szereg zbiega
jednostajnie, to suma $S$ jest ciągła:
każda suma częściowa jest ciągła, a twierdzenie
\ref{ma2:thm:uniform-continuous-limit} daje tezę.
Jeżeli $X$ jest metryczna, wystarcza również
zbieżność jednostajna na zbiorach zwartych
(twierdzenie \ref{ma2:thm:compact-uniform-continuity}).

\section{Kryterium Weierstrassa i zbieżność normalna}

\begin{definicja}
\label{ma2:def:normal-series}
Szereg $\sum_n u_n$ jest zbieżny normalnie na $X$,
jeśli dla każdego $n$ istnieje skończona liczba
$M_n=\sup_{x\in X}\|u_n(x)\|$ i
\[
 \sum_{n=1}^{\infty}M_n<\infty.
\]
Dla przestrzeni topologicznej $X$ mówimy o zbieżności
normalnej na zbiorach zwartych, jeśli warunek
zachodzi po ograniczeniu do każdego zwartego
$K\subseteq X$. Dla $K=\varnothing$ przyjmujemy
wszystkie normy supremum równe zero.
\end{definicja}

\begin{twierdzenie}[Kryterium Weierstrassa]
\label{ma2:thm:weierstrass-test}
Niech $E$ będzie przestrzenią Banacha. Jeśli
istnieje zbieżny szereg nieujemnych liczb
$\sum_n a_n$ taki, że
$\|u_n(x)\|\le a_n$ dla wszystkich $n,x$,
to szereg $\sum_n u_n$ jest zbieżny normalnie
i jednostajnie, a szereg funkcji
$\sum_n\|u_n(\cdot)\|$ jest jednostajnie zbieżny.
Dla sumy $S$ i każdego $N\ge0$ zachodzi
\[
 \sup_{x\in X}\|S(x)-S_N(x)\|
 \le\sum_{n>N}a_n.
\]
\end{twierdzenie}
\begin{proof}
Z definicji supremum mamy $0\le M_n\le a_n$.
Z porównania szeregów nieujemnych (AM I,
rozdział 09, etykieta \texttt{ma1:thm:series-comparison})
wynika $\sum_n M_n\le\sum_n a_n<\infty$.
Każde $u_n$ jest więc elementem przestrzeni $B(X,E)$,
która jest przestrzenią Banacha na mocy
twierdzenia \ref{ma2:thm:bounded}. Zastosowanie
twierdzenia \ref{ma2:thm:absolute-series} w tej
przestrzeni daje zbieżność w normie supremum
i oszacowanie reszty przez $\sum_{n>N}M_n$.
Z twierdzenia \ref{ma2:thm:uniform-criteria}
jest to zbieżność jednostajna.
Analogicznie, w przestrzeni Banacha $B(X,\R)$,
funkcje $v_n(x)=\|u_n(x)\|$ mają normy $M_n$,
więc ich szereg jest jednostajnie zbieżny.
\end{proof}
Nie jest konieczne, aby oszacowanie
$\|u_n(x)\|\le a_n$ zachodziło dla skończonego
początku szeregu: dla zbieżności jednostajnej
wystarczy oszacowanie ogona. Wtedy stosujemy
twierdzenie do ogona, a następnie dodajemy
skończoną sumę funkcji początkowych, która
nie zależy od indeksu sumy częściowej.
Jeśli choć jedna z tych początkowych funkcji
jest nieograniczona, cały szereg nie jest zbieżny
normalnie w sensie przyjętej definicji.
Na przykład dla $X=\R$, $E=\R$, $u_1(x)=x$
i $u_n(x)=0$ dla $n\ge2$, sumy częściowe
od indeksu $1$ są identyczne, ale
$\sup_{x\in\R}|u_1(x)|=+\infty$.
Dla niezupełnej przestrzeni $E$ oszacowanie
Weierstrassa wciąż daje jednostajny warunek
Cauchy'ego i zbieżność szeregu norm. Jeżeli
istnieje suma punktowa w $E$, daje także
zbieżność jednostajną i to samo oszacowanie reszty:
stosujemy nierówność trójkąta do skończonych
ogonów i przechodzimy do granicy przy ustalonym $x$,
po czym bierzemy supremum.

\begin{przyklad}[Szereg geometryczny]
\label{ma2:ex:function-geometric}
Dla $0\le r<1$ szereg $\sum_{n=0}^\infty z^n$
jest zbieżny normalnie na $\{z\in\C:|z|\le r\}$,
ponieważ $\sup_{|z|\le r}|z^n|=r^n$ (przy $r^0=1$).
Z twierdzenia o szeregu geometrycznym
(AM I, rozdział 09, etykieta \texttt{ma1:thm:geometric-series})
suma jest równa $1/(1-z)$. Błąd po wyrazie
o indeksie $N$ jest nie większy niż
$r^{N+1}/(1-r)$.
Każdy niepusty zwarty podzbiór otwartego koła
jednostkowego mieści się w kole o promieniu $r<1$,
bo ciągła funkcja $z\mapsto |z|$ osiąga na nim maksimum.
Stąd zbieżność normalna na zbiorach zwartych.
Na całym otwartym kole zbieżność nie jest jednostajna:
dla każdego $n\ge1$ supremum $|z|^n$ wynosi $1$,
co przeczy wnioskowi \ref{ma2:cor:uniform-series-term}.
\end{przyklad}

\section{Różne rodzaje zbieżności szeregów}

Z kryterium Weierstrassa w przestrzeni Banacha
wynika ciąg implikacji:
zbieżność normalna implikuje jednostajną
zbieżność szeregu norm, a ta implikuje
zbieżność jednostajną szeregu wektorowego.
Istotnie, jednostajny warunek Cauchy'ego dla
szeregu norm i nierówność
$\|\sum_{n=p}^q u_n(x)\|\le\sum_{n=p}^q\|u_n(x)\|$
dają warunek z twierdzenia
\ref{ma2:thm:function-series-cauchy}.
Zbieżność szeregu norm jednostajnie implikuje
też jego zbieżność punktową.
Poniższe przykłady pokazują, że własności te
nie są równoważne. W szczególności zbieżność
jednostajna i punktowa zbieżność bezwzględna
nie implikują jedna drugiej. Dla szeregu o wartościach
w przestrzeni niezupełnej zbieżność normalna
zapewnia jednostajny warunek Cauchy'ego, lecz
istnienie sumy wymaga dodatkowego założenia
(por. przykład \ref{ma2:ex:incomplete-series}
przeniesiony na funkcje stałe).

\begin{przyklad}[Jednostajna zbieżność bez zbieżności normalnej]
\label{ma2:ex:uniform-nonnormal}
Na $[0,1]$ niech $u_n(x)=(-1)^{n-1}/n$.
Szereg jest jednostajnie zbieżny, ponieważ
każda suma częściowa jest funkcją stałą,
a szereg liczbowy jest zbieżny z twierdzenia
\ref{ma1:thm:leibniz-series}.
Nie jest punktowo bezwzględnie zbieżny,
gdyż szereg modułów to szereg harmoniczny.
\end{przyklad}

\begin{przyklad}[Jednostajna zbieżność szeregu norm bez zbieżności normalnej]
\label{ma2:ex:absolute-uniform-nonnormal}
Dla $n\ge1$ niech
$I_n=[2^{-n},2^{1-n}]$,
$c_n=3\cdot2^{-n-1}$ i $h_n=2^{-n-1}$.
Na $[0,1]$ określamy
\[
 u_n(x)=\frac1n\max\{1-|x-c_n|/h_n,0\}.
\]
Funkcja $u_n$ jest ciągła i nieujemna, znika poza
wnętrzem $I_n$, a $u_n(c_n)=1/n$.
Wnętrza przedziałów $I_n$ są parami rozłączne.
W każdym punkcie co najwyżej jeden wyraz jest
niezerowy, więc od pewnego indeksu ciąg sum
częściowych w tym punkcie jest stały i suma
punktowa $S$ istnieje.
Dla $N\ge0$ mamy
\[
 \sup_{x\in[0,1]}|S(x)-S_N(x)|=\frac1{N+1}.
\]
Nierówność górna wynika z tego, że jedyny możliwy
niezerowy wyraz ogona ma indeks $n>N$,
a równość otrzymujemy dla $x=c_{N+1}$.
Zatem szereg norm (tu równy szeregowi funkcji)
jest jednostajnie zbieżny, lecz
$\sum_n\|u_n\|_\infty=\sum_n1/n=+\infty$.
\end{przyklad}

\begin{figure}[htbp]
\centering
\begin{tikzpicture}[x=11cm,y=3cm,>=Stealth]
 \draw[->] (0,0)--(1.04,0) node[right]{$x$};
 \draw[->] (0,0)--(0,1.13) node[above]{$y$};
 \foreach \n in {1,...,5}{
  \pgfmathsetmacro{\aa}{pow(2,-\n)}
  \pgfmathsetmacro{\bb}{2*\aa}
  \pgfmathsetmacro{\cc}{1.5*\aa}
  \pgfmathsetmacro{\hh}{1/\n}
  \draw[blue!70!black,thick,fill=blue!10]
       (\aa,0)--(\cc,\hh)--(\bb,0)--cycle;
  \node[above,font=\scriptsize] at (\cc,\hh) {$1/\n$};
 }
 \foreach \xx/\lab in {.25/{1/4},.5/{1/2},1/1}{
  \draw (\xx,0)--(\xx,-.025) node[below] {$\lab$};
 }
 \node[below left] at (0,0) {$0$};
\end{tikzpicture}
\caption{Rozłączne wnętrza przedziałów z przykładu
\ref{ma2:ex:absolute-uniform-nonnormal}.
Supremum ogona wynosi $1/(N+1)$, chociaż suma
supremów wyrazów jest nieskończona. Pokazano pięć
pierwszych funkcji.}
\label{ma2:fig:normal-hats}
\end{figure}
\FloatBarrier

\begin{przyklad}[Punktowa zbieżność bezwzględna bez jednostajnej]
\label{ma2:ex:pointwise-absolute-not-uniform}
Na $[0,1]$ połóżmy $u_n(x)=(1-x)x^n$ dla $n\ge1$.
Wyrazy są nieujemne i
$S_N(x)=x-x^{N+1}$ na mocy teleskopowania.
Dla $0\le x<1$ granica wynosi $x$, dla $x=1$
wynosi $0$. Zatem szereg jest punktowo bezwzględnie
zbieżny, ale $\sup_{[0,1]}|S-S_N|=1$:
dla $x<1$ różnica wynosi $x^{N+1}$, a w punkcie
$1$ wynosi zero. Zbieżność nie jest jednostajna.
\end{przyklad}

\paragraph{Rysunek 3D: sumy częściowe szeregu.}
Krzywe przedstawiają sumy częściowe $S_N(x)=x-x^{N+1}$.
Dla każdego $N$ wartość w $x=1$ wynosi zero;
przekroje uwidaczniają brak zbieżności jednostajnej.
Poniższy kod wymaga bibliotek NumPy i Plotly. Zapisuje samodzielny
plik HTML, w którym można obracać i powiększać wykres.
Rysunek \ref{ma2:fig:partial-sums-3d} przedstawia te same
krzywe; dla czytelności używamy na osi indeksu współrzędnej
$\log_2N$. Indeksy nadal są dyskretne.

\begin{figure}[htbp]
\centering
\begin{tikzpicture}
\definecolor{curve0}{HTML}{440154}
\definecolor{curve1}{HTML}{443983}
\definecolor{curve2}{HTML}{31688E}
\definecolor{curve3}{HTML}{21918C}
\definecolor{curve4}{HTML}{35B779}
\definecolor{curve5}{HTML}{90D743}
\definecolor{curve6}{HTML}{B8DE29}
\begin{axis}[width=11cm,height=7.4cm,view={40}{29},axis lines=box,grid=major,major grid style={draw=gray!22},axis line style={gray!65},tick style={gray!65},xmin=0,xmax=1.05,ymin=0,ymax=6.5,zmin=0,zmax=1.1,xlabel=$x$,ylabel={$N$ (log.)},zlabel={$S_N(x)$},ytick={0,1,2,3,4,5,6},yticklabels={1,2,4,8,16,32,64}]
\addplot3[very thick,curve0,domain=0:1,samples=201,samples y=1] (x,0,{x-x^(2)});
\addplot3[only marks,mark=*,mark size=1.2pt,curve0] coordinates {(1,0,0)};
\addplot3[very thick,curve1,domain=0:1,samples=201,samples y=1] (x,1,{x-x^(3)});
\addplot3[only marks,mark=*,mark size=1.2pt,curve1] coordinates {(1,1,0)};
\addplot3[very thick,curve2,domain=0:1,samples=201,samples y=1] (x,2,{x-x^(5)});
\addplot3[only marks,mark=*,mark size=1.2pt,curve2] coordinates {(1,2,0)};
\addplot3[very thick,curve3,domain=0:1,samples=201,samples y=1] (x,3,{x-x^(9)});
\addplot3[only marks,mark=*,mark size=1.2pt,curve3] coordinates {(1,3,0)};
\addplot3[very thick,curve4,domain=0:1,samples=201,samples y=1] (x,4,{x-x^(17)});
\addplot3[only marks,mark=*,mark size=1.2pt,curve4] coordinates {(1,4,0)};
\addplot3[very thick,curve5,domain=0:1,samples=201,samples y=1] (x,5,{x-x^(33)});
\addplot3[only marks,mark=*,mark size=1.2pt,curve5] coordinates {(1,5,0)};
\addplot3[very thick,curve6,domain=0:1,samples=201,samples y=1] (x,6,{x-x^(65)});
\addplot3[only marks,mark=*,mark size=1.2pt,curve6] coordinates {(1,6,0)};
\end{axis}
\end{tikzpicture}
\caption{Sumy częściowe $S_N(x)=x-x^{N+1}$ dla $N=1,2,4,8,16,32,64$. Oś indeksu ma skalę logarytmiczną, a każda krzywa kończy się w punkcie o wartości $S_N(1)=0$.}
\label{ma2:fig:partial-sums-3d}
\end{figure}

\begin{lstlisting}[language=Python]
import numpy as np
import plotly.graph_objects as go

x = np.linspace(0, 1, 501)
fig = go.Figure()
for n in [1, 2, 4, 8, 16, 32, 64]:
    fig.add_trace(go.Scatter3d(
        x=x, y=np.full_like(x, n), z=x-x**(n+1),
        mode="lines", name=f"N={n}"))
fig.update_layout(scene=dict(
    xaxis_title="x", yaxis_title="N", zaxis_title="S_N(x)"))
axis_style = dict(showbackground=True,
    backgroundcolor="#f5f6f7", gridcolor="#d5dbe0",
    zerolinecolor="#b0b8bf", showspikes=False)
fig.update_scenes(xaxis=axis_style, yaxis=axis_style,
                  zaxis=axis_style)
fig.write_html("am2_09_sums.html", include_plotlyjs=True)
fig.show()
\end{lstlisting}


\begin{przyklad}[Dwie własności łącznie nie wystarczają]
\label{ma2:ex:uniform-pointwise-absolute}
Jednostajna zbieżność szeregu i jego punktowa
zbieżność bezwzględna, nawet rozpatrywane łącznie,
nie gwarantują jednostajnej zbieżności szeregu norm.
Niech $v_k$ będzie funkcją trójkątną z przykładu
\ref{ma2:ex:absolute-uniform-nonnormal} pomnożoną
przez $k$. Mamy $0\le v_k\le1$, $v_k(c_k)=1$
i w każdym punkcie co najwyżej jedna z funkcji $v_k$
jest niezerowa.

Dla $k\ge1$ definiujemy blok $2k$ wyrazów
\[
 w_{k(k-1)+j}(x)=(-1)^{j-1}\frac{v_k(x)}k,
 \qquad 1\le j\le2k.
\]
Blok zaczyna się od indeksu $k(k-1)+1$,
a kończy na $k(k+1)$; bloki pokrywają wszystkie
dodatnie indeksy bez powtórzeń.
Suma każdego pełnego bloku jest zerowa.
Jeśli suma częściowa kończy się wewnątrz bloku $k$,
to jest równa zeru albo $v_k/k$ zależnie od
parzystości liczby wziętych wyrazów tego bloku.
Jej norma supremum nie przekracza $1/k$.
Ponieważ indeks bloku dąży do nieskończoności
wraz z indeksem sumy częściowej, szereg zbiega
jednostajnie do zera.

Dla ustalonego $x$ co najwyżej jeden blok
zawiera niezerowe wyrazy, więc szereg jest
punktowo bezwzględnie zbieżny. Jednak
\[
 \sum_{n=k(k-1)+1}^{k(k+1)}|w_n(c_k)|=2
 \qquad(k\ge1).
\]
Początki tych bloków dążą do nieskończoności.
Jednostajny warunek Cauchy'ego dla szeregu norm
nie zachodzi, na przykład dla $\varepsilon=1$.
\end{przyklad}

\section{Sumowanie przez części}

\begin{lemat}[Skończone sumowanie przez części]
\label{ma2:lem:summation-parts}
Niech $a_p,\ldots,a_q$ będą elementami $E$,
a $b_p,\ldots,b_q$ skalarami, gdzie $1\le p\le q$.
Połóżmy $A_{p-1}=0$ i $A_k=\sum_{n=p}^k a_n$.
Wówczas
\[
 \sum_{n=p}^q b_n a_n
 = b_q A_q+\sum_{n=p}^{q-1}(b_n-b_{n+1})A_n.
\]
\end{lemat}
\begin{proof}
Podstawiamy $a_n=A_n-A_{n-1}$:
\[
 \sum_{n=p}^q b_n a_n
 =\sum_{n=p}^q b_n A_n
   -\sum_{j=p-1}^{q-1}b_{j+1}A_j.
\]
Wyraz $j=p-1$ znika, a wyraz $b_qA_q$
występuje tylko w pierwszej sumie. Pozostałe
wyrazy łączymy przy wspólnym indeksie.
Dla $q=p$ druga suma w tezie jest pusta.
\end{proof}

\begin{twierdzenie}[Jednostajne kryterium Dirichleta]
\label{ma2:thm:uniform-dirichlet}
Niech $E$ będzie przestrzenią Banacha,
$a_n:X\to E$, $b_n:X\to[0,\infty)$.
Załóżmy, że istnieje $M\ge0$ takie, że
$\|\sum_{n=1}^k a_n(x)\|\le M$ dla wszystkich $k,x$,
że $b_{n+1}(x)\le b_n(x)$ dla wszystkich $n,x$
oraz że $b_n\to0$ jednostajnie.
Wtedy $\sum_n b_n a_n$ jest jednostajnie zbieżny.
Dla $q\ge p\ge1$ zachodzi oszacowanie
\[
 \left\|\sum_{n=p}^q b_n(x)a_n(x)\right\|
 \le 2M b_p(x).
\]
\end{twierdzenie}
\begin{proof}
Dla ustalonego $x$ i $k\ge p$ suma od $p$ do $k$
jest różnicą dwóch sum początkowych (przy czym
suma do indeksu $0$ wynosi zero), więc jej norma
nie przekracza $2M$: dla $P_k(x)=\sum_{n=1}^k a_n(x)$
mamy $\|P_k(x)-P_{p-1}(x)\|\le
\|P_k(x)\|+\|P_{p-1}(x)\|\le2M$.
Lemat \ref{ma2:lem:summation-parts} i monotoniczność
dają
\[
 \left\|\sum_{n=p}^q b_n(x)a_n(x)\right\|
 \le 2M\left(b_q(x)+
              \sum_{n=p}^{q-1}(b_n(x)-b_{n+1}(x))\right)
 =2M b_p(x).
\]
Jeśli $M=0$, wszystkie sumy początkowe
$P_k$ są zerowe, więc $a_k=P_k-P_{k-1}=0$
i szereg jest jednostajnie zbieżny do zera.
Jeśli $M>0$, z jednostajnej zbieżności $b_n$ do zera
wybieramy $N$ tak, że $b_p(x)<\varepsilon/(2M)$
dla wszystkich $p\ge N$ i $x\in X$.
Twierdzenie \ref{ma2:thm:function-series-cauchy}
daje jednostajną zbieżność. Przechodząc do granicy
$q\to\infty$ przy ustalonych $x,p$, otrzymujemy
takie samo oszacowanie nieskończonego ogona.
\end{proof}

\begin{przyklad}[Szereg trygonometryczny]
\label{ma2:ex:dirichlet-sine}
Niech $0<\delta<\pi$. Dla $x\in[\delta,2\pi-\delta]$
i $k\ge1$ z tożsamości dla skończonego szeregu
geometrycznego otrzymujemy
\[
 \left|\sum_{n=1}^k e^{inx}\right|
 =\left|\frac{e^{ix}(1-e^{ikx})}{1-e^{ix}}\right|
 \le\frac1{|\sin(x/2)|}
 \le\frac1{\sin(\delta/2)}.
\]
Korzystamy z tożsamości
$1-e^{ix}=e^{ix/2}(e^{-ix/2}-e^{ix/2})
=-2i e^{ix/2}\sin(x/2)$.
Wzór Eulera i $|e^{it}|=1$ wynikają z AM I,
rozdział 18, sekcja 3. Ponadto na
$[\delta/2,\pi-\delta/2]$ sinus jest nie mniejszy
niż $\sin(\delta/2)$ dzięki monotoniczności
na $[0,\pi/2]$ i symetrii $\sin(\pi-t)=\sin t$
(AM I, rozdział 18, sekcja 4).
Część urojona sumy
ma moduł nie większy od modułu sumy.
Twierdzenie \ref{ma2:thm:uniform-dirichlet}
dla $a_n(x)=\sin(nx)$ i $b_n(x)=1/n$
daje jednostajną zbieżność $\sum_n\sin(nx)/n$
na wskazanym przedziale oraz oszacowanie reszty
po $N$ wyrazach przez $2/((N+1)\sin(\delta/2))$.

Szereg nie jest jednostajnie zbieżny na $[0,2\pi]$.
Dla $p\ge1$ bierzemy $x_p=1/p$. Gdy $p\le n\le2p$,
mamy $1\le nx_p\le2$. Niech
$c=\min_{t\in[1,2]}\sin t>0$; minimum istnieje
z twierdzenia Weierstrassa i jest dodatnie,
gdyż $[1,2]\subset(0,\pi)$. Wtedy
\[
 \sum_{n=p}^{2p}\frac{\sin(nx_p)}n
 \ge c\frac{p+1}{2p}\ge c/2.
\]
Przeczy to jednostajnemu kryterium Cauchy'ego.
W każdym punkcie z $(0,2\pi)$ szereg jest jednak
zbieżny z pierwszej części przykładu, a na obu
końcach wszystkie wyrazy są zerowe.
\end{przyklad}

\begin{twierdzenie}[Jednostajne kryterium Abela]
\label{ma2:thm:uniform-abel}
Niech $E$ będzie przestrzenią Banacha,
$a_n:X\to E$, $b_n:X\to\R$.
Załóżmy, że $\sum_n a_n$ jest jednostajnie zbieżny,
że dla każdego $x$ ciąg $(b_n(x))$ jest monotoniczny
(kierunek może zależeć od $x$) oraz
że istnieje $B\ge0$ z $|b_n(x)|\le B$ dla
wszystkich $n,x$. Wówczas $\sum_n b_na_n$
jest jednostajnie zbieżny. Ponadto dla każdego
$p$ takiego, że
\[
 C_p=\sup_{\substack{x\in X\\k\ge p}}
       \left\|\sum_{n=p}^k a_n(x)\right\|<\infty,
\]
mamy $\sup_x\|\sum_{n=p}^\infty b_n(x)a_n(x)\|
\le3BC_p$. Wielkości $C_p$ są skończone od
pewnego indeksu i dążą do zera.
\end{twierdzenie}
\begin{proof}
Najpierw zauważmy, że z jednostajnego kryterium
Cauchy'ego dla $\sum a_n$ wynika $C_p\to0$:
dla dowolnego $\eta>0$ istnieje próg $N$, dla
którego wszystkie skończone ogony od $p\ge N$
mają normy mniejsze niż $\eta/2$, stąd
$C_p\le\eta/2$. W szczególności $C_p$ jest
skończone od pewnego indeksu.
Jeśli $B=0$, wszystkie iloczyny są zerowe,
więc zarówno zbieżność, jak i oszacowanie
dla każdego skończonego $C_p$ zachodzą.
Niech $B>0$ i $\varepsilon>0$.
Z jednostajnego kryterium Cauchy'ego dla
$\sum_n a_n$ wybieramy $N$ takie, że
$\|\sum_{n=p}^k a_n(x)\|<\varepsilon/(3B)$
dla wszystkich $k\ge p\ge N$ i $x$.
Dla ustalonych $x$ i $q\ge p\ge N$ monotoniczność
daje
\[
 \sum_{n=p}^{q-1}|b_n(x)-b_{n+1}(x)|
 =|b_p(x)-b_q(x)|\le2B.
\]
Z lematu \ref{ma2:lem:summation-parts}
i nierówności trójkąta otrzymujemy
\[
 \left\|\sum_{n=p}^q b_n(x)a_n(x)\right\|
 <\varepsilon
\]
jeśli suma modułów współczynników jest dodatnia,
gdyż nie przekracza $|b_q(x)|+2B\le3B$.
Jeżeli wszystkie współczynniki są zerowe,
lewa strona jest zerem. Twierdzenie
\ref{ma2:thm:function-series-cauchy} daje zbieżność.
Z tego samego lematu dla dowolnego $p$ z
$C_p<\infty$ otrzymujemy ograniczenie skończonych
ogonów przez $3BC_p$, niezależnie od $q\ge p$ i $x$.
Przejście do granicy po $q$ i wzięcie supremum
daje wskazane oszacowanie.
\end{proof}

\begin{przyklad}
\label{ma2:ex:abel-alternating-power}
Na $[0,1]$ szereg
$\sum_{n\ge1}(-1)^{n-1}x^n/n$ jest jednostajnie
zbieżny z kryterium Abela: $a_n(x)=(-1)^{n-1}/n$
daje jednostajnie zbieżny szereg stałych funkcji,
a $b_n(x)=x^n$ jest nierosnący i ograniczony przez $1$.
Szereg nie jest zbieżny normalnie, ponieważ
$\sup_{[0,1]}|a_n b_n|=1/n$. W tym konkretnym
przykładzie mamy nawet dokładniejsze oszacowanie
reszty $1/(N+1)$, bezpośrednio z twierdzenia
\ref{ma1:thm:leibniz-series} dla nierosnącego
ciągu $x^n/n$ (dla każdego ustalonego $x$).
\end{przyklad}
Kryterium Abela dotyczy tutaj zwykłej zbieżności
szeregu, a nie metody sumowania Abela omówionej
w AM I, rozdział 13.

\section{Zbieżność jednostajna na zbiorach zwartych}

Jeżeli $X$ jest przestrzenią topologiczną, to
szereg $\sum_n u_n$ jest jednostajnie zbieżny
na zbiorach zwartych, gdy ciąg jego sum częściowych
ma tę własność. Kryteria z tego rozdziału można
stosować oddzielnie na każdym niepustym zwartym
$K\subseteq X$, z progami i stałymi zależnymi od $K$.
Na zbiorach jednoelementowych uzyskujemy istnienie
sumy punktowej, więc sumy na różnych zwartych
podzbiorach są ograniczeniami tej samej funkcji.
Przykładowe zastosowanie kryterium Weierstrassa
do szeregów Fouriera znajduje się w skrypcie
równań różniczkowych cząstkowych, rozdział 06,
sekcja 1, lemat o zanikaniu współczynników gładkiej
funkcji (etykieta \texttt{lem:pde6-decay}).
Dowód tego lematu wykorzystamy w części o analizie
Fouriera, po wprowadzeniu współczynników Fouriera.

W szczególności zbieżność normalna na zbiorach
zwartych implikuje zbieżność jednostajną na zbiorach
zwartych w przestrzeni Banacha, z kryterium
\ref{ma2:thm:weierstrass-test}.

Jeżeli wszystkie wyrazy są ciągłe, suma takiego
szeregu jest ciągła dla metrycznej dziedziny
$X$ na mocy twierdzenia
\ref{ma2:thm:compact-uniform-continuity}.
Zachowanie różniczkowalności i możliwość
całkowania wyraz po wyrazie wymagają odrębnych
twierdzeń, którymi zajmiemy się w kolejnym rozdziale.

% AM I 21, 29, 32; AF 06, 14, 15A; AM II 08–09.
\chapter{Ciągłość, różniczkowanie i całkowanie granic}
\label{ma2:ch:limits-calculus}

W rozdziałach \ref{ma2:ch:uniform-convergence}
i \ref{ma2:ch:function-series} ustaliliśmy warunki,
pod którymi granica funkcji ciągłych pozostaje
ciągła. Zbadamy teraz przemienność granicy
z innymi operacjami analizy. W przypadku
wyników udowodnionych w tomie I podajemy
odsyłacze zamiast ponownych dowodów.

\section{Zamiana kolejności granic}

\begin{twierdzenie}
\label{ma2:thm:interchange-limits}
Niech $X$ będzie przestrzenią metryczną,
$A\subseteq X$, $x_0$ punktem skupienia $A$,
a $E$ przestrzenią Banacha. Niech
$f_n:A\setminus\{x_0\}\to E$ zbiegają
jednostajnie do $f$ oraz niech dla każdego $n$
istnieje granica $a_n=\lim_{x\to x_0,\ x\in A\setminus\{x_0\}}f_n(x)$.
Wtedy istnieją obie granice
\[
 \lim_{x\to x_0,\ x\in A\setminus\{x_0\}} f(x)
 =\lim_{n\to\infty} a_n.
\]
\end{twierdzenie}
\begin{proof}
Z jednostajnej zbieżności wynika, że dla
$\varepsilon>0$ istnieje $N$ z
$\|f_n(x)-f_m(x)\|<\varepsilon/2$
dla wszystkich $m,n\ge N$ i wszystkich
$x\in A\setminus\{x_0\}$.
Dla ustalonych $n,m$ funkcja $f_n-f_m$ ma
granicę $a_n-a_m$ z ciągłości odejmowania,
a norma jest ciągła (twierdzenie
\ref{ma2:thm:metric}). Przechodząc do granicy
przy $x\to x_0$, otrzymujemy
$\|a_n-a_m\|\le\varepsilon/2$.
Z zupełności $E$ istnieje $a=\lim_n a_n$.
Dla $\varepsilon>0$ wybieramy $n$ tak duże,
aby $\|a_n-a\|<\varepsilon/3$ oraz
$\|f_n(x)-f(x)\|<\varepsilon/3$ dla wszystkich $x$.
Z definicji $a_n$ istnieje $\delta>0$,
dla którego $0<d(x,x_0)<\delta$, $x\in A$,
implikuje $\|f_n(x)-a_n\|<\varepsilon/3$.
Wtedy
\[
 \|f(x)-a\|\le
 \|f(x)-f_n(x)\|+\|f_n(x)-a_n\|+\|a_n-a\|
 <\varepsilon.
\]
\end{proof}
Wystarcza zbieżność jednostajna na pewnym
ustalonym przebitym otoczeniu $x_0$ w $A$.
Jeżeli granica $(a_n)$ już istnieje w $E$,
ostatni argument nie wymaga zupełności.
Bez jednostajności zamiana może zawodzić:
dla $f_n(x)=x^n$ na $[0,1)$ mamy
$\lim_n\lim_{x\to1^-} f_n(x)=1$ oraz
$\lim_{x\to1^-}\lim_n f_n(x)=0$ (por.
przykład \ref{ma2:ex:powers-pointwise}).

\section{Granica pod znakiem całki Riemanna}

Przypomnijmy pojęcie całki Riemanna z tomu I,
rozdział 29. Dla $a<b$ i podziału z punktami
pośrednimi $P=(a=t_0<\cdots<t_m=b;\
\xi_j\in[t_{j-1},t_j])$ piszemy
\[
 |P|=\max_j(t_j-t_{j-1}),\qquad
 R(f,P)=\sum_{j=1}^{m}(t_j-t_{j-1})f(\xi_j).
\]
Jeśli $f:[a,b]\to E$, gdzie $E$ jest przestrzenią
unormowaną, to całkowalność w sensie Riemanna
oznacza istnienie $I\in E$ takiego, że
dla każdego $\varepsilon>0$ istnieje $\delta>0$
i $\|R(f,P)-I\|<\varepsilon$ dla wszystkich
podziałów z punktami pośrednimi o $|P|<\delta$.
Piszemy $I=\int_a^b f(t)\,\mathrm dt$.
Granica jest jednoznaczna: jeżeli $I,J$
spełniają definicję, to dla $\varepsilon>0$
wybieramy podział o średnicy tak małej, aby
$\|R(f,P)-I\|<\varepsilon/2$ oraz
$\|R(f,P)-J\|<\varepsilon/2$. Wtedy
$\|I-J\|<\varepsilon$, a dowolność
$\varepsilon$ daje $I=J$.
Każda funkcja całkowalna w tym sensie jest
ograniczona. Istotnie, wybieramy jeden podział
o średnicy dostatecznie małej, aby wszystkie jego
sumy z dowolnymi punktami pośrednimi różniły się
od $I$ w normie o mniej niż $1$. Ustalamy punkty
pośrednie i zmieniamy tylko punkt w $j$-tym
podprzedziale. Różnica sum ma normę
$(t_j-t_{j-1})\|f(t)-f(\xi_j)\|<2$, więc
$\|f(t)\|\le\|f(\xi_j)\|+2/(t_j-t_{j-1})$
na tym podprzedziale. Maksimum skończonej
liczby takich ograniczeń daje ograniczoność.
Dla funkcji rzeczywistych jest to zatem definicja
równoważna definicji Darboux (AM I, twierdzenie
29.2.2, etykieta \texttt{ma1:thm:riemann-darboux-equivalence}).
W przypadku funkcji ciągłych o wartościach
w przestrzeni Banacha istnienie całki zostało
udowodnione w analizie funkcjonalnej, rozdział 15A,
twierdzenie o istnieniu całki przy zupełności
ciągowej (\texttt{thm:riemann-existence15a}).
Pełniejszą teorię całki wektorowej omówimy
w rozdziale 18.

\begin{twierdzenie}[Skalarna granica pod znakiem całki]
\label{ma1:thm:riemann-uniform-limit}
Niech $a<b$ i niech $f_n:[a,b]\to\R$ będą całkowalne
w sensie Riemanna i niech $f_n\to f$ jednostajnie.
Wówczas $f$ jest całkowalna w sensie Riemanna i
\[
 \left|\int_a^b f_n(t)\,\mathrm dt-
              \int_a^b f(t)\,\mathrm dt\right|
 \le (b-a)\sup_{t\in[a,b]}|f_n(t)-f(t)|
 \longrightarrow0.
\]
\end{twierdzenie}
\emph{Dowód:} AM I, twierdzenie 32.1.1, wraz
z oszacowaniem w jego dowodzie. Zachowano
oryginalną etykietę.

\begin{twierdzenie}[Granica całek wektorowych]
\label{ma2:thm:vector-riemann-uniform}
Niech $E$ będzie przestrzenią Banacha,
$a<b$, a $f_n:[a,b]\to E$ będą całkowalne
w sensie Riemanna. Jeżeli $f_n\to f$
jednostajnie, to $f$ jest całkowalna oraz
\[
 \left\|\int_a^b f_n(t)\,\mathrm dt-
                 \int_a^b f(t)\,\mathrm dt\right\|
 \le (b-a)\sup_{t\in[a,b]}\|f_n(t)-f(t)\|.
\]
W szczególności całki $f_n$ dążą do całki $f$.
\end{twierdzenie}
\begin{proof}
Oznaczmy $L=b-a>0$ i $I_n=\int_a^b f_n(t)\,\mathrm dt$.
Dla dowolnego podziału z punktami pośrednimi $P$
i funkcji $g,h$ mamy
\[
 \|R(g,P)-R(h,P)\|\le L\sup_{t\in[a,b]}\|g(t)-h(t)\|.
\]
Istotnie, różnica sum jest sumą wektorów
$(t_j-t_{j-1})(g(\xi_j)-h(\xi_j))$.
Z nierówności trójkąta jej norma nie przekracza
$\sum_j(t_j-t_{j-1})\|g(\xi_j)-h(\xi_j)\|$,
co po oszacowaniu każdej normy przez supremum
daje tezę, gdyż $\sum_j(t_j-t_{j-1})=L$.
Jeżeli $g,h$ są całkowalne, przechodzimy do
granicy przy $|P|\to0$, uzyskując tę samą
nierówność dla różnicy ich całek.
Dla $g=f_n$, $h=f_m$ z jednostajnego
warunku Cauchy'ego (twierdzenie
\ref{ma2:thm:uniform-cauchy}) wynika,
że $(I_n)$ jest ciągiem Cauchy'ego w $E$.
Z zupełności istnieje $I=\lim_n I_n$.
Ponadto $f$ jest ograniczona, bo z jednostajnej
zbieżności dla pewnego $n$ mamy
$\sup_t\|f(t)-f_n(t)\|\le1$, a $f_n$ jest ograniczona.
Ustalmy $\varepsilon>0$ i wybierzmy $n$ tak,
aby $\sup_t\|f_n(t)-f(t)\|<\varepsilon/(3L)$
i $\|I_n-I\|<\varepsilon/3$.
Z całkowalności $f_n$ istnieje $\delta>0$
takie, że dla $|P|<\delta$ mamy
$\|R(f_n,P)-I_n\|<\varepsilon/3$.
Wtedy
\[
 \|R(f,P)-I\|
 \le \|R(f,P)-R(f_n,P)\|
    +\|R(f_n,P)-I_n\|+\|I_n-I\|<\varepsilon.
\]
To jest definicja całkowalności $f$ i jej całki $I$.
Stosując początkowe oszacowanie do całkowalnych
$f_n,f$, otrzymujemy tezę.
\end{proof}
Liniowość całki w tej definicji wynika z liniowości
sum Riemanna i ciągłości działań w $E$:
dla całkowalnych $g,h$ oraz skalarów $\alpha,\beta$
mamy $R(\alpha g+\beta h,P)=\alpha R(g,P)+\beta R(h,P)$;
granica po obu stronach daje istnienie całki
kombinacji i jej wartość.

\begin{wniosek}[Całkowanie szeregu wyraz po wyrazie]
\label{ma2:cor:termwise-integral}
Niech $E$ będzie przestrzenią Banacha,
$a<b$, a $u_n:[a,b]\to E$ funkcjami całkowalnymi
w sensie Riemanna. Jeżeli $\sum_n u_n$
jest jednostajnie zbieżny do $S$, to $S$
jest całkowalna i
\[
 \int_a^b S(t)\,\mathrm dt
 =\sum_{n=1}^\infty\int_a^b u_n(t)\,\mathrm dt.
\]
Błąd po $N$ wyrazach ma normę nie większą niż
$(b-a)\sup_t\|S(t)-S_N(t)\|$.
\end{wniosek}
\begin{proof}
Sumy częściowe $S_N$ są całkowalne
z liniowości całki, a z założenia $S_N\to S$
jednostajnie. Zastosowanie twierdzenia
\ref{ma2:thm:vector-riemann-uniform} i ponowne
użycie liniowości daje równość oraz oszacowanie.
\end{proof}
Dla $a=b$ całka jest zerowa, a dla $b<a$
zmieniamy orientację; w oszacowaniach
używamy wtedy $|b-a|$.

\begin{przyklad}
\label{ma2:ex:termwise-geometric-integral}
Dla $0\le r<1$ szereg $\sum_{n=0}^\infty t^n$
jest jednostajnie zbieżny na $[0,r]$ do $1/(1-t)$
(przykład \ref{ma2:ex:function-geometric}).
Całkując wyraz po wyrazie, otrzymujemy
\[
 -\log(1-r)=\sum_{n=0}^\infty\frac{r^{n+1}}{n+1}.
\]
Dla $r>0$ użyto funkcji pierwotnych
$-\log(1-t)$ i $t^{n+1}/(n+1)$
(AM I, stwierdzenie 28.4.1)
oraz twierdzenia Newtona--Leibniza
(AM I, twierdzenie 29.6.3).
Dla $r=0$ obie strony równości są zerowe. Błąd po wyrazie
o indeksie $N$ nie przekracza
$r^{N+2}/(1-r)$, z oszacowania w przykładzie
\ref{ma2:ex:function-geometric} i wniosku
\ref{ma2:cor:termwise-integral}.
\end{przyklad}

\begin{przyklad}[Sama zbieżność punktowa nie wystarcza]
\label{ma1:ex:parameter-moving-spike}
Dla $n\ge2$ niech
\[
 f_n(t)=\begin{cases}
 n^2t,&0\le t\le1/n,\\
 2n-n^2t,&1/n\le t\le2/n,\\
 0,&2/n\le t\le1.
 \end{cases}
\]
Każda funkcja jest ciągła i $f_n\to0$ punktowo
na $[0,1]$, lecz $\int_0^1 f_n(t)\,\mathrm dt=1$.
Weryfikację rachunków i granicy zawiera AM I,
przykład 32.1.3. Zatem granica całek
nie jest całką granicy.
\end{przyklad}

\section{Granica pod znakiem pochodnej}

Pochodną $f:I\to E$ względem rzeczywistej
zmiennej definiujemy w normie $E$ jako granicę
ilorazów różnicowych. Na końcach domkniętego
przedziału rozumiemy ją jednostronnie,
o ile ją rozważamy.

\begin{lemat}[Oszacowanie przyrostu]
\label{ma2:lem:vector-increment}
Niech $E$ będzie przestrzenią unormowaną nad
$\R$ lub $\C$, $g:[u,v]\to E$ ciągła,
różniczkowalna w $(u,v)$, gdzie $u<v$, $M\ge0$.
Jeżeli $\|g'(t)\|\le M$ na $(u,v)$, to
\[
 \|g(v)-g(u)\|\le M(v-u).
\]
\end{lemat}
\begin{proof}
Dla $g(v)=g(u)$ lewa strona jest zerem,
a prawa jest nieujemna, więc teza zachodzi.
W przeciwnym przypadku traktujemy $E$ jako
rzeczywistą przestrzeń unormowaną i wybieramy
rzeczywisty ciągły funkcjonał liniowy $\ell$
o normie $1$, taki że
$\ell(g(v)-g(u))=\|g(v)-g(u)\|$.
Istnienie wynika z twierdzenia Hahna--Banacha
(analiza funkcjonalna, rozdział 06, wniosek
o funkcjonale normującym, \texttt{cor:norming}).
Skalarna funkcja $\ell\circ g$ jest ciągła
na odcinku i różniczkowalna wewnątrz, a jej
pochodna wynosi $\ell(g'(t))$ z ciągłości
$\ell$ i definicji pochodnej.
Z twierdzenia Lagrange'a (AM I, twierdzenie 21.3.1)
istnieje $c\in(u,v)$ z
\[
 \|g(v)-g(u)\|=\ell(g'(c))(v-u)
 \le\|g'(c)\|(v-u)\le M(v-u).
\]
Pierwsza nierówność wynika z $\|\ell\|=1$
i $v-u>0$.
\end{proof}
W przypadku skalarnym rzeczywistym jest to
szczególny przypadek twierdzenia 21.6.1 z AM I.
Nie zakładamy tutaj ciągłości pochodnej.

\begin{twierdzenie}[Jednostajna zbieżność pochodnych]
\label{ma2:thm:uniform-derivative}
Niech $I$ będzie przedziałem o niepustym wnętrzu,
a $E$ przestrzenią Banacha. Niech $f_n:I\to E$
będą ciągłe i różniczkowalne we wnętrzu $I$.
Załóżmy, że dla pewnego $x_0\in I$
ciąg $(f_n(x_0))$ jest zbieżny w $E$ i że
$f_n'\to g:\operatorname{int}I\to E$ jednostajnie
na $\operatorname{int}I$.
Wtedy $f_n$ zbiegają jednostajnie na każdym
ograniczonym podprzedziale $J\subseteq I$ do funkcji
$f:I\to E$ ciągłej i różniczkowalnej we wnętrzu,
a $f'=g$. Jeśli $I$ jest ograniczony, zbieżność
jest jednostajna na całym $I$.
\end{twierdzenie}
\begin{proof}
Z jednostajnej zbieżności pochodnych dla każdego
$\eta>0$ istnieje $N$ taki, że dla $n,m\ge N$
i $t\in\operatorname{int}I$ mamy
$\|f_n'(t)-f_m'(t)\|<\eta$.
Lemat \ref{ma2:lem:vector-increment} zastosowany
do $f_n-f_m$ na odcinku o końcach $x_0,x$
daje
\[
 \|f_n(x)-f_m(x)\|
 \le\|f_n(x_0)-f_m(x_0)\|+\eta|x-x_0|.
\]
Dla $x=x_0$ nierówność zachodzi bez lematu.
Na dowolnym niepustym ograniczonym
podprzedziale $J\subseteq I$ liczba $D=\sup_{x\in J}|x-x_0|$ jest skończona.
Dla $\varepsilon>0$ wybieramy
$\eta=\varepsilon/(2(1+D))$ oraz zwiększamy
$N$, aby pierwszy składnik był mniejszy niż
$\varepsilon/2$. Otrzymujemy jednostajny
warunek Cauchy'ego na $J$.
Z twierdzenia \ref{ma2:thm:uniform-cauchy}
istnieje granica jednostajna na $J$.
Granice na różnych podprzedziałach zgadzają się
w punktach wspólnych z jednoznaczności granicy
punktowej; definiują $f$ na całym $I$.
Każdy punkt $I$ ma ograniczone otoczenie względne,
więc ciągłość $f$ wynika z twierdzenia
\ref{ma2:thm:uniform-continuous-limit}.
Dla ograniczonego $I$ bierzemy $J=I$.
Otrzymujemy też oszacowanie błędu. Dla każdego
$n$ takiego, że
$A_n=\sup_{t\in\operatorname{int}I}\|f_n'(t)-g(t)\|$
jest skończone, zachodzi
\[
 \|f_n(x)-f(x)\|
 \le\|f_n(x_0)-f(x_0)\|+A_n|x-x_0|
 \qquad(x\in I).
\]
Istotnie, dla dostatecznie dużych $m$ norma
pochodnej $f_n-f_m$ jest nie większa niż
$A_n+A_m$, a $A_m\to0$ z jednostajnej
zbieżności pochodnych. Stosujemy oszacowanie
przyrostu do $f_n-f_m$ i przechodzimy po $m$
do granicy. W szczególności na ograniczonym
$J$ błąd jednostajny jest nie większy niż
$\|f_n(x_0)-f(x_0)\|+DA_n$.

Ustalmy $x\in\operatorname{int}I$ i $\varepsilon>0$.
Wybieramy $n$ tak, aby
$\sup_t\|f_m'(t)-f_n'(t)\|\le\varepsilon/3$
dla wszystkich $m\ge n$ i
$\|f_n'(x)-g(x)\|<\varepsilon/3$.
Dla $h\ne0$ z $x+h\in I$ lemat daje
\[
 \|(f_m-f_n)(x+h)-(f_m-f_n)(x)\|
 \le(\varepsilon/3)|h|.
\]
Przechodząc po $m$ do granicy, otrzymujemy
\[
 \left\|\frac{f(x+h)-f(x)}h-
             \frac{f_n(x+h)-f_n(x)}h\right\|
 \le\varepsilon/3.
\]
Przy ustalonym $n$ dla dostatecznie małego
$|h|>0$ iloraz różnicowy $f_n$ różni się
od $f_n'(x)$ w normie o mniej niż $\varepsilon/3$.
Nierówność trójkąta pokazuje, że iloraz
różnicowy $f$ różni się od $g(x)$ o mniej niż
$\varepsilon$. Stąd $f'(x)=g(x)$.
\end{proof}
Jeżeli wszystkie $f_n'$ są ciągłe, to $g$
jest ciągła jako granica jednostajna, zatem $f$
jest klasy $C^1$ we wnętrzu $I$.
Jeżeli $x_0\in\operatorname{int}I$, wystarczy
założyć zbieżność jednostajną pochodnych na
każdym zwartym podprzedziale zawartym we
wnętrzu $I$, aby otrzymać zbieżność $f_n$
jednostajną na zbiorach zwartych tego wnętrza
i równość $f'=g$ we wnętrzu. Istotnie, każdy
niepusty zwarty podzbiór wnętrza wraz z $x_0$
mieści się w zwartym podprzedziale wnętrza;
stosujemy na nim twierdzenie i łączymy
granice punktowe. W celu obliczenia pochodnej
w ustalonym punkcie wybieramy podprzedział
zawierający także jego otoczenie.
Nie wynika stąd zbieżność na końcach $I$,
jeśli do niego należą: na $[0,1]$ funkcje
$f_n(x)=x^n$ spełniają $f_n(0)=0$,
a na każdym zwartym podprzedziale $(0,1)$
ich pochodne $nx^{n-1}$ zbiegają jednostajnie
do zera. Dla $0<r<1$ mamy bowiem
$\sup_{0\le x\le r}|f_n'(x)|=nr^{n-1}\to0$
(iloraz kolejnych dodatnich wyrazów $nr^{n-1}$
dąży do $r<1$, więc szereg tych wyrazów jest
zbieżny z kryterium ilorazowego, AM I,
twierdzenie 9.6.1).
Mimo to granica punktowa na $[0,1]$ nie jest
ciągła w $1$. W tym przykładzie jako punkt
normalizacji można wybrać dowolny $x_0\in(0,1)$,
bo $x_0^n\to0$. Co więcej, dla $F_n(x)=nx^n$
i $x_0=1/2$ mamy $F_n(x_0)\to0$,
a na każdym zwartym podprzedziale $(0,1)$
pochodne $n^2x^{n-1}$ zbiegają jednostajnie
do zera: dla $0<r<1$ ciąg $n^2r^{n-1}$
dąży do zera z tego samego kryterium
ilorazowego (iloraz dąży do $r$).
Jednak $F_n(1)=n$ nie ma skończonej granicy.
Wariant lokalny nie zapewnia więc nawet
istnienia granicy na końcu przedziału.

\begin{wniosek}[Różniczkowanie szeregu wyraz po wyrazie]
\label{ma2:cor:termwise-derivative}
Niech $I$ będzie przedziałem o niepustym
wnętrzu, $E$ przestrzenią Banacha, a funkcje
$u_n:I\to E$ ciągłe i różniczkowalne
we wnętrzu $I$.
Jeżeli $\sum_n u_n(x_0)$ jest zbieżny dla pewnego
$x_0\in I$, a $\sum_n u_n'$ zbiega jednostajnie
we wnętrzu $I$, to $\sum_n u_n$ jest jednostajnie
zbieżny na każdym ograniczonym podprzedziale $J\subseteq I$
do różniczkowalnej we wnętrzu funkcji $S$ i
\[
 S'(x)=\sum_{n=1}^\infty u_n'(x)
 \qquad(x\in\operatorname{int}I).
\]
\end{wniosek}
\begin{proof}
Dla sum częściowych $S_N$ mamy
$S_N'=\sum_{n=1}^N u_n'$ z liniowości pochodnej.
Hipotezy są dokładnie hipotezami twierdzenia
\ref{ma2:thm:uniform-derivative} dla $(S_N)$.
\end{proof}

\section{Kontrprzykłady i zakres hipotez}

\begin{przyklad}[Jednostajna granica może nie być różniczkowalna]
\label{ma2:ex:uniform-limit-cusp}
Na $[-1,1]$ niech $f_n(x)=\sqrt{x^2+n^{-2}}$.
Funkcje są klasy $C^1$ oraz
\[
 0\le f_n(x)-|x|
 =\frac{n^{-2}}{\sqrt{x^2+n^{-2}}+|x|}
 \le\frac1n.
\]
W punkcie $x=0$ różnica wynosi $1/n$,
więc $\|f_n-|\cdot|\|_\infty=1/n$.
Zatem $f_n\to |x|$ jednostajnie. Funkcja $|x|$
nie jest różniczkowalna w zerze, ponieważ
ilorazy różnicowe z prawej i lewej strony
wynoszą odpowiednio $1$ i $-1$.
Pochodne $f_n'(x)=x/\sqrt{x^2+n^{-2}}$ zbiegają
punktowo do funkcji znaku z wartością $0$ w zerze.
Ta funkcja jest nieciągła w zerze, więc
zbieżność pochodnych nie jest jednostajna.
\end{przyklad}

\begin{figure}[htbp]
\centering
\begin{tikzpicture}[x=5cm,y=3.3cm,>=Stealth]
 \draw[->] (-1.08,0)--(1.1,0) node[right]{$x$};
 \draw[->] (0,-.06)--(0,1.28) node[above]{$y$};
 \draw[black,very thick] (-1,1)--(0,0)--(1,1);
 \foreach \n/\kol in {2/blue!75!black,4/teal!75!black,8/orange!85!black}{
  \draw[\kol,thick,domain=-1:1,samples=151]
      plot (\x,{sqrt(\x*\x+1/(\n*\n))});
  \fill[\kol] (0,{1/\n}) circle[radius=1.2pt];
 }
 % Legenda poniżej prawego ramienia granicy, bez przecinania krzywych.
 \draw[blue!75!black,thick] (.56,.40)--(.68,.40)
      node[right,black,font=\small] {$n=2$};
 \draw[teal!75!black,thick] (.56,.27)--(.68,.27)
      node[right,black,font=\small] {$n=4$};
 \draw[orange!85!black,thick] (.56,.14)--(.68,.14)
      node[right,black,font=\small] {$n=8$};
 \node[below] at (-1,0) {$-1$};
 \node[below] at (1,0) {$1$};
 \node[below left] at (0,0) {$0$};
\end{tikzpicture}
\caption{Gładkie funkcje z przykładu
\ref{ma2:ex:uniform-limit-cusp} i ich granica
jednostajna $|x|$ (czarna). Różnica w zerze
wynosi dokładnie $1/n$.}
\label{ma2:fig:uniform-limit-cusp}
\end{figure}
\FloatBarrier

\begin{przyklad}
\label{ma2:ex:derivative-mismatch}
Na $[-1,1]$ funkcje $f_n(x)=\sin(nx)/n$
zbiegają jednostajnie do zera, gdyż
$|f_n(x)|\le1/n$. Tymczasem $f_n'(0)=1$
dla każdego $n$, podczas gdy pochodna granicy
w zerze wynosi $0$. Nawet przy różniczkowalnej
granicy nie można bez dodatkowych założeń
zamienić pochodnej z granicą.
Co więcej, sama punktowa zbieżność pochodnych
także nie wystarcza. Na $[-1,1]$ funkcje
$h_n(x)=x/(1+nx^2)$ spełniają
$|h_n(x)|\le 1/(2\sqrt n)$, gdyż
$1+nx^2\ge 2\sqrt n|x|$.
Ich pochodne są równe
\[
 h_n'(x)=\frac{1-nx^2}{(1+nx^2)^2}.
\]
Dla $x=0$ mają wartość $1$, dla $x\ne0$
dążą do zera po podzieleniu licznika
i mianownika przez $n^2$.
Zatem $h_n\to0$ jednostajnie, pochodne
mają granicę punktową wszędzie, lecz
granica ta w zerze nie jest pochodną
funkcji granicznej.
\end{przyklad}

\begin{przyklad}[Znaczenie wartości w jednym punkcie]
\label{ma2:ex:derivative-anchor}
Na dowolnym niepustym przedziale funkcje stałe
$f_n(x)=n$ mają pochodne identycznie zerowe,
lecz nie mają granicy punktowej. Warunek
zbieżności wartości w pewnym punkcie
w twierdzeniu \ref{ma2:thm:uniform-derivative}
eliminuje tę dowolność stałej.
Na nieograniczonym przedziale $\R$ funkcje
$g_n(x)=x/n$ spełniają $g_n(0)=0$,
a pochodne $g_n'=1/n$ zbiegają jednostajnie
do zera. Funkcje $g_n$ zbiegają jednostajnie
na ograniczonych podprzedziałach, ale nie
na całej $\R$: dla $x=n$ mamy $g_n(n)=1$.
\end{przyklad}

W następnym rozdziale zastosujemy powyższe
twierdzenia do szeregów potęgowych.

\chapter{Granice i ciągłość funkcji wielu zmiennych}
\label{ma2:ch:multivariable-continuity}

Ogólne wyniki o granicach i ciągłości w przestrzeniach
metrycznych znamy z AM I, rozdziały 14--16. Przytaczamy je
z odsyłaczami do dowodów; skupimy się na konsekwencjach
wielowymiarowych, przykładach i kontrprzykładach.

\section{Przestrzenie współrzędnych i granice}

Niech $n,m\ge1$, $A\subseteq\R^n$, $f:A\to\R^m$,
$f=(f_1,\ldots,f_m)$, gdzie $f_j:A\to\R$.
Domyślnie stosujemy normę euklidesową
$\|x\|_2=(\sum_{j=1}^n x_j^2)^{1/2}$. Z twierdzeń
\ref{ma2:thm:equiv} i \ref{ma2:thm:finite}
wynika niezależność zbieżności, ciągłości i zupełności
od wyboru normy. W szczególności
\[
 \|x\|_\infty\le\|x\|_2\le\sqrt n\,\|x\|_\infty,
 \qquad \|x\|_\infty=\max_{1\le j\le n}|x_j|.
\]
Pierwsza nierówność wynika z uwzględnienia największego
składnika sumy kwadratów, druga z ograniczenia każdego
składnika przez $\|x\|_\infty^2$.

Przypomnijmy definicje AM I, 14.1.1 i 14.1.3.
Punkt $p\in\R^n$ jest punktem skupienia zbioru $A$, gdy
dla każdego $\delta>0$ istnieje $x\in A$ z
$0<\|x-p\|_2<\delta$. Zbiór tych punktów oznaczamy $A'$.
Dla $p\in A'$ i $L\in\R^m$ warunek
$\lim_{A\ni x\to p}f(x)=L$ oznacza
\[
 \forall\varepsilon>0\ \exists\delta>0\ \forall x\in A:
 \quad 0<\|x-p\|_2<\delta\Longrightarrow
 \|f(x)-L\|_2<\varepsilon.
\]
Punkt $p$ nie musi należeć do $A$.

\begin{twierdzenie}[Kryteria granicy]
\label{ma2:thm:multi-limit}
Niech $p\in A'$ i $L=(L_1,\ldots,L_m)\in\R^m$.
Następujące warunki są równoważne:
\begin{enumerate}[label=(\roman*)]
 \item $\lim_{A\ni x\to p}f(x)=L$;
 \item dla każdego ciągu $(x_k)$ w $A\setminus\{p\}$
 zbieżnego do $p$ zachodzi $f(x_k)\to L$;
 \item dla każdego $j=1,\ldots,m$ zachodzi
 $\lim_{A\ni x\to p}f_j(x)=L_j$.
\end{enumerate}
\end{twierdzenie}
\begin{proof}
Równoważność (i) i (ii) jest szczególnym przypadkiem
twierdzenia AM I, 14.1.5.
Dla (i)$\Rightarrow$(iii) korzystamy z nierówności
$|f_j(x)-L_j|\le\|f(x)-L\|_2$.
Dla (iii)$\Rightarrow$(i) ustalamy $\varepsilon>0$ i
dla każdej współrzędnej wybieramy $\delta_j>0$ dla
dokładności $\varepsilon/\sqrt m$.
Przy $\delta=\min_j\delta_j>0$ i
$0<\|x-p\|_2<\delta$ mamy
\[
 \|f(x)-L\|_2^2=\sum_{j=1}^m|f_j(x)-L_j|^2
 < m(\varepsilon/\sqrt m)^2=\varepsilon^2.
\]
To daje (i).
\end{proof}

\begin{wniosek}
\label{ma2:cor:multi-limit-cauchy}
Dla $p\in A'$ granica $f$ w $p$ istnieje w $\R^m$ wtedy
i tylko wtedy, gdy dla każdego $\varepsilon>0$ istnieje
$\delta>0$ takie, że $\|f(x)-f(y)\|_2<\varepsilon$
dla wszystkich $x,y\in A$ spełniających
$0<\|x-p\|_2<\delta$, $0<\|y-p\|_2<\delta$.
\end{wniosek}
Jest to kryterium Cauchy'ego z AM I, twierdzenie 14.2.4:
jego założenie zupełności przestrzeni wartości jest spełnione
na mocy twierdzenia \ref{ma2:thm:finite}.
Rachunek granic skalarnych i twierdzenie o trzech funkcjach
stosujemy z AM I, twierdzenia 14.3.1--2, a do odwzorowań
wektorowych stosujemy kryterium współrzędnych.

\section{Granice wzdłuż krzywych i granice iterowane}

Jeżeli $\lim_{A\ni x\to p}f(x)=L$ oraz
$\gamma:(0,r)\to A\setminus\{p\}$ spełnia
$\lim_{t\to0^+}\gamma(t)=p$, to
$\lim_{t\to0^+}f(\gamma(t))=L$ na mocy twierdzenia
AM I, 14.2.2: argument $\gamma(t)$ unika $p$.
Różne granice wzdłuż dwóch takich krzywych wykluczają
granicę wspólną.

\begin{przyklad}[Proste nie wystarczają]
\label{ma2:ex:multivariable-paths}
Niech $f:\R^2\to\R$ będzie dana wzorem
\[
 f(x,y)=
 \begin{cases}
 \dfrac{x^2y}{x^4+y^2},&(x,y)\ne(0,0),\\
 0,&(x,y)=(0,0).
 \end{cases}
\]
Dla $(a,b)\ne(0,0)$, $t\ne0$ i $b\ne0$ mamy
\[
 f(ta,tb)=\frac{ta^2b}{t^2a^4+b^2}\longrightarrow0.
\]
Gdy $b=0$, wartości na prostej są zerowe.
Ograniczenie do każdej prostej przez początek jest
ciągłe w początku. Jednak dla $t\ne0$ mamy
$f(t,t^2)=1/2$, podczas gdy $f(t,0)=0$.
Granica w $(0,0)$ nie istnieje, a $f$ nie jest tam ciągła.
\end{przyklad}

\begin{figure}[htbp]
\centering
\begin{tikzpicture}[>=Stealth,scale=1.25]
 \draw[->] (-1.6,0)--(1.7,0) node[right] {$x$};
 \draw[->] (0,-.5)--(0,1.9) node[above] {$y$};
 \draw[blue!70!black,thick,domain=-1.2:1.2,samples=65]
 plot (\x,{\x*\x});
 \draw[red!70!black,very thick,->] (1.25,0)--(.45,0);
 \draw[blue!70!black,very thick,->]
 plot[domain=1.05:.65,samples=20] (\x,{\x*\x});
 \fill (0,0) circle (1.4pt) node[below left] {$0$};
 \node[blue!70!black,anchor=west] at (1.0,1.52) {$y=x^2$};
 \node[red!70!black] at (.9,-.3) {$y=0$};
 \node[blue!70!black,anchor=east] at (-.95,1.5) {$f=1/2$};
 \node[red!70!black] at (-.9,.2) {$f=0$};
\end{tikzpicture}
\caption{Dwie drogi do początku układu z różnymi granicami
funkcji z przykładu \ref{ma2:ex:multivariable-paths}.
Wartości na paraboli dotyczą punktów różnych od początku.}
\label{ma2:fig:multivariable-paths}
\end{figure}

\begin{przyklad}[Granice iterowane]
Dla tej samej funkcji i każdego ustalonego $x$ zachodzi
$\lim_{y\to0}f(x,y)=0$: dla $x\ne0$ mianownik dąży
do $x^4>0$, a dla $x=0$ funkcja jest zerowa.
Dla każdego ustalonego $y$ analogicznie
$\lim_{x\to0}f(x,y)=0$ (dla $y\ne0$ mianownik dąży
do $y^2>0$, dla $y=0$ wartości są zerowe).
Obie granice iterowane są więc równe $0$,
ale granica wspólna nie istnieje.
\end{przyklad}

\begin{twierdzenie}[Jednostajność względem kierunku]
\label{ma2:thm:radial-limit}
Niech $p\in\R^n$, $\rho>0$,
$B(p,\rho)\setminus\{p\}\subseteq A$ i
$S^{n-1}=\{u\in\R^n:\|u\|_2=1\}$.
Wtedy $\lim_{A\ni x\to p}f(x)=L$ zachodzi wtedy
i tylko wtedy, gdy
\[
 \lim_{r\to0^+}\sup_{u\in S^{n-1}}\|f(p+ru)-L\|_2=0,
 \qquad 0<r<\rho.
\]
Dopuszczamy supremum $+\infty$ dla niektórych $r$.
\end{twierdzenie}
\begin{proof}
Jeśli granica istnieje, dla $\varepsilon>0$ wybieramy
$\delta>0$ z definicji granicy dla dokładności $\varepsilon/2$.
Dla $0<r<\min(\rho,\delta)$ i każdego $u\in S^{n-1}$
mamy $\|p+ru-p\|_2=r<\delta$, zatem supremum nie przekracza
$\varepsilon/2$. To dowodzi granicy supremum.
Odwrotnie, dla $\varepsilon>0$ wybieramy $\delta\in(0,\rho)$,
dla którego supremum jest mniejsze od $\varepsilon$ dla
$0<r<\delta$. Dla $x\in A$ z $0<\|x-p\|_2<\delta$
kładziemy $r=\|x-p\|_2$, $u=(x-p)/r$.
Wówczas $u\in S^{n-1}$, $x=p+ru$ i
$\|f(x)-L\|_2<\varepsilon$.
\end{proof}

\begin{przyklad}
Dla $(x,y)\ne(0,0)$ niech $g(x,y)=x^2y/(x^2+y^2)$.
Ponieważ $x^2/(x^2+y^2)\le1$, mamy
$|g(x,y)|\le |y|\le\|(x,y)\|_2$.
Stąd granica jest równa $0$; przedłużenie $g(0,0)=0$
jest ciągłe. Oszacowanie nie zależy od kierunku,
w przeciwieństwie do przykładu \ref{ma2:ex:multivariable-paths}.
\end{przyklad}

\section{Ciągłość i działania na odwzorowaniach}

Dla $p\in A$ ciągłość oznacza
\[
 \forall\varepsilon>0\ \exists\delta>0\ \forall x\in A:
 \quad\|x-p\|_2<\delta\Longrightarrow
 \|f(x)-f(p)\|_2<\varepsilon.
\]
Jest to definicja AM I, 15.1.1. W punkcie skupienia
jest równoważna istnieniu granicy równej $f(p)$,
a w punkcie izolowanym zachodzi automatycznie.
Twierdzenie AM I, 15.1.3, daje równoważność z warunkiem
$f(x_k)\to f(p)$ dla każdego ciągu $(x_k)$ w $A$
zbieżnego do $p$. Z twierdzenia \ref{ma2:thm:multi-limit}
i przypadku punktów izolowanych wynika, że ciągłość
$f$ jest równoważna ciągłości wszystkich współrzędnych.

Według twierdzenia AM I, 15.2.4, $f$ jest ciągłe na $A$
wtedy i tylko wtedy, gdy przeciwobrazy zbiorów otwartych
w $\R^m$ są otwarte względnie w $A$; równoważnie
przeciwobrazy zbiorów domkniętych w $\R^m$ są domknięte
względnie w $A$. Jeśli $f:A\to B$ jest ciągłe w $p\in A$,
$B\subseteq\R^m$, a $g:B\to\R^k$ jest ciągłe w $f(p)$,
to $g\circ f$ jest ciągłe w $p$ (AM I, 15.3.1).

\begin{wniosek}
\label{ma2:cor:multi-polynomials}
Wielomiany $P:\R^n\to\R$ są ciągłe, a ilorazy
$P/Q$ wielomianów są ciągłe na otwartym zbiorze
$\{x:Q(x)\ne0\}$.
\end{wniosek}
\begin{proof}
Rzuty współrzędnych są ciągłe, gdyż
$|x_j-y_j|\le\|x-y\|_2$. Stałe są ciągłe.
Każdy wielomian jest skończoną sumą iloczynów stałych
i rzutów, więc stosujemy rachunek funkcji ciągłych
(AM I, twierdzenie 15.3.2).
Zbiór, gdzie $Q\ne0$, jest przeciwobrazem otwartego
zbioru $\R\setminus\{0\}$, a ciągłość ilorazu na nim
wynika z tego samego twierdzenia.
\end{proof}

\begin{przyklad}
Odwzorowanie $F:\R^2\to\R^3$,
$F(x,y)=(x+y,xy,\sin(x^2+y^2))$, jest ciągłe:
pierwsze dwie współrzędne są wielomianami, trzecia
jest złożeniem ciągłego sinusa z wielomianem
(AM I, rozdział 18).
\end{przyklad}

\section{Zwartość i jednostajna ciągłość}

\begin{twierdzenie}
\label{ma2:thm:euclidean-compact}
Zbiór $K\subseteq\R^n$ jest ciągowo zwarty wtedy i tylko wtedy,
gdy jest domknięty i ograniczony.
\end{twierdzenie}
\begin{proof}
Jeśli $K$ jest ciągowo zwarty, to każdy ciąg w $K$
zbieżny w $\R^n$ ma podciąg zbieżny do punktu $K$.
Jednoznaczność granicy daje domkniętość przez kryterium
ciągowe (AM I, lemat 7.5.2).
Gdyby $K$ nie był ograniczony, można by wybrać
$x_k\in K$ z $\|x_k\|_2>k$. Żaden podciąg nie byłby
zbieżny, gdyż każdy zbieżny ciąg jest ograniczony.
Otrzymujemy sprzeczność.

Niech teraz $K$ będzie domknięty i ograniczony,
a $(x_k)$ ciągiem w $K$. Stosujemy twierdzenie
Bolzana--Weierstrassa (AM I, 7.2.4) do ograniczonego
ciągu pierwszych współrzędnych, wybierając podciąg
zbieżny w tej współrzędnej. Z niego wybieramy podciąg
zbieżny w drugiej współrzędnej itd. Po $n$ krokach
otrzymujemy jeden podciąg zbieżny w każdej współrzędnej;
wybieranie dalszych podciągów nie zmienia wcześniejszych
granic. Nierówność między normą maksimum i euklidesową
daje zbieżność wektorową. Domkniętość $K$ zapewnia
przynależność granicy do $K$. Dla pustego $K$
warunek ciągowej zwartości zachodzi w sposób pusty.
\end{proof}

W przestrzeniach metrycznych zwartość i zwartość ciągowa
są równoważne (topologia ogólna, rozdział 07).
Przytaczamy trzy wyniki AM I:
\begin{enumerate}[label=(\roman*)]
 \item Jeśli $K\subseteq\R^n$ jest domknięty i ograniczony,
 a $f:K\to\R^m$ ciągłe, to $f(K)$ jest ciągowo zwarty
 (AM I, 15.4.3), zatem domknięty i ograniczony.
 \item Jeśli dodatkowo $K\ne\varnothing$ i $m=1$, to $f$
 jest ograniczona i osiąga kres dolny oraz górny
 (twierdzenie Weierstrassa, AM I, 15.4.4).
 \item Jeśli $K\subseteq\R^n$ jest domknięty i ograniczony,
 a $f:K\to\R^m$ ciągłe, to $f$ jest jednostajnie ciągłe
 (twierdzenie Heinego--Cantora, AM I, 16.2.1).
\end{enumerate}
Hipotezę ciągowej zwartości we wszystkich trzech wynikach
zapewnia twierdzenie \ref{ma2:thm:euclidean-compact}.
Jednostajna ciągłość oznacza, że dla każdego $\varepsilon>0$
istnieje $\delta>0$ wspólne dla wszystkich par $x,y\in K$
takie, że $\|x-y\|_2<\delta$ implikuje
$\|f(x)-f(y)\|_2<\varepsilon$.

\begin{przyklad}
Funkcja $P(x,y)=x^2+xy+y^2$ jest ciągła na
domkniętej kuli jednostkowej $K\subseteq\R^2$.
Osiąga tam minimum i maksimum oraz jest jednostajnie ciągła.
Na całej płaszczyźnie nie jest jednostajnie ciągła:
dla $u_k=(k,0)$, $v_k=(k+1/k,0)$ mamy
$\|u_k-v_k\|_2=1/k\to0$, lecz
$P(v_k)-P(u_k)=2+1/k^2$, co przeczy definicji
jednostajnej ciągłości z $\varepsilon=1$.
\end{przyklad}

\begin{przyklad}
Dla $x,y\in B(0,R)$, $R>0$, nierówność trójkąta daje
\[
 |\|x\|_2^2-\|y\|_2^2|
 =|\|x\|_2-\|y\|_2|(\|x\|_2+\|y\|_2)
 \le 2R\|x-y\|_2.
\]
Kwadrat normy spełnia warunek Lipschitza na każdej takiej
kuli, więc jest tam jednostajnie ciągły (AM I, 16.3.3).
\end{przyklad}

\begin{wniosek}[Przedłużenie do domknięcia]
\label{ma2:cor:multi-uniform-extension}
Jeśli $A\subseteq\R^n$ i $f:A\to\R^m$ jest
jednostajnie ciągłe, to istnieje dokładnie jedno
ciągłe przedłużenie $F:\overline A\to\R^m$.
Jest ono jednostajnie ciągłe, a jeśli dla $x,y\in A$
zachodzi $\|f(x)-f(y)\|_2\le C\|x-y\|_2^\alpha$,
$C\ge0$, $0<\alpha\le1$, to ta sama nierówność
z tymi samymi stałymi zachodzi dla $F$ na $\overline A$.
\end{wniosek}
Jest to twierdzenie AM I, 16.6.3, ponieważ $\R^m$ jest
zupełna (twierdzenie \ref{ma2:thm:finite}).
Wartość przedłużenia w $p$ jest granicą $f(x_k)$ dla
dowolnego ciągu $x_k\in A$ zbieżnego do $p$.

Przykład \ref{ma2:ex:multivariable-paths} będzie ważny
przy badaniu pochodnych kierunkowych: regularność
na każdej prostej nie gwarantuje nawet ciągłości.

\chapter{Pochodne cząstkowe i kierunkowe}
\label{ma2:ch:directional-derivatives}

W rozdziale \ref{ma2:ch:multivariable-continuity}
badaliśmy granice przy dowolnym sposobie zbliżania się
do punktu. Teraz rozpatrujemy ilorazy różnicowe wzdłuż
prostych. Korzystamy z ogólnych definicji i przykładów
analizy funkcjonalnej, rozdział 14 (dalej AF 14),
oraz z rachunku pochodnych jednej zmiennej w AM I 20.

\section{Pochodne kierunkowe i cząstkowe}

Niech $n,m$ będą dodatnimi liczbami całkowitymi,
$U\subseteq\R^n$ zbiorem otwartym, $f:U\to\R^m$
i $p\in U$. Dla $h\in\R^n$ istnieje $\delta>0$ takie,
że $p+th\in U$ dla $|t|<\delta$: jeśli
$B(p,r)\subseteq U$, dla $h\ne0$ można wziąć
$\delta=r/\|h\|_2$, a dla $h=0$ dowolne $\delta>0$.

\begin{definicja}[Pochodna kierunkowa i cząstkowa]
\label{ma2:def:directional}
Dwustronną pochodną kierunkową odwzorowania $f$ w punkcie
$p$ w kierunku $h$ nazywamy granicę
\[
 \partial_h f(p)=\lim_{\substack{t\to0\\t\in\R,\ t\ne0}}
 \frac{f(p+th)-f(p)}{t},
\]
o ile istnieje ona w $\R^m$. Jest to definicja AF 14.1.1
w przestrzeniach euklidesowych.
Dla $j=1,\ldots,n$, oznaczając przez $e_j$ $j$-ty wektor
standardowej bazy $\R^n$, definiujemy pochodną cząstkową
\[
 \partial_j f(p)=\frac{\partial f}{\partial x_j}(p)
 =\partial_{e_j}f(p).
\]
Pochodną jednostronną w kierunku $h$ otrzymujemy,
zastępując $t\to0$ przez $t\to0^+$.
\end{definicja}

Nie wymagamy $\|h\|_2=1$. Gdy $h=0$, iloraz jest zerowy
i $\partial_0 f(p)=0$. Jeśli $\gamma_h(t)=f(p+th)$,
to $\partial_h f(p)=\gamma_h'(0)$ z definicji pochodnej
krzywej. W szczególności pochodna cząstkowa jest
pochodną funkcji jednej zmiennej, w której pozostałe
współrzędne są ustalone.

\begin{stwierdzenie}[Rachunek współrzędnych]
\label{ma2:prop:directional-components}
Niech $f=(f_1,\ldots,f_m)$. Pochodna $\partial_h f(p)$
istnieje wtedy i tylko wtedy, gdy istnieje każda
$\partial_h f_i(p)$; wówczas
\[
 \partial_h f(p)=(\partial_h f_1(p),\ldots,
 \partial_h f_m(p)).
\]
Jeśli $\partial_h f(p)$ istnieje, to dla $\lambda\in\R$
istnieje $\partial_{\lambda h} f(p)$ i
$\partial_{\lambda h}f(p)=\lambda\partial_h f(p)$.
\end{stwierdzenie}
\begin{proof}
Pierwsza teza wynika z zastosowania kryterium
\ref{ma2:thm:multi-limit} do odwzorowania jednej
zmiennej $t\mapsto(f(p+th)-f(p))/t$.
Dla $\lambda=0$ obie strony drugiej równości są zerowe.
Dla $\lambda\ne0$ podstawiamy $s=\lambda t$:
\[
 \frac{f(p+t\lambda h)-f(p)}{t}
 =\lambda\frac{f(p+sh)-f(p)}s.
\]
Przy $t\to0$ parametr $s$ dąży do zera dwustronnie,
również dla $\lambda<0$. Ciągłość mnożenia przez skalar
(twierdzenie \ref{ma2:thm:metric}) daje tezę.
\end{proof}

\begin{przyklad}
Dla $f(x,y)=x^2y+\sin y$ mamy
\[
 \partial_1 f(x,y)=2xy,\qquad
 \partial_2 f(x,y)=x^2+\cos y.
\]
Pierwszy wzór uzyskujemy, traktując $y$ jako stałą,
a drugi, traktując $x$ jako stałą; stosujemy AM I,
20.3.1--2 i 20.5.2.
Dla $F(x,y)=(x^2+y,xy)$ kryterium współrzędnych
daje $\partial_1F(x,y)=(2x,y)$ i
$\partial_2F(x,y)=(1,x)$.
\end{przyklad}

\section{Interpretacja geometryczna}

Dla funkcji skalarnej $f:U\to\R$ oraz jednostkowego
$u$ pochodna $\partial_u f(p)$ opisuje zmianę wartości
na jednostkę przesunięcia wzdłuż prostej $p+tu$.
Krzywa przekroju wykresu ma postać
$t\mapsto(p+tu,f(p+tu))\in\R^{n+1}$.
Jeśli pochodna istnieje, jej wektor styczny w $t=0$
to $(u,\partial_u f(p))$, ponieważ iloraz różnicowy
pierwszej części jest stale równy $u$, a ostatniej
dąży do $\partial_u f(p)$.
Nie wyciągamy stąd wniosku o istnieniu wspólnej
płaszczyzny stycznej do całego wykresu; potrzebna
będzie różniczkowalność z następnego rozdziału.

\begin{figure}[htbp]
\centering
\begin{tikzpicture}[>=Stealth,x=1.35cm,y=1.1cm]
 \draw[->] (-1.65,0)--(2.05,0) node[right] {$t$};
 \draw[->] (0,-.25)--(0,3.45) node[above] {$f(p+tu)$};
 \draw[blue!70!black,thick,domain=-1.3:1.55,samples=65]
 plot (\x,{1.15+.65*\x+.36*\x*\x});
 \draw[red!65!black,thick] (-1.3,.305)--(1.45,2.0925)
 node[right] {styczna};
 \coordinate (P) at (0,1.15);
 \coordinate (Q) at (1.2,2.4484);
 \draw[gray,densely dashed] (P)--(1.2,1.15)--(Q);
 \draw[gray,thick] (P)--(Q);
 \fill (P) circle (1.5pt) node[left] {$f(p)$};
 \fill (Q) circle (1.5pt) node[above left] {$f(p+tu)$};
 \node[below] at (.6,1.15) {$t$};
 \node[right] at (1.2,1.78) {$\Delta f$};
 \node[blue!70!black,anchor=west] at (-1.5,2.7)
 {$\gamma_u(t)=f(p+tu)$};
\end{tikzpicture}
\caption{Pochodna wzdłuż kierunku jednostkowego jako
nachylenie stycznej do przekroju. Iloraz $\Delta f/t$
dąży do $\partial_u f(p)$.}
\label{ma2:fig:directional-section}
\end{figure}

\section{Reguły rachunkowe}

\begin{twierdzenie}
\label{ma2:thm:directional-rules}
Niech $f,g:U\to\R$ mają pochodne w kierunku $h$
w punkcie $p$. Dla $\alpha,\beta\in\R$ zachodzą
\begin{align*}
 \partial_h(\alpha f+\beta g)(p)
 &=\alpha\partial_hf(p)+\beta\partial_hg(p),\\
 \partial_h(fg)(p)
 &=\partial_hf(p)g(p)+f(p)\partial_hg(p),\\
 \partial_h(f/g)(p)
 &=\frac{\partial_hf(p)g(p)-f(p)\partial_hg(p)}
 {g(p)^2},\qquad g(p)\ne0.
\end{align*}
Jeśli $J\subseteq\R$ jest otwartym przedziałem,
$f(U)\subseteq J$, a $\varphi:J\to\R$ jest
różniczkowalna w $f(p)$, to
\[
 \partial_h(\varphi\circ f)(p)
 =\varphi'(f(p))\partial_h f(p).
\]
\end{twierdzenie}
\begin{proof}
Stosujemy twierdzenie AM I 20.3.1 do funkcji
$t\mapsto f(p+th)$ i $t\mapsto g(p+th)$ w $t=0$.
Istnienie ich pochodnych zapewnia ciągłość w zerze;
jeśli $g(p)\ne0$, mianownik pozostaje niezerowy
dla dostatecznie małych $|t|$. Ostatnia równość
wynika z jednowymiarowej reguły łańcuchowej
(AM I 20.4.1) dla $t\mapsto\varphi(f(p+th))$.
Nie jest to ogólna reguła łańcuchowa dla dwóch
odwzorowań wielowymiarowych.
\end{proof}

\begin{przyklad}
Niech $f(x,y)=\exp(xy)$. Dla $p=(a,b)$ i $h=(u,v)$
mamy $(a+tu)(b+tv)=ab+t(av+bu)+t^2uv$.
Z AM I 20.5.1 i twierdzenia \ref{ma2:thm:directional-rules}
otrzymujemy
\[
 \partial_h f(a,b)=e^{ab}(bu+av).
\]
W szczególności $\partial_1 f(a,b)=be^{ab}$
i $\partial_2 f(a,b)=ae^{ab}$.
\end{przyklad}

\section{Różniczkowalność Gâteaux i gradient}

\begin{definicja}
\label{ma2:def:gateaux}
Odwzorowanie $f:U\to\R^m$ jest różniczkowalne
w sensie Gâteaux w $p$, jeżeli istnieje operator
liniowy $A:\R^n\to\R^m$ taki, że
$\partial_hf(p)=Ah$ dla każdego $h\in\R^n$.
Oznaczamy $D_Gf(p)=A$. Jest to definicja AF 14.1.2:
w skończonym wymiarze operator liniowy jest ciągły,
gdyż, pisząc $h=\sum_jh_je_j$, mamy
\[
 \|Ah\|_2\le \sum_j|h_j|\|Ae_j\|_2
 \le \|h\|_2\sum_j\|Ae_j\|_2.
\]
\end{definicja}

\begin{stwierdzenie}
\label{ma2:prop:gateaux-matrix}
Jeśli $D_Gf(p)$ istnieje, jego macierz w standardowych
bazach ma postać
\[
 J_f(p)=(\partial_j f_i(p))_
 {1\le i\le m,\ 1\le j\le n}.
\]
W szczególności
$\partial_h f(p)=\sum_{j=1}^n h_j\partial_j f(p)$.
\end{stwierdzenie}
\begin{proof}
Z definicji $Ae_j=\partial_j f(p)$, więc $j$-tą
kolumną macierzy $A$ jest wektor $\partial_j f(p)$.
Liniowość daje wzór na $Ah$ i jednoznaczność $A$.
\end{proof}

Dla funkcji skalarnej różniczkowalnej w sensie Gâteaux
definiujemy gradient
\[
 \nabla f(p)=(\partial_1 f(p),\ldots,\partial_n f(p)).
\]
Wtedy $\partial_h f(p)=\langle\nabla f(p),h\rangle$
w standardowym iloczynie skalarnym, na mocy stwierdzenia
\ref{ma2:prop:gateaux-matrix}. W szczególności, gdy istnieje
pochodna Frécheta, jest to gradient z AF 14.1.24,
ponieważ pochodne Frécheta i Gâteaux są wtedy równe
(AF 14.1.5). Jeśli istnieją tylko pochodne cząstkowe,
ich wektor także bywa nazywany gradientem, ale sam
w sobie nie daje wzoru na pochodną w dowolnym kierunku.

\begin{wniosek}[Największa zmiana kierunkowa]
\label{ma2:cor:gradient-steepest}
Niech $f:U\to\R$ będzie różniczkowalna w sensie
Gâteaux w $p$ i niech $v=\nabla f(p)$.
Wówczas
\[
 \max_{\|u\|_2=1}\partial_u f(p)=\|v\|_2,\qquad
 \min_{\|u\|_2=1}\partial_u f(p)=-\|v\|_2.
\]
Jeśli $v\ne0$, maksimum osiągane jest dokładnie dla
$u=v/\|v\|_2$, a minimum dla $u=-v/\|v\|_2$.
Dla $v=0$ wszystkie pochodne kierunkowe są zerowe.
\end{wniosek}
\begin{proof}
Z nierówności Cauchy'ego--Schwarza
$|\langle v,u\rangle|\le\|v\|_2\|u\|_2=\|v\|_2$.
Podstawienie wskazanych wektorów jednostkowych daje
odpowiednie wartości. Dla $v\ne0$ równość w oszacowaniu
wymaga liniowej zależności $u,v$, a znak iloczynu
skalarnego rozstrzyga, który z dwóch wektorów otrzymujemy.
Dla $v=0$ wzór na pochodną daje zero w każdym kierunku.
\end{proof}

\section{Zależności i kontrprzykłady}

Istnienie pochodnych cząstkowych w $p$ nie oznacza
istnienia pochodnej w każdym kierunku. Dla
\[
 q(x,y)=
 \begin{cases}
 xy/(x^2+y^2),&(x,y)\ne(0,0),\\
 0,&(x,y)=(0,0)
 \end{cases}
\]
ograniczenia do osi są zerowe, więc obie pochodne
cząstkowe w zerze istnieją i wynoszą zero. Jednak
$q(t,t)=1/2$ dla $t\ne0$ i iloraz różnicowy
$q(t,t)/t=1/(2t)$ nie ma skończonej granicy dwustronnej.

Istnienie wszystkich pochodnych kierunkowych nie
oznacza addytywności względem kierunku. Wykorzystujemy
przykład AF 14.1.10 (etykieta źródłowa
\texttt{ex:nonlinear-directional14}): funkcja
$q_1(x,y)=x^3/(x^2+y^2)$ poza zerem i $q_1(0,0)=0$
ma pochodne
$\partial_{(a,b)}q_1(0,0)=a^3/(a^2+b^2)$
dla $(a,b)\ne0$, a w kierunku zerowym pochodną zero.
Dla kierunków $e_1,e_2,e_1+e_2$ wartości wynoszą
odpowiednio $1,0,1/2$, więc nie istnieje pochodna
Gâteaux. Sama funkcja jest ciągła w zerze, gdyż
$|q_1(x,y)|\le |x|$.

Nawet pochodna Gâteaux nie gwarantuje ciągłości
w punkcie. Z AF 14.1.8 (etykieta
\texttt{ex:directional-discontinuous14}) znamy
funkcję $g(x,y)=x^3y/(x^6+y^2)$ poza zerem,
$g(0,0)=0$: jej pochodna Gâteaux w zerze jest
operatorem zerowym, lecz $g(x,x^3)=1/2$ dla $x\ne0$.
Pełne obliczenia ilorazów różnicowych znajdują się
w przywołanym przykładzie. Zatem ciągłość ani tym
bardziej różniczkowalność Frécheta nie wynika z
różniczkowalności Gâteaux.

W następnym rozdziale zdefiniujemy pochodną Frécheta
jako liniowe przybliżenie całego przyrostu, z błędem
małym względem normy przyrostu. Wykorzystamy AF 14.1.3,
14.1.5--6 oraz 14.1.9, a także udowodnimy kryterium
różniczkowalności za pomocą ciągłych pochodnych
cząstkowych.

\chapter{Pochodna Frécheta i różniczka zupełna}
\label{ma2:ch:frechet}

Od tego rozdziału stosujemy ogólne twierdzenia rachunku
różniczkowego w rzeczywistych przestrzeniach unormowanych.
Zupełność zakładamy tylko tam, gdzie jest potrzebna.
Podstawowe wyniki pochodzą z analizy funkcjonalnej,
rozdział 14 (AF 14); wykorzystujemy ich dowody, a w razie
osłabienia hipotez wskazujemy uzasadnienie.
Przypadek euklidesowy pozostaje szczególnie ważny z uwagi
na rachunek współrzędnych i geometrię.

\section{Liniowe przybliżenie przyrostu}

Niech $E,F$ będą rzeczywistymi przestrzeniami unormowanymi,
$U\subseteq E$ otwarty i $p\in U$.
Przestrzeń operatorów liniowych ciągłych
$\mathcal L(E,F)$ wyposażamy w normę operatorową
(stwierdzenie \ref{prop:operator-norm12}).

\begin{definicja}[Pochodna Frécheta]
\label{ma2:def:frechet}
Odwzorowanie $f:U\to F$ nazywamy różniczkowalnym
w sensie Frécheta w $p$, jeśli istnieje
$A\in\mathcal L(E,F)$ takie, że
\[
 f(p+h)=f(p)+Ah+r(h),\qquad
 \lim_{\substack{h\to0\\h\ne0}}
 \frac{\|r(h)\|_F}{\|h\|_E}=0.
\]
Reszta jest określona dla $h$ z pewnej kuli wokół zera
i $r(0)=0$. Piszemy $r(h)=o(\|h\|_E)$.
Operator $A$ oznaczamy przez $Df(p)$, a także
$\mathrm df_p$ (różniczka zupełna w $p$).
Jest to definicja AF 14.3.
\end{definicja}

Dla $E=\{0\}$ warunek na resztę rozumiemy w sensie
$\varepsilon$--$\delta$ (warunek dla niezerowych $h$ jest pusty); jedyny operator
liniowy jest wtedy zerowy.
Dla $E=F=\R$ otrzymujemy definicję pochodnej i różniczki
z AM I 20.1.1--2 oraz 20.2.1, gdyż każdy operator liniowy ma
postać $h\mapsto ah$ i iloraz reszty przez $h$ dąży
do zera dokładnie wtedy, gdy jej moduł podzielony
przez $|h|$ dąży do zera.

W ogólnych przestrzeniach unormowanych pochodną
kierunkową definiujemy tym samym ilorazem co w definicji
\ref{ma2:def:directional}, z granicą w normie $F$.
Pochodna Gâteaux wymaga dodatkowo ograniczoności
operatora liniowego; w nieskończonym wymiarze sama
liniowość nie wystarcza do jego ciągłości (AF 14.2).

\begin{twierdzenie}[Ciągłość i pochodna Gâteaux]
\label{ma2:thm:frechet-gateaux}
Jeśli $f$ jest różniczkowalne w sensie Frécheta w $p$,
to jest ciągłe w $p$. Dla każdego $v\in E$ istnieje
dwustronna pochodna kierunkowa i
$\partial_v f(p)=Df(p)v$.
W szczególności pochodna Frécheta jest jednoznaczna
i $D_Gf(p)=Df(p)$, gdzie pochodna Gâteaux oznacza
ciągły operator liniowy reprezentujący wszystkie
pochodne kierunkowe (AF 14.2).
\end{twierdzenie}

Jest to dokładnie twierdzenie AF 14.5
(\texttt{thm:frechet-implies-gateaux14}); tam znajduje
się dowód dla przestrzeni unormowanych. W szczególności
nie wymaga on zupełności. Dla $E=\R^n$, $F=\R^m$
zastosowanie stwierdzenia \ref{ma2:prop:gateaux-matrix}
daje macierz $J_f(p)$ operatora $Df(p)$.

\begin{stwierdzenie}[Jednostajność względem kierunku]
\label{ma2:prop:frechet-uniform}
Jeśli $E\ne\{0\}$ i $A\in\mathcal L(E,F)$, to
$Df(p)=A$ wtedy i tylko wtedy, gdy
\[
 \lim_{t\to0,\ t\ne0}
 \sup_{\|v\|_E=1}
 \left\|\frac{f(p+tv)-f(p)}{t}-Av\right\|_F=0.
\]
Parametr $t$ jest rzeczywisty i ma moduł dostatecznie
mały, aby $p+tv\in U$ dla wszystkich jednostkowych $v$.
Supremum może a priori przyjąć wartość $+\infty$.
\end{stwierdzenie}

Stosujemy AF 14.6 (\texttt{prop:uniform-directions14}).
W porównaniu z definicją pochodnej Gâteaux nie wystarcza
granica dla każdego ustalonego $v$: potrzebna jest
jedna kontrola błędu dla wszystkich jednostkowych
kierunków (por. twierdzenie \ref{ma2:thm:radial-limit}).

\begin{przyklad}[Ciągłość i pochodna Gâteaux nie wystarczają]
\label{ma2:ex:frechet-counterexample}
W przykładzie AF 14.9 funkcja $f:\R^2\to\R$ określona
przez $f(0,0)=0$ oraz
$f(x,y)=x^3y/(x^4+y^2)$ poza zerem jest ciągła w zerze
i $D_Gf(0,0)=0$. Na paraboli $y=x^2$ dla $x\ne0$ mamy jednak
\[
 \frac{|f(x,x^2)|}{\|(x,x^2)\|_2}
 =\frac{1}{2\sqrt{1+x^2}}\longrightarrow\frac12.
\]
Pochodna Frécheta, gdyby istniała, musiałaby być zerowa
na mocy twierdzenia \ref{ma2:thm:frechet-gateaux},
co przeczy warunkowi na resztę. Pełne obliczenia
ciągłości i pochodnych kierunkowych są w AF 14.9.
\end{przyklad}

\section{Reguły różniczkowania}

\begin{twierdzenie}[Kombinacje liniowe i reguła łańcuchowa]
\label{ma2:thm:frechet-chain}
Niech $E,F,G$ będą rzeczywistymi przestrzeniami
unormowanymi, $U\subseteq E$, $V\subseteq F$ otwarte,
$p\in U$. Jeżeli $f_1,f_2:U\to F$ są różniczkowalne
w $p$, to dla $\alpha,\beta\in\R$ zachodzi
\[
 D(\alpha f_1+\beta f_2)(p)=
 \alpha Df_1(p)+\beta Df_2(p).
\]
Jeżeli $f:U\to V$ jest różniczkowalne w $p$,
a $g:V\to G$ w $f(p)$, to
\[
 D(g\circ f)(p)=Dg(f(p))\circ Df(p).
\]
Odwzorowanie afiniczne $x\mapsto Lx+b$ dla
$L\in\mathcal L(E,F)$, $b\in F$ ma pochodną $L$.
\end{twierdzenie}
\begin{proof}
Korzystamy ze stwierdzenia AF 14.14 oraz twierdzenia
AF 14.15. W źródle założono zupełność przestrzeni,
lecz dowody jej nie wykorzystują: w pierwszym
sumuje się reszty i używa nierówności trójkąta,
a w drugim oznacza się $A=Df(p)$, $B=Dg(f(p))$,
$k=f(p+h)-f(p)=Ah+r_f(h)$. Dla małych $h$ mamy
$\|r_f(h)\|\le\|h\|$, stąd
$\|k\|\le(\|A\|+1)\|h\|$.
Reszta po odjęciu $BAh$ od przyrostu złożenia
to $Br_f(h)+r_g(k)$. Pierwszy składnik ma normę
co najwyżej $\|B\|\|r_f(h)\|=o(\|h\|)$.
Dla $k\ne0$ zachodzi
\[
 \frac{\|r_g(k)\|}{\|h\|}
 \le(\|A\|+1)\frac{\|r_g(k)\|}{\|k\|}\longrightarrow0,
\]
a dla $k=0$ reszta $r_g(k)$ jest zerowa.
Operator $BA$ jest liniowy i ograniczony przez
$\|B\|\|A\|$.
Wszystkie działania odbywają się na istniejących
wektorach i ograniczonych operatorach, bez przejścia
do granicy ciągu Cauchy'ego. W przypadku afinicznym
reszta jest identycznie zerowa. Te same dowody
obowiązują zatem dla przestrzeni unormowanych.
\end{proof}

\begin{uwaga}[Dlaczego pochodna Gâteaux nie wystarcza?]
W AF 14.11 (stwierdzenie
\texttt{prop:chain-fails14}) do funkcji $f$ z przykładu
\ref{ma2:ex:frechet-counterexample} stosujemy krzywą
$\gamma(t)=(t,t^2)$. Wtedy $(f\circ\gamma)(t)=t/2$
dla każdego $t$, więc $(f\circ\gamma)'(0)=1/2$,
podczas gdy $D_Gf(0,0)\gamma'(0)=0$.
Krzywa $\gamma$ jest nawet różniczkowalna w sensie
Frécheta, gdyż $\gamma(t+s)-\gamma(t)=(s,2ts+s^2)$.
Oznacza to, że reguła łańcuchowa może zawieść już przy
zastąpieniu pochodnej zewnętrznego odwzorowania
samą pochodną Gâteaux.
\end{uwaga}

\begin{twierdzenie}[Reguła dla odwzorowań dwuliniowych]
\label{ma2:thm:frechet-bilinear}
Niech $E,E_1,E_2,G$ będą rzeczywistymi przestrzeniami
unormowanymi i $B:E_1\times E_2\to G$ dwuliniowym
odwzorowaniem spełniającym
$\|B(u,v)\|\le C\|u\|\|v\|$ dla pewnego $C\ge0$.
Wówczas $B$ jest różniczkowalne na całym iloczynie
$E_1\times E_2$ wyposażonym w normę maksimum i
\[
 DB(u,v)(h,k)=B(h,v)+B(u,k).
\]
Jeżeli $f:U\to E_1$ i $g:U\to E_2$ są różniczkowalne
w $p$, gdzie $U\subseteq E$ jest otwarty i $p\in U$, to
$z\mapsto B(f(z),g(z))$ jest różniczkowalne w $p$ i
\[
 D(B(f,g))(p)h
 =B(Df(p)h,g(p))+B(f(p),Dg(p)h).
\]
\end{twierdzenie}
\begin{proof}
Dowód AF 14.16 jest ważny bez założenia zupełności:
w normie maksimum produktu reszta jest równa
$B(h,k)$ i ma normę co najwyżej
$C\max(\|h\|,\|k\|)^2$.
Wskazany operator liniowy jest ograniczony, ponieważ
$\|B(h,v)+B(u,k)\|\le C(\|u\|+\|v\|)
\max(\|h\|,\|k\|)$.
W dowodzie drugiej części pochodna pary $(f,g)$
jest operatorem $h\mapsto(Df(p)h,Dg(p)h)$
o normie nie większej niż $\max(\|Df(p)\|,\|Dg(p)\|)$.
Iloraz normy reszty przez $\|h\|$ jest maksimum
dwóch ilorazów dążących do zera. Stosujemy następnie
twierdzenie \ref{ma2:thm:frechet-chain}.
\end{proof}

Dla $E_1=E_2=G=\R$ oraz $B(u,v)=uv$ otrzymujemy
regułę iloczynu. Niech $f,g:U\to\R$ będą
różniczkowalne w $p\in U$ i $g(p)\ne0$. Wówczas ciągłość $g$ w $p$ zapewnia,
że mianownik nie znika w pewnej kuli wokół $p$.
Stosując regułę łańcuchową do $s\mapsto1/s$
(AM I 20.3.2), a następnie regułę iloczynu, dostajemy
\[
 D(f/g)(p)h
 =\frac{g(p)Df(p)h-f(p)Dg(p)h}{g(p)^2}.
\]

\begin{przyklad}[Nieskończony wymiar]
Niech $E=C([0,1],\R)$ z normą supremum
(por. twierdzenie \ref{ma2:thm:continuous}).
Odwzorowanie $P:E\to E$, $P(u)=u^2$, jest
różniczkowalne, a $DP(u)h=2uh$.
Istotnie, mnożenie punktowe jest dwuliniowe
i $\|uv\|_\infty\le\|u\|_\infty\|v\|_\infty$.
Twierdzenie \ref{ma2:thm:frechet-bilinear} zastosowane
do $B(u,u)$ daje wskazany operator. Reszta $h^2$
ma normę $\|h\|_\infty^2$.
\end{przyklad}

\section{Pochodne częściowe na iloczynach}

Niech $E=E_1\times\cdots\times E_N$, gdzie $N\ge1$,
a $E_j$ są rzeczywistymi przestrzeniami unormowanymi.
Stosujemy normę $\|h\|=\max_j\|h_j\|$.
Przez $i_j:E_j\to E$ oznaczamy włożenie wektora
na $j$-tą pozycję, z zerami na pozostałych.
Niech $F$ będzie rzeczywistą przestrzenią unormowaną,
$U\subseteq E$ otwarty i $f:U\to F$.
Dla każdego $j=1,\ldots,N$ oraz $p\in U$ definiujemy $D_jf(p)$ jako pochodną
Frécheta funkcji $z\mapsto f(p+i_jz)$ w $z=0$,
jeśli istnieje. Dziedzina tej funkcji,
$\{z\in E_j:p+i_jz\in U\}$, jest otwarta i zawiera zero,
gdyż $i_j$ jest liniową izometrią.
Jest to pochodna częściowa względem całego bloku
zmiennych, a nie pojedynczej współrzędnej skalarnej.

\begin{stwierdzenie}
\label{ma2:prop:frechet-blocks}
Jeżeli $Df(p)$ istnieje, to istnieją wszystkie
$D_jf(p)\in\mathcal L(E_j,F)$ i
\[
 D_jf(p)=Df(p)\circ i_j,\qquad
 Df(p)h=\sum_{j=1}^N D_jf(p)h_j.
\]
\end{stwierdzenie}
\begin{proof}
Odwzorowanie $z\mapsto p+i_jz$ jest afiniczne
i jego pochodną jest $i_j$ (operator ograniczony,
gdyż $\|i_jz\|=\|z\|$). Reguła łańcuchowa daje
pierwszą równość. Druga wynika z
$h=\sum_j i_jh_j$ i liniowości $Df(p)$.
\end{proof}

\begin{twierdzenie}[Kryterium ciągłych pochodnych częściowych]
\label{ma2:thm:continuous-block-partials}
Niech $F$ będzie rzeczywistą przestrzenią unormowaną,
$U\subseteq E$ zbiorem otwartym, $f:U\to F$ i $p\in U$.
Jeżeli wszystkie pochodne częściowe $D_jf$ istnieją
w pewnym otoczeniu $p$ i są ciągłe w $p$ w normach
operatorowych, to $f$ jest różniczkowalne w sensie
Frécheta w $p$ i
\[
 Df(p)h=\sum_{j=1}^N D_jf(p)h_j.
\]
\end{twierdzenie}
\begin{proof}
Oznaczmy $A_j=D_jf(p)$ i
$Ah=\sum_j A_jh_j$. Operator $A$ jest liniowy
i ograniczony, bo $\|Ah\|\le(\sum_j\|A_j\|)\|h\|$.
Dla $\varepsilon>0$ wybieramy $\delta>0$ takie,
że kula $B(p,\delta)$ zawiera się w otoczeniu
z założenia i dla $z$ w tej kuli oraz każdego $j$
zachodzi $\|D_jf(z)-A_j\|<\varepsilon/N$.
Można użyć wspólnego promienia, bo bloków jest skończenie wiele.

Niech $\|h\|<\delta$, $p_0=p$ oraz
$p_j=p+\sum_{k=1}^j i_k h_k$.
Odcinek $p_{j-1}+t i_jh_j$, $0\le t\le1$,
zawiera się w $B(p,\delta)$, ponieważ w normie
maksimum odległość od $p$ nie przekracza $\|h\|$.
Funkcja
\[
 \gamma_j(t)=f(p_{j-1}+t i_jh_j)-t A_jh_j
\]
jest różniczkowalna (włącznie z pochodnymi jednostronnymi
na końcach): z definicji pochodnej częściowej,
po przyroście parametru $s$ przyrost $j$-tego bloku
wynosi $s h_j$. Jeśli $h_j=0$, krzywa jest stała.
W przeciwnym przypadku, dla $z=p_{j-1}+t i_jh_j$,
z definicji pochodnej częściowej mamy
$f(z+i_jk)-f(z)=D_jf(z)k+r_{j,z}(k)$,
gdzie $\|r_{j,z}(k)\|=o(\|k\|)$ dla $k\to0$.
Dla $k=sh_j$ otrzymujemy
\[
 \frac{\|r_{j,z}(sh_j)\|}{|s|}
 =\|h_j\|\frac{\|r_{j,z}(sh_j)\|}{\|sh_j\|}
 \longrightarrow0.
\]
Daje to w obu przypadkach
\[
 \gamma_j'(t)=
 (D_jf(p_{j-1}+t i_jh_j)-A_j)h_j,
 \qquad
 \|\gamma_j'(t)\|\le\frac{\varepsilon}{N}\|h_j\|.
\]
Różniczkowalność zapewnia ciągłość $\gamma_j$ na odcinku.
Lemat \ref{ma2:lem:vector-increment}, który nie wymaga
zupełności przestrzeni wartości, daje
\[
 \|f(p_j)-f(p_{j-1})-A_jh_j\|
 \le\frac{\varepsilon}{N}\|h_j\|.
\]
Sumując po $j$ i używając sumy teleskopowej, mamy
\[
 \|f(p+h)-f(p)-Ah\|
 \le\frac{\varepsilon}{N}\sum_{j=1}^N\|h_j\|
 \le\varepsilon\|h\|.
\]
Ponieważ $\varepsilon$ było dowolne, jest to
warunek różniczkowalności z pochodną $A$.
\end{proof}

\section{Klasa \texorpdfstring{$C^1$}{C1} i rachunek współrzędnych}

\begin{definicja}
\label{ma2:def:c1}
Piszemy $f\in C^1(U,F)$, jeśli $f$ jest różniczkowalne
w sensie Frécheta w każdym punkcie $U$ oraz
$Df:U\to\mathcal L(E,F)$ jest ciągłe w normie operatorowej.
Jest to definicja AF 14.18 rozszerzona na przestrzenie
unormowane; zupełność nie jest konieczna do jej
sformułowania.
\end{definicja}

\begin{wniosek}
\label{ma2:cor:c1-blocks}
Niech $U$ będzie otwartym podzbiorem skończonego
iloczynu $E$ rzeczywistych przestrzeni unormowanych
i $f:U\to F$, gdzie $F$ jest rzeczywistą przestrzenią
unormowaną. Wówczas $f$ jest klasy $C^1$ wtedy i tylko wtedy, gdy wszystkie
$D_jf$ istnieją na $U$ i są ciągłe w normach operatorowych.
\end{wniosek}
\begin{proof}
Jeśli $f$ jest klasy $C^1$, to
$D_jf(p)=Df(p)\circ i_j$ i
$\|D_jf(p)-D_jf(q)\|\le\|Df(p)-Df(q)\|$.
Odwrotnie, twierdzenie
\ref{ma2:thm:continuous-block-partials} daje pochodną
w każdym punkcie, a z jej wzoru wynika
$\|Df(p)-Df(q)\|\le\sum_j\|D_jf(p)-D_jf(q)\|$.
Skończona suma dąży do zera przy $q\to p$.
\end{proof}

\begin{stwierdzenie}[Niezależność od wyboru norm równoważnych]
\label{ma2:prop:frechet-equivalent-norms}
Zastąpienie norm w $E,F$ normami równoważnymi nie
zmienia istnienia ani wartości pochodnej Frécheta,
ani przynależności odwzorowania do klasy $C^1$.
\end{stwierdzenie}
\begin{proof}
Niech $a_E\|h\|\le\|h\|_*\le b_E\|h\|$
i $a_F\|y\|\le\|y\|_*\le b_F\|y\|$,
gdzie wszystkie cztery stałe są dodatnie.
Otwartość dziedziny nie zmienia się na mocy twierdzenia
\ref{ma2:thm:equiv}. Dla dowolnego liniowego $T:E\to F$ i
$r(h)=f(p+h)-f(p)-Th$ oraz $h\ne0$ mamy
\[
 \frac{a_F}{b_E}\frac{\|r(h)\|}{\|h\|}
 \le\frac{\|r(h)\|_*}{\|h\|_*}
 \le\frac{b_F}{a_E}\frac{\|r(h)\|}{\|h\|}.
\]
Stosując te same nierówności do $Th$ zamiast $r(h)$,
widzimy najpierw, że operator liniowy jest ograniczony
w normach pierwotnych wtedy i tylko wtedy, gdy jest
ograniczony w nowych normach. Identyczne stałe dają
$(a_F/b_E)\|T\|\le\|T\|_*\le(b_F/a_E)\|T\|$,
korzystając z równoważnego wzoru na normę operatorową
jako supremum $\|Th\|/\|h\|$ po $h\ne0$.
Dla przestrzeni zerowej operator jest zerowy i
nierówności pozostają prawdziwe.
Zatem ograniczoność operatora, zbieżność ilorazu
reszty do zera i ciągłość pochodnej są niezależne
od wyboru tych norm.
\end{proof}

W szczególności dla $E=\R^n$, $F=\R^m$
istnienie i ciągłość wszystkich pochodnych cząstkowych
funkcji współrzędnych na $U$ są równoważne
przynależności do klasy $C^1$.
Istotnie, dla $E_j=\R$ istnienie pochodnej cząstkowej
$\partial_j f(p)$ jest równoważne istnieniu pochodnej
Frécheta przekroju jednej zmiennej, gdyż
\[
 \left\|\frac{f(p+te_j)-f(p)}{t}-\partial_j f(p)\right\|
 =\frac{\|f(p+te_j)-f(p)-t\partial_j f(p)\|}{|t|}.
\]
Wtedy $D_jf(p)t=t\partial_j f(p)$, a norma różnicy
dwóch takich operatorów jest równa normie różnicy
ich wektorów współczynników. Istnienie i ciągłość
wektora pochodnej cząstkowej sprowadzają się do
odpowiednich własności jego współrzędnych
(twierdzenie \ref{ma2:thm:multi-limit} zastosowane
do ilorazu różnicowego oraz do granicy wektora).
Operator z twierdzenia
\ref{ma2:thm:continuous-block-partials} ma macierz
Jacobiego. Zmiana normy maksimum na euklidesową
jest dozwolona na mocy stwierdzenia
\ref{ma2:prop:frechet-equivalent-norms} i twierdzenia
\ref{ma2:thm:finite}.

\begin{wniosek}[Złożenie odwzorowań klasy $C^1$]
\label{ma2:cor:c1-chain}
Niech $E,F,G$ będą rzeczywistymi przestrzeniami
unormowanymi, $U\subseteq E$, $V\subseteq F$ otwarte.
Jeżeli $f\in C^1(U,F)$, $f(U)\subseteq V$ oraz
$g\in C^1(V,G)$, to $g\circ f\in C^1(U,G)$.
\end{wniosek}
\begin{proof}
Z twierdzenia \ref{ma2:thm:frechet-chain} pochodna
złożenia istnieje w każdym punkcie i wynosi
$Dg(f(x))Df(x)$. Dla ustalonego $p\in U$ zachodzi
\begin{align*}
 &\|Dg(f(x))Df(x)-Dg(f(p))Df(p)\|\\
 &\quad\le \|Dg(f(x))\|\|Df(x)-Df(p)\|\\
 &\qquad+\|Dg(f(x))-Dg(f(p))\|\|Df(p)\|.
\end{align*}
Ponieważ $f$ jest ciągłe, $Dg(f(x))\to Dg(f(p))$,
więc normy $\|Dg(f(x))\|$ są ograniczone w pewnym
otoczeniu $p$. Ponadto $Df(x)\to Df(p)$. Oba składniki
prawej strony dążą do zera, co daje ciągłość pochodnej.
\end{proof}

Istnienie pochodnej w każdym punkcie nie daje jeszcze
klasy $C^1$. Przykład AF 14.22 pokazuje to już na
$\R$: $f(0)=0$, $f(x)=x^2\sin(1/x^2)$ dla $x\ne0$.
Pochodna w zerze istnieje i wynosi zero, ale pochodna
nie jest ciągła w zerze (pełne obliczenia w AF 14.22).

Dla $f:\R^n\supseteq U\to\R^m$ różniczkowalnego w $p$
różniczka ma postać
\[
 \mathrm df_p=\sum_{j=1}^n\partial_j f(p)\,\mathrm dx_j,
 \qquad \mathrm dx_j(h)=h_j.
\]
Iloczyn wektora i funkcjonału oznacza tu operator
$h\mapsto h_j\partial_j f(p)$. Dla funkcji skalarnej
otrzymujemy $\mathrm df_p(h)=\langle\nabla f(p),h\rangle$.
W rzeczywistej przestrzeni Hilberta reprezentację
$Df(p)h=\langle h,\nabla f(p)\rangle$ zapewnia
twierdzenie Riesza (AF 14.24); w ogólnej przestrzeni
unormowanej różniczka funkcji skalarnej jest funkcjonałem
liniowym ciągłym, nie zaś kanonicznie wektorem.

\begin{przyklad}
Niech $f(x,y)=(x^2+y,xy)$ i $g(u,v)=u+v^2$.
Wszystkie pochodne cząstkowe są ciągłe, więc
$f,g$ są klasy $C^1$ na odpowiednich przestrzeniach.
Z reguły łańcuchowej
\[
 J_{g\circ f}(x,y)
 =\begin{pmatrix}1&2xy\end{pmatrix}
  \begin{pmatrix}2x&1\\y&x\end{pmatrix}
 =\begin{pmatrix}2x+2xy^2&1+2x^2y\end{pmatrix}.
\]
Jest to ten sam wynik, który otrzymujemy przez
różniczkowanie funkcji $x^2+y+x^2y^2$.
\end{przyklad}

\begin{przyklad}[Pochodne częściowe funkcji dwóch zmiennych funkcyjnych]
Niech $X=C([0,1],\R)$ z normą supremum.
Dla $H:X\times X\to X$, $H(u,v)=uv+u^2$,
mamy $D_1H(u,v)h=(v+2u)h$ oraz
$D_2H(u,v)k=uk$. Istotnie, przyrost pierwszego
bloku daje resztę $h^2$ o normie $\|h\|_\infty^2$,
a przyrost drugiego bloku jest dokładnie liniowy.
Pochodne częściowe są ciągłe w normie operatorowej:
dla operatora mnożenia $M_w:h\mapsto wh$ zachodzi
$\|M_w\|\le\|w\|_\infty$, więc różnica operatorów
jest oszacowana przez normę różnicy ich mnożników.
Z twierdzenia \ref{ma2:thm:continuous-block-partials}
otrzymujemy $DH(u,v)(h,k)=(v+2u)h+uk$.
\end{przyklad}

\section{Interpretacja geometryczna}

Dla $f:U\subseteq\R^n\to\R$ różniczkowalnej
w $p$ afiniczny wykres przybliżenia jest zbiorem
\[
 \{(p+h,f(p)+Df(p)h):h\in\R^n\}.
\]
Jego odchylenie pionowe od wykresu $f$ wynosi
$r(h)=o(\|h\|)$ na mocy definicji
\ref{ma2:def:frechet}. Dla $n=2$ jest to płaszczyzna
styczna w punkcie $(p,f(p))$, a dla dowolnego $n$
hiperpłaszczyzna styczna. Nie wnioskujemy jeszcze,
że cały wykres jest gładką podrozmaitością:
do tego potrzebna jest odpowiednia regularność w otoczeniu.

\begin{przyklad}
Dla $f(x,y)=x^2+2y^2$ w $(1,1)$ otrzymujemy
$\mathrm df_{(1,1)}(h,k)=2h+4k$.
Płaszczyzna styczna do wykresu ma równanie
$z=3+2(x-1)+4(y-1)$.
Po podstawieniu $x=1+h$, $y=1+k$ różnica wysokości
wynosi $h^2+2k^2$, co nie przekracza $2(h^2+k^2)$.
Po podzieleniu przez $\sqrt{h^2+k^2}$ iloraz dąży
do zera. Jest to bezpośrednie sprawdzenie
liniowego przybliżenia, a nie tylko obliczenie
pochodnych cząstkowych.
\end{przyklad}

W następnym rozdziale wykorzystamy regułę łańcuchową
i lemat o przyroście do ogólnych nierówności o wartości
średniej. Wzór całkowy dla krzywych o wartościach
w przestrzeni Banacha wymaga osobnego uzasadnienia
istnienia całki.

% Tymczasowa luka w numeracji: rozdziały 14–35 są w opracowaniu.
% Usuń setcounter, gdy zostaną włączone wszystkie rozdziały 11–35.
% Zamknij stronę poprzedniego rozdziału przed zmianą licznika.
\clearpage
\setcounter{chapter}{35}
% AM I 09, 18; AM II 02, 04, 07, 09–10.
\chapter{Szeregi potęgowe}
\label{ma2:ch:power-series}

Niech $\K\in\{\R,\C\}$, $E$ będzie przestrzenią
Banacha nad $\K$, a $z_0\in\K$. Rozważamy szeregi
o współczynnikach wektorowych i zmiennej skalarnej.
Przypadek $E=\K$ daje zwykłe szeregi potęgowe.

\section{Promień zbieżności}

\begin{definicja}
\label{ma2:def:power-series}
Szeregiem potęgowym o środku $z_0$ nazywamy
szereg funkcyjny
\[
 \sum_{n=0}^\infty a_n(z-z_0)^n,\qquad a_n\in E.
\]
Przyjmujemy $w^0=1$ również dla $w=0$.
Liczbę
\[
 R=\left(\limsup_{n\to\infty}\|a_n\|^{1/n}\right)^{-1}
 \in[0,+\infty]
\]
nazywamy promieniem zbieżności; przyjmujemy
$1/0=+\infty$ i $1/(+\infty)=0$.
\end{definicja}

\begin{twierdzenie}[Wzór Cauchy'ego--Hadamarda]
\label{ma2:thm:cauchy-hadamard}
Szereg potęgowy jest bezwzględnie zbieżny
dla $|z-z_0|<R$, a rozbieżny dla $|z-z_0|>R$.
W środku $z_0$ ma zawsze sumę $a_0$.
Dla każdego skończonego $0\le r<R$ szereg
jest zbieżny normalnie na zbiorze
$|z-z_0|\le r$. Jeśli $R>0$, suma szeregu jest ciągła
na zbiorze $|z-z_0|<R$, a szereg jest
zbieżny normalnie i jednostajnie na każdym
jego zwartym podzbiorze.
\end{twierdzenie}
\begin{proof}
Dla $z=z_0$ wszystkie wyrazy poza $a_0$
są zerowe. Dla $\rho=|z-z_0|>0$ mamy
\[
 \limsup_n \|a_n(z-z_0)^n\|^{1/n}
 =\rho\limsup_n\|a_n\|^{1/n}.
\]
Mnożenie przez dodatnią skończoną liczbę
zachowuje granicę górną, także nieskończoną:
dla każdego $N$ supremum ogona ciągu
$(\rho\|a_n\|^{1/n})_{n\ge N}$ jest równe
$\rho\sup_{n\ge N}\|a_n\|^{1/n}$, a następnie
bierzemy granicę po $N$.
Jeśli $0<\rho<R$, prawa strona jest mniejsza
niż $1$, także gdy $R=+\infty$. Jeśli
$\rho>R$, jest większa niż $1$, także gdy $R=0$.
Stosujemy kryterium pierwiastkowe
(twierdzenie \ref{ma2:thm:vector-ratio-root};
wersja skalarna: AM I, twierdzenie 9.6.2,
etykieta \texttt{ma1:thm:root-test}).
Dla $0<r<R$ szereg liczbowy $\sum_n\|a_n\|r^n$
jest zbieżny z tego samego kryterium.
Na domkniętym kole lub przedziale promienia $r$
supremum normy $n$-tego wyrazu wynosi
$\|a_n\|r^n$, gdyż zbiór zawiera punkt $z_0+r$.
Dla $r=0$ supremum jest zerowe dla $n\ge1$
i równe $\|a_0\|$ dla $n=0$.
Kryterium Weierstrassa
\ref{ma2:thm:weierstrass-test} daje zbieżność
jednostajną. Jeśli $K$ jest niepustym zwartym
podzbiorem otwartego koła lub przedziału,
funkcja $z\mapsto|z-z_0|$ osiąga na $K$
maksimum $r<R$. Stosujemy wynik na zbiorze
promienia $r$ i ograniczamy do $K$.
Dla pustego $K$ warunek jest automatyczny.
Każdy punkt otwartego koła lub przedziału
należy do wnętrza pewnego zbioru promienia
$r<R$, na którym zbieżność jest jednostajna.
Ciągłość sum częściowych oraz twierdzenie
\ref{ma2:thm:uniform-continuous-limit}
dają ciągłość sumy w tym punkcie.
\end{proof}

Zatem dla $0<R<+\infty$ zbieżność zachodzi
wewnątrz koła (przedziału w przypadku rzeczywistym),
a na brzegu wymaga osobnego badania.
Promień $R$ jest również supremum zbioru
liczb $r\ge0$, dla których $\sum_n\|a_n\|r^n$
jest skończony. Istotnie, zbiór ten zawiera
$0$ i wszystkie dodatnie liczby mniejsze od $R$,
a nie zawiera żadnej liczby większej od $R$.
To rozstrzyga również przypadki $R=0$ oraz
$R=+\infty$.
Ponadto dla $0<R<+\infty$ bezwzględna
zbieżność w choć jednym punkcie brzegu
jest równoważna zbieżności normalnej na
całym domkniętym kole (przedziale) promienia
$R$: oba warunki oznaczają dokładnie
$\sum_{n\ge0}\|a_n\|R^n<\infty$.
Zbieżność warunkowa w jednym punkcie nie
rozstrzyga zbieżności w innych punktach brzegu.

\begin{wniosek}[Oszacowanie reszty]
\label{ma2:cor:power-remainder}
Jeśli $0\le r<s<R$ i $s<\infty$, to przy
$C_s=\sum_{n\ge0}\|a_n\|s^n$ zachodzi
\[
 \sup_{|z-z_0|\le r}
 \left\|f(z)-\sum_{n=0}^Na_n(z-z_0)^n\right\|
 \le C_s(r/s)^{N+1},\qquad N\ge0,
\]
gdzie $f$ jest sumą szeregu.
\end{wniosek}
\begin{proof}
Dla $n>N$ mamy
$\|a_n\|r^n=\|a_n\|s^n(r/s)^n
\le\|a_n\|s^n(r/s)^{N+1}$.
Sumujemy nierówności i stosujemy oszacowanie
reszty z twierdzenia \ref{ma2:thm:weierstrass-test}.
\end{proof}

\begin{wniosek}[Kryterium ilorazowe]
\label{ma2:cor:power-ratio}
Jeżeli $a_n\ne0$ od pewnego indeksu oraz
$\lim_n\|a_{n+1}\|/\|a_n\|=L\in[0,+\infty]$,
to $R=1/L$ z powyższymi konwencjami.
\end{wniosek}
\begin{proof}
Dla $\rho>0$ iloraz norm kolejnych wyrazów
szeregu w punkcie $z_0+\rho$ ma granicę $\rho L$.
Z twierdzenia \ref{ma2:thm:vector-ratio-root}
szereg zbiega bezwzględnie dla $\rho L<1$,
a jest rozbieżny dla $\rho L>1$.
Porównując z twierdzeniem
\ref{ma2:thm:cauchy-hadamard}, otrzymujemy
$R=1/L$, ponieważ różne promienie wyznaczałyby
sprzeczne wnioski dla pewnego dodatniego $\rho$
ściśle pomiędzy nimi. Gdy $L=0$, wszystkie
dodatnie $\rho$ dają zbieżność, a gdy $L=+\infty$,
żadne dodatnie $\rho$ nie daje zbieżności.
\end{proof}

\begin{przyklad}
\label{ma2:ex:power-radii}
Szeregi $\sum_{n\ge0}n!z^n$, $\sum_{n\ge0}z^n$
oraz $\sum_{n\ge0}z^n/n!$ mają odpowiednio
promienie $0$, $1$ i $+\infty$. Ilorazy norm
współczynników wynoszą odpowiednio $n+1$,
$1$ oraz $1/(n+1)$, więc stosujemy wniosek
\ref{ma2:cor:power-ratio}.
W trzecim przypadku otrzymujemy funkcję
wykładniczą z AM I, rozdział 18, definicja
18.1.1; jej definicji nie wprowadzamy ponownie.
Wielomian ma promień nieskończony, gdyż
ciąg jego współczynników jest od pewnego
indeksu zerowy.
W przypadku współczynników z nieskończenie
wieloma zerami kryterium ilorazowe może nie
mieć zastosowania. Dla szeregu
$\sum_{k\ge0}4^kz^{2k}$ współczynnik o indeksie
$2k$ jest równy $4^k$, a o indeksie nieparzystym
jest zerowy. Dla parzystych dodatnich indeksów
pierwiastki norm wynoszą $2$, dla nieparzystych
$0$, zatem granica górna jest równa $2$ i
$R=1/2$ z twierdzenia \ref{ma2:thm:cauchy-hadamard}.
\end{przyklad}

\paragraph{Rysunek 3D: Moduł sumy szeregu geometrycznego.}
W przykładzie \ref{ma2:ex:power-radii} szereg
geometryczny ma promień zbieżności $1$. Dla $z=x+iy$
i $|z|<1$ jego suma wynosi $(1-z)^{-1}$.
Rysujemy moduł sumy na kole $|z|\le0{,}97$ oraz
moduły wybranych sum częściowych na tym samym kole.
Nie rozszerzamy wykresu sumy na obszar rozbieżności.
Kod wymaga bibliotek NumPy i Plotly; zapisuje samodzielny
interaktywny plik HTML.
\begin{lstlisting}[language=Python]
import numpy as np
import plotly.graph_objects as go

r, t = np.meshgrid(np.linspace(0, 0.97, 61),
                   np.linspace(0, 2*np.pi, 121))
Z = r*np.exp(1j*t)
fig = go.Figure(go.Surface(x=Z.real, y=Z.imag,
    z=np.abs(1/(1-Z)), colorscale="Viridis",
    name="suma", showlegend=True))
for n in [4, 16, 64]:
    S = (1-Z**(n+1))/(1-Z)
    fig.add_trace(go.Surface(x=Z.real, y=Z.imag,
        z=np.abs(S), name=f"N={n}", opacity=0.65,
        showscale=False, showlegend=True,
        visible="legendonly"))
fig.update_layout(scene=dict(xaxis_title="Re z",
    yaxis_title="Im z", zaxis_title="modul"))
fig.write_html("am2_36_geometric.html", include_plotlyjs=True)
fig.show()
\end{lstlisting}

\section{Zbieżność na brzegu}

\begin{przyklad}[Trzy zachowania na brzegu]
\label{ma2:ex:power-boundary}
Każdy z szeregów
\[
 \sum_{n\ge1}z^n,\qquad
 \sum_{n\ge1}\frac{z^n}{n},\qquad
 \sum_{n\ge1}\frac{z^n}{n^2}
\]
ma promień $1$ z wniosku \ref{ma2:cor:power-ratio}.
Pierwszy jest rozbieżny we wszystkich punktach
$|z|=1$, ponieważ normy wyrazów są równe $1$.
Trzeci jest zbieżny bezwzględnie w każdym
takim punkcie i normalnie na całym domkniętym
kole, ponieważ $\sum 1/n^2<\infty$
(AM I, rozdział 09, przykład
\texttt{ma1:ex:integer-p-series}).
Drugi jest rozbieżny w $z=1$ (szereg harmoniczny).
Dla $|z|=1$, $z\ne1$, zachodzi
\[
 \left|\sum_{n=1}^k z^n\right|
 =\left|\frac{z(1-z^k)}{1-z}\right|
 \le\frac{2}{|1-z|}.
\]
Z kryterium Dirichleta
\ref{ma2:thm:uniform-dirichlet} na zbiorze
jednoelementowym wynika zbieżność. Nie jest
ona bezwzględna, gdyż $\sum |z^n/n|=\sum1/n$.
Dokładniej, dla $\delta>0$ na zbiorze
$D_\delta=\{z:|z|\le1,\ |1-z|\ge\delta\}$
to samo oszacowanie sum geometrycznych
obowiązuje ze stałą $2/\delta$, ponieważ
$|1-z^k|\le1+|z|^k\le2$.
Kryterium Dirichleta daje jednostajną zbieżność
na $D_\delta$ i błąd po $N$ wyrazach
nie większy niż $4/(\delta(N+1))$.
W szczególności suma jest ciągła na
domkniętym kole z wyłączeniem punktu $1$
w topologii względnej: każdy punkt tego zbioru
ma względne otoczenie zawarte w pewnym
$D_\delta$, a sumy częściowe są ciągłe.
Na całym otwartym kole zbieżność nie jest
jednostajna: dla $p\ge1$ wybieramy
$z_p=2^{-1/(2p)}\in(0,1)$.
Wtedy dla $p\le n\le2p$ mamy
$z_p^n\ge1/2$ i
\[
 \sum_{n=p}^{2p}\frac{z_p^n}{n}
 \ge\frac{p+1}{4p}\ge\frac14.
\]
Przeczy to jednostajnemu kryterium Cauchy'ego
\ref{ma2:thm:function-series-cauchy}. W przypadku rzeczywistym
drugi szereg zbiega warunkowo w $-1$
i jest rozbieżny w $1$.
\end{przyklad}

\begin{figure}[htbp]
\centering
\begin{tikzpicture}[x=1.5cm,y=1.5cm,>=Stealth]
 \foreach \p/\opis in {0/{A},3/{B},6/{C}}{
  \begin{scope}[shift={(\p,0)}]
   \fill[blue!10] (0,0) circle[radius=1];
   \draw[->,gray] (-1.15,0)--(1.2,0);
   \draw[->,gray] (0,-1.15)--(0,1.2);
   \node[below] at (0,-1.2) {\opis};
   \node[font=\scriptsize,above left] at (0,0) {$0$};
  \end{scope}
 }
 \draw[red!75!black,very thick,dashed] (0,0) circle[radius=1];
 \draw[orange!85!black,very thick] (3,0) circle[radius=1];
 \filldraw[fill=white,draw=red!75!black,thick] (4,0) circle[radius=2pt];
 \node[font=\scriptsize,above right] at (4,0) {$1$};
 \draw[blue!70!black,very thick] (6,0) circle[radius=1];
\end{tikzpicture}

\small
A: $\sum z^n$; \quad
B: $\sum z^n/n$; \quad
C: $\sum z^n/n^2$.
\caption{Koła o promieniu zbieżności $1$.
Wewnątrz zachodzi zbieżność bezwzględna.
Brzeg przerywany oznacza rozbieżność,
pomarańczowy --- zbieżność warunkową
(z wyłączeniem punktu $1$), niebieski ---
zbieżność bezwzględną.}
\label{ma2:fig:power-boundary}
\end{figure}
\FloatBarrier

\begin{przyklad}[Zbieżność warunkowa na całym brzegu]
\label{ma2:ex:power-vector-boundary}
W przestrzeni Banacha $E=\ell^2(\K)$
(por. twierdzenie \ref{thm:seq-banach12})
niech $e_n$ oznacza ciąg o jedynce na $n$-tej
współrzędnej i zerach na pozostałych.
Szereg $\sum_{n\ge1}e_n z^n/n$ ma promień
$1$, ponieważ $\|e_n/n\|_2=1/n$ i iloraz
kolejnych norm dąży do $1$.
Dla $|z|\le1$ i $q\ge p\ge1$ z definicji
normy $\ell^2$ otrzymujemy
\[
 \left\|\sum_{n=p}^q\frac{e_n z^n}{n}\right\|_2^2
 =\sum_{n=p}^q\frac{|z|^{2n}}{n^2}
 \le\sum_{n=p}^{\infty}\frac1{n^2}.
\]
Prawa strona dąży do zera niezależnie od $z,q$.
Jednostajne kryterium Cauchy'ego
\ref{ma2:thm:function-series-cauchy}
i zupełność $\ell^2$ dają jednostajną zbieżność
na domkniętym kole. W każdym punkcie
$|z|=1$ szereg norm jest szeregiem harmonicznym,
więc zbieżność jest warunkowa. W szczególności
jednostajna zbieżność na całym domkniętym kole
nie wymaga zbieżności normalnej.
\end{przyklad}

\begin{twierdzenie}[Twierdzenie Abela o granicy radialnej]
\label{ma2:thm:power-abel-boundary}
Niech $0<R<+\infty$ i $|\zeta-z_0|=R$.
Jeżeli szereg $\sum_n a_n(\zeta-z_0)^n$
jest zbieżny do $S$, to
\[
 \lim_{t\to1^-}f(z_0+t(\zeta-z_0))=S.
\]
Co więcej, szereg po podstawieniu
$z=z_0+t(\zeta-z_0)$ zbiega jednostajnie
dla $t\in[0,1]$.
\end{twierdzenie}
\begin{proof}
Połóżmy $b_n=a_n(\zeta-z_0)^n$.
Szereg $\sum b_n$ rozpatrywany jako szereg
funkcji stałych jest jednostajnie zbieżny
na $[0,1]$. Ciąg mnożników $t^n$ jest
dla każdego $t\in[0,1]$ nierosnący
i ograniczony przez $1$.
Jednostajne kryterium Abela
\ref{ma2:thm:uniform-abel}, po przesunięciu
indeksu początkowego, daje jednostajną
zbieżność szeregu $\sum b_nt^n$ na $[0,1]$.
Jego suma jest ciągła z twierdzenia
\ref{ma2:thm:uniform-continuous-limit},
gdyż sumy częściowe są ciągłe.
W $t=1$ wynosi $S$, a w $0\le t<1$
jest równa $f(z_0+t(\zeta-z_0))$.
Ciągłość tej sumy na $[0,1]$ w punkcie $1$
daje żądaną granicę jednostronną.
\end{proof}
Odwrotność twierdzenia nie zachodzi.
Dla $f(z)=1/(1+z)=\sum_{n\ge0}(-z)^n$
wewnątrz koła jednostkowego granica radialna
w $z=1$ wynosi $1/2$, lecz szereg na brzegu
$\sum_{n\ge0}(-1)^n$ jest rozbieżny.

\section{Różniczkowanie i całkowanie}

\begin{twierdzenie}[Promień szeregu pochodnych]
\label{ma2:thm:power-derivative-radius}
Szereg
$\sum_{n\ge1}na_n(z-z_0)^{n-1}$
ma ten sam promień zbieżności co
$\sum_{n\ge0}a_n(z-z_0)^n$.
\end{twierdzenie}
\begin{proof}
Niech $R'$ będzie promieniem szeregu pochodnych.
Jeżeli $0<r<s<R$, to zbiega szereg
$\sum_n\|a_n\|s^n$. Liczby
\[
 n r^{n-1}/s^n=\frac{n}{s}(r/s)^{n-1}
\]
są ograniczone, gdyż iloraz kolejnych
dodatnich liczb dąży do $r/s<1$,
więc od pewnego indeksu maleją.
Jeśli $M$ jest wspólnym ograniczeniem tych
liczb, to $n\|a_n\|r^{n-1}\le M\|a_n\|s^n$.
Z porównania szeregów nieujemnych (AM I,
twierdzenie 9.3.2) wynika zbieżność
$\sum_{n\ge1}n\|a_n\|r^{n-1}$. Dla każdego $r<R$ dodatniego
można wybrać takie $s$, więc $R'\ge R$
(dla $R=0$ nierówność jest automatyczna,
a dla $R=+\infty$ otrzymujemy $R'\ge r$
dla każdej skończonej liczby $r>0$).
Odwrotnie, dla $0<r<R'$ mamy
\[
 \sum_{n\ge1}\|a_n\|r^n
 \le r\sum_{n\ge1}n\|a_n\|r^{n-1}<\infty.
\]
Po dodaniu $\|a_0\|$ wynika $R\ge R'$.
Gdy $R'=0$, ta nierówność jest automatyczna.
\end{proof}

\begin{twierdzenie}[Różniczkowanie wyraz po wyrazie]
\label{ma2:thm:power-derivative}
Jeżeli $R>0$, suma $f$ szeregu potęgowego
jest różniczkowalna w każdym punkcie
$|z-z_0|<R$ (w przypadku zespolonym
w sensie granicy zespolonych ilorazów różnicowych)
oraz
\[
 f'(z)=\sum_{n\ge1}na_n(z-z_0)^{n-1}.
\]
Dla każdego $k\ge1$ istnieje pochodna $k$-tego rzędu
i
\[
 f^{(k)}(z)=
 \sum_{n=k}^\infty\frac{n!}{(n-k)!}a_n(z-z_0)^{n-k}.
\]
Wszystkie te szeregi mają promień $R$
i zbiegają normalnie na zbiorach zwartych
wewnątrz obszaru zbieżności. Każda pochodna
$f^{(k)}$ jest tam ciągła.
\end{twierdzenie}
\begin{proof}
Ustalmy $z$ wewnątrz obszaru zbieżności.
Wybieramy skończone $r$ z $|z-z_0|<r<R$.
Dla $0<|h|<r-|z-z_0|$ mamy
$|z+h-z_0|\le |z-z_0|+|h|<r$. Dla $w=z-z_0$ i $n\ge1$
tożsamość różnicy potęg daje
\[
 \frac{(w+h)^n-w^n}{h}
 =\sum_{j=0}^{n-1}(w+h)^{n-1-j}w^j.
\]
Każdy składnik ma moduł nie większy niż
$r^{n-1}$, a dla ustalonego $n$ prawa
strona dąży do $nw^{n-1}$ przy $h\to0$.
Z twierdzenia \ref{ma2:thm:power-derivative-radius}
szereg $\sum_{n\ge1}n\|a_n\|r^{n-1}$ jest
zbieżny. Kryterium Weierstrassa
\ref{ma2:thm:weierstrass-test} zastosowane
do ilorazów różnicowych (jako funkcji $h$)
zapewnia jednostajną zbieżność ich szeregu
na dostatecznie małym przebitym otoczeniu zera.
Z liniowości granic sum częściowych suma
tego szeregu jest równa $(f(z+h)-f(z))/h$.
Z twierdzenia \ref{ma2:thm:interchange-limits},
stosowanego do sum częściowych, możemy
przejść do granicy przy $h\to0$ wyraz po wyrazie.
Otrzymujemy wzór na $f'$.
Kolejne wzory wynikają indukcyjnie:
szereg pochodnej ma ten sam promień
z twierdzenia \ref{ma2:thm:power-derivative-radius},
więc na każdym kroku stosujemy już
udowodnioną część twierdzenia. Współczynniki
mnożą się przez kolejny wykładnik potęgi,
co daje $n!/(n-k)!$: przy przejściu
od $k$ do $k+1$ wyraz stały $n=k$ znika,
a dla $n\ge k+1$ mamy
$(n-k)n!/(n-k)!=n!/(n-k-1)!$.
Normalność i ciągłość każdej pochodnej
wynikają z twierdzenia \ref{ma2:thm:cauchy-hadamard}.
\end{proof}
Równość promieni nie oznacza identycznego
zachowania na brzegu. Szereg $\sum_{n\ge1}z^n/n^2$
jest normalnie zbieżny na domkniętym kole
jednostkowym, lecz jego szereg pochodnych
$\sum_{n\ge1}z^{n-1}/n$ jest rozbieżny w $z=1$,
gdzie staje się szeregiem harmonicznym.

\begin{wniosek}[Jednoznaczność współczynników]
\label{ma2:cor:power-coefficients}
Dla szeregu o dodatnim promieniu zbieżności
zachodzi $a_k=f^{(k)}(z_0)/k!$ dla $k\ge0$,
gdzie $f^{(0)}=f$.
W szczególności dwa szeregi o tym samym środku
i dodatnich promieniach, których sumy zgadzają
się w otoczeniu środka, mają identyczne współczynniki.
\end{wniosek}
\begin{proof}
Wstawiamy $z=z_0$ do wzoru z twierdzenia
\ref{ma2:thm:power-derivative}; wszystkie składniki
poza $n=k$ znikają. Dla $k=0$ korzystamy
z $f(z_0)=a_0$. Równość funkcji w otoczeniu
środka pociąga równość wszystkich ich pochodnych
w środku, z definicji pochodnej i indukcji.
\end{proof}

\begin{wniosek}[Funkcja pierwotna i całkowanie]
\label{ma2:cor:power-primitive}
Jeśli $R>0$ i $c\in E$, szereg
\[
 F(z)=c+\sum_{n\ge0}\frac{a_n}{n+1}(z-z_0)^{n+1}
\]
ma promień $R$, a $F'=f$ dla $|z-z_0|<R$.
Dla rzeczywistej zmiennej i dowolnych
$u<v$ z $[u,v]\subseteq\{t\in\R:|t-z_0|<R\}$
zachodzi
\[
 \int_u^v f(t)\,\mathrm dt
 =\sum_{n\ge0}\frac{a_n}{n+1}
      ((v-z_0)^{n+1}-(u-z_0)^{n+1}).
\]
\end{wniosek}
\begin{proof}
Formalny szereg pochodnych szeregu $F$ jest
szeregiem $f$, więc równość promieni wynika
z twierdzenia \ref{ma2:thm:power-derivative-radius},
a wzór na pochodną z twierdzenia
\ref{ma2:thm:power-derivative}.
Na $[u,v]$ szereg $f$ zbiega jednostajnie.
Wniosek \ref{ma2:cor:termwise-integral} pozwala
całkować wyraz po wyrazie; całka wektora stałego
pomnożonego przez skalarną funkcję jest iloczynem
tego wektora i całki skalarnej, co wynika wprost
z sum Riemanna. Całkę potęgi obliczamy
z twierdzenia Newtona--Leibniza (AM I, 29.6.3).
\end{proof}

\section{Działania algebraiczne}

\begin{twierdzenie}
\label{ma2:thm:power-algebra}
Niech $f=\sum a_n(z-z_0)^n$ i
$g=\sum b_n(z-z_0)^n$ będą szeregami
o promieniach $R_f,R_g$, ze współczynnikami
odpowiednio w przestrzeniach Banacha $E,F$
nad tym samym ciałem $\K$.
Jeśli $T:E\times F\to G$ jest ciągłym
odwzorowaniem dwuliniowym do przestrzeni
Banacha $G$, to szereg
\[
 \sum_{n\ge0}c_n(z-z_0)^n,\qquad
 c_n=\sum_{j=0}^n T(a_j,b_{n-j}),
\]
ma promień co najmniej $\min(R_f,R_g)$
i wewnątrz tego obszaru sumuje się do $T(f(z),g(z))$.
Jeśli $E=F$, analogicznie dla $\alpha,\beta\in\K$
szereg o współczynnikach $\alpha a_n+\beta b_n$
ma promień co najmniej $\min(R_f,R_g)$
i sumę $\alpha f+\beta g$ w tym obszarze.
\end{twierdzenie}
\begin{proof}
Dla $0<r<\min(R_f,R_g)$ oba szeregi
$\sum\|a_n\|r^n$ i $\sum\|b_n\|r^n$ są zbieżne.
Z lematu \ref{ma2:lem:bilinear-bound}
istnieje $K\ge0$ z
$\|T(a,b)\|\le K\|a\|\|b\|$.
Dla każdego $N$ mamy
\[
 \sum_{n=0}^N\|c_n\|r^n
 \le K\left(\sum_{j\ge0}\|a_j\|r^j\right)
      \left(\sum_{k\ge0}\|b_k\|r^k\right).
\]
Wynika to z nierówności trójkąta,
$r^n=r^jr^{n-j}$ i rozszerzenia sumowania
na wszystkie pary nieujemnych indeksów.
Sumy częściowe szeregu norm są niemalejące
i ograniczone, zatem szereg norm jest zbieżny.
Wynika stąd dolne ograniczenie promienia;
gdy minimum jest zerem, jest ono automatyczne.
Dla $|z-z_0|<\min(R_f,R_g)$ stosujemy
wniosek \ref{ma2:cor:bilinear-absolute}
do obu szeregów w tym punkcie.
Dwuliniowość daje
$T(a_j(z-z_0)^j,b_k(z-z_0)^k)
=T(a_j,b_k)(z-z_0)^{j+k}$.
Zatem iloczyn Cauchy'ego jest dokładnie
wskazanym szeregiem. Dla kombinacji liniowej i $0<r<\min(R_f,R_g)$
mamy
\[
 \sum_{n\ge0}\|\alpha a_n+\beta b_n\|r^n
 \le|\alpha|\sum_{n\ge0}\|a_n\|r^n
   +|\beta|\sum_{n\ge0}\|b_n\|r^n<\infty.
\]
Zatem promień jest nie mniejszy niż minimum,
a liniowość granic sum częściowych daje wzór
na sumę. Jeżeli minimum jest zerem, oszacowanie
promienia jest automatyczne. W środku $z_0$
obie równości dla sum zachodzą niezależnie
od promieni, ponieważ pozostaje tylko wyraz
o indeksie zero.
\end{proof}

\begin{przyklad}
\label{ma2:ex:power-cancellation}
Promień wynikowego szeregu może wzrosnąć.
Dla $f(z)=\sum_{n\ge0}z^n$ i $g(z)=1-z$
iloczyn Cauchy'ego ma współczynnik $1$ dla
$n=0$ i $1-1=0$ dla $n\ge1$.
Ma więc promień $+\infty$, chociaż
$\min(R_f,R_g)=1$.
Podobnie $f-f$ jest szeregiem zerowym.
\end{przyklad}
W przypadku algebry Banacha możemy jako $T$
wziąć mnożenie, nawet gdy algebra nie jest
przemienna. Kolejność czynników we współczynnikach
iloczynu jest wtedy istotna. W rozdziale
\ref{ma2:ch:banach-algebras}, twierdzenie
\ref{ma2:thm:algebra-power-series}, rozpatrzono
odmienną sytuację: skalarne współczynniki i argument
będący elementem algebry Banacha.

\section{Przykłady rozwinięć}

\begin{przyklad}
\label{ma2:ex:power-log-arctan}
Z szeregu geometrycznego i podstawienia $z=-t^2$
dla $|t|<1$ wynika
\[
 \frac1{1+t^2}=\sum_{n\ge0}(-1)^nt^{2n}.
\]
Całkując na odcinku o końcach $0,x$,
gdzie $|x|<1$, dostajemy
\[
 \arctan x=\sum_{n\ge0}\frac{(-1)^nx^{2n+1}}{2n+1}.
\]
Wykorzystano pochodną arcus tangens z AM I
(twierdzenie 20.5.3), wartość $\arctan0=0$
oraz twierdzenie Newtona--Leibniza 29.6.3.
Dla $x<0$ stosujemy konwencję odwróconej
orientacji całki z rozdziału
\ref{ma2:ch:limits-calculus}.
Szereg dla $1/(1+t^2)$ ma promień $1$,
bo zbiega bezwzględnie dla $|t|<1$, a dla
$|t|>1$ jego wyrazy nie dążą do zera.
Szereg dla arcus tangens ma zatem promień
$1$ z twierdzenia \ref{ma2:thm:power-derivative-radius},
gdyż jego formalna pochodna jest właśnie
tym szeregiem. Dla $x=1$ szereg po prawej stronie jest
zbieżny z kryterium Leibniza
\ref{ma1:thm:leibniz-series}, ponieważ
$1/(2n+1)$ maleje do zera. Twierdzenie
\ref{ma2:thm:power-abel-boundary} i ciągłość
arcus tangens dają
\[
 \frac{\pi}{4}=\sum_{n\ge0}\frac{(-1)^n}{2n+1}.
\]
Bezpośrednio całkując od $0$ do $x$ szereg
geometryczny $\sum_{n\ge0}(-t)^n=1/(1+t)$
na odcinku o końcach $0,x$, otrzymujemy
\[
 \log(1+x)=\sum_{n\ge1}\frac{(-1)^{n-1}x^n}{n},
 \qquad |x|<1.
\]
Korzystamy z pochodnej logarytmu
(AM I, twierdzenie 20.5.1) i $\log1=0$,
podobnie jak w przykładzie
\ref{ma2:ex:termwise-geometric-integral}.
Promień szeregu logarytmicznego wynosi $1$
z wniosku \ref{ma2:cor:power-ratio}.
Dla $x=1$ kryterium Leibniza
i twierdzenie Abela dają sumę $\log2$.
\end{przyklad}

Powyższe rozwinięcia są podstawą dalszego
badania zespolonych funkcji elementarnych,
funkcji analitycznych i szeregów Taylora.

\end{document}
